% Recherche de nouvelles sources d'énergie — SVT Tle D — Cours signé OnBuch+
% Programme officiel MINESEC. Compiler avec : tectonic N19-energies-renouvelables-biocarburants.tex
\newcommand{\DOCMATIERE}{SVT}
\newcommand{\DOCNIVEAU}{Tle D}
\newcommand{\DOCMODULE}{Module IV : La Biotechnologie — Énergies renouvelables}
\newcommand{\DOCLECON}{N19}
\newcommand{\DOCTITRE}{Recherche de nouvelles sources d'énergie}
% OnBuch+ — préambule « cours signés » ENRICHI (compilé avec tectonic / XeLaTeX)
% Système visuel « Pop » d'OnBuch+ : fond crème (jamais blanc), bordures encre
% épaisses, ombres « dures » décalées, titres Archivo Black en capitales.
% Version MATHÉMATIQUES : graphes (pgfplots/tikz), boîtes pédagogiques
% (définition, propriété, méthode, attention, le savais-tu, exercice résolu,
% exercice, TP, bilan), page de garde et signature OnBuch+ ancrée en bas.
%
% Macros attendues dans main.tex AVANT \input{preamble} :
%   \DOCMATIERE  \DOCNIVEAU  \DOCMODULE  \DOCLECON (numéro)  \DOCTITRE
\documentclass[11pt]{article}

\usepackage{fontspec}
\usepackage[french]{babel}
\usepackage[a4paper,margin=2cm,top=3.4cm,bottom=2.8cm,headheight=1.4cm,footskip=1.3cm]{geometry}
\usepackage{amsmath,amssymb,amsfonts}
\usepackage{mathtools} % \xmapsto, \coloneqq... souvent utilisés par les modèles
\usepackage{cancel} % simplifications barrées (bilans de forces)
% \abs{z} (module, valeur absolue) et \norm{u} (norme), utilisés par les modèles
\providecommand{\abs}{}\let\abs\relax\DeclarePairedDelimiter{\abs}{\lvert}{\rvert}
\providecommand{\norm}{}\let\norm\relax\DeclarePairedDelimiter{\norm}{\lVert}{\rVert}
\usepackage[table]{xcolor} % [table] : fournit aussi \rowcolors (alternance de lignes)
\usepackage{pagecolor}
\usepackage{tikz}
\usetikzlibrary{arrows.meta,positioning,calc,shapes.geometric,shapes.misc,shapes.arrows,decorations.pathmorphing,decorations.pathreplacing,decorations.markings,patterns,shadows,mindmap,trees,fit,backgrounds,matrix,intersections,angles,quotes}
\usepackage{pgfplots}
\usepackage[version=4]{mhchem} % équations nucléaires \ce{^{235}_{92}U}
\usepackage[european,siunitx]{circuitikz} % schémas électriques
\pgfplotsset{compat=1.18}
\usepgfplotslibrary{fillbetween,groupplots} % aires entre courbes (calcul intégral)
\usepackage{enumitem}
\usepackage{fancyhdr}
% [most] charge toutes les bibliothèques tcolorbox (listings, minted, raster,
% documentation...) et gonfle inutilement la mémoire TeX sur un document long
% (~300 boîtes) : on ne charge que ce qui est réellement utilisé.
\usepackage[skins,breakable]{tcolorbox}
\usepackage{graphicx}
\usepackage{colortbl}
\usepackage{booktabs}
\usepackage{tabularx}
\usepackage{diagbox} % case d'en-tête diagonale des tableaux à double entrée
% Colonnes extensibles centrée / gauche / droite (C, L, R), souvent utilisées par les modèles
\newcolumntype{C}{>{\centering\arraybackslash}X}
\newcolumntype{L}{>{\raggedright\arraybackslash}X}
\newcolumntype{R}{>{\raggedleft\arraybackslash}X}
\usepackage{array}
\usepackage{multicol}
\usepackage{siunitx}
% parse-numbers=false : accepte aussi \SI{2,5 \times 10^{-3}}{...}
\sisetup{locale=FR,output-decimal-marker={,},group-minimum-digits=4,parse-numbers=false}
\DeclareSIUnit\ohm{\ensuremath{\symup{\Omega}}} % U+2126 absent de la police
\DeclareSIUnit\division{div} % réglages d'oscilloscope (ms/div, V/div)
\DeclareSIUnit\tr{tr} % tours (vitesse de rotation, ex. \SI{1500}{\tr\per\minute})
\DeclareSIUnit\molar{M} % concentration molaire (ex. \SI{3}{\molar}, \SI{1}{\micro\molar})
\DeclareSIUnit\an{an} % année (le modèle confond parfois avec \year, un compteur TeX) — ex. \SI{5}{\kilo\meter\per\an}
\DeclareSIUnit\FCFA{FCFA} % franc CFA (ex. \SI{500}{\FCFA\per\kilo\gram})
\DeclareSIUnit\degreeBrix{\textdegree Brix} % teneur en sucre d'un sirop/jus (réfractométrie)
\DeclareSIUnit\month{mois} % le modèle utilise parfois \month, un compteur TeX, comme unité de durée
\usepackage{qrcode}
% \degree : commande fréquemment utilisée par les modèles pour un angle ou une
% température en dehors de \SI{}{} (non fournie par les paquets chargés).
\providecommand{\degree}{^{\circ}}
% \qed (fin de démonstration), utilisé par les modèles sans amsthm
\providecommand{\qed}{\hfill$\square$}
% \barre : double barre des valeurs interdites dans un tableau de variations (tkz-tab)
\providecommand{\barre}{\ensuremath{\|}}
% environnement proof (amsthm non chargé) utilisé par les modèles
\ifdefined\proof\else\newenvironment{proof}[1][Démonstration]{\par\noindent\textit{#1.}\ }{\qed\par}\fi
\usepackage{lastpage}
\usepackage{hyperref}
\hypersetup{hidelinks}

% ============================================================
% Palette « Pop » OnBuch+ — mêmes hex que la classe P (plus_ui.dart)
% ============================================================
\definecolor{poporange}{HTML}{F59321}
\definecolor{popdark}{HTML}{E07A0C}
\definecolor{popcream}{HTML}{FFF6E8}
\definecolor{popink}{HTML}{1C1714}
\definecolor{poppurple}{HTML}{7A5AE0}
\definecolor{popgreen}{HTML}{2E9E6A}
\definecolor{poppink}{HTML}{E0457B}
\definecolor{popblue}{HTML}{2F7DE1}
\definecolor{popgold}{HTML}{C98A12}
\definecolor{popmuted}{HTML}{8A7F76}
\definecolor{popline}{HTML}{EADFD2}
\definecolor{popgray}{HTML}{8A7F76}
\colorlet{popgrey}{popgray}
% Alias tolérants (le modèle nomme parfois les couleurs différemment)
\colorlet{popwhite}{white}
\colorlet{popblack}{popink}
\colorlet{popred}{poppink}
\colorlet{popyellow}{popgold}
% Teintes claires (fonds de tableaux/encadrés) — le modèle nomme parfois
% les couleurs avec un suffixe L (ex. popblueL) sans les avoir définies.
\colorlet{poporangeL}{poporange!20}
\colorlet{popdarkL}{popdark!20}
\colorlet{poppurpleL}{poppurple!20}
\colorlet{popgreenL}{popgreen!20}
\colorlet{poppinkL}{poppink!20}
\colorlet{popblueL}{popblue!20}
\colorlet{popgoldL}{popgold!20}
\colorlet{popredL}{popred!20}
\colorlet{popgrayL}{popgray!20}
% Teintes claires pour fonds de tableaux / zones
\colorlet{popblueL}{popblue!10!white}
\colorlet{poporangeL}{poporange!14!white}
\colorlet{popgreenL}{popgreen!12!white}
\colorlet{poppurpleL}{poppurple!12!white}
\colorlet{poppinkL}{poppink!12!white}
\colorlet{popgoldL}{popgold!14!white}

\pagecolor{popcream}

% ============================================================
% Typographie
% ============================================================
\defaultfontfeatures{Ligatures=TeX}
\setmainfont{PlusJakartaSans-Medium.ttf}[Path=../../../fonts/,BoldFont=PlusJakartaSans-Bold.ttf]
% Maths en sans-serif (Fira Math) avec chiffres et lettres droites de Plus
% Jakarta, pour que les formules se fondent dans le texte.
\usepackage[math-style=ISO,bold-style=ISO]{unicode-math}
\setmathfont{FiraMath-Regular.otf}[Path=../../../fonts/]
\setmathfont{PlusJakartaSans-Medium.ttf}[Path=../../../fonts/,range={up/num,up/{Latin,latin},\mathup}]
\usepackage{icomma} % virgule décimale sans espace en mode math
% Caractères mathématiques Unicode absents des polices de texte (Plus Jakarta,
% Archivo Black) : rendus par leur équivalent mathématique, en texte comme en
% formule — sinon ils s'affichent en carré vide (ex. « Arithmétique dans ℤ »).
\usepackage{newunicodechar}
\newunicodechar{ℝ}{\ensuremath{\mathbb{R}}}
\newunicodechar{ℤ}{\ensuremath{\mathbb{Z}}}
\newunicodechar{ℂ}{\ensuremath{\mathbb{C}}}
\newunicodechar{ℕ}{\ensuremath{\mathbb{N}}}
\newunicodechar{ℚ}{\ensuremath{\mathbb{Q}}}
\newunicodechar{𝔻}{\ensuremath{\mathbb{D}}}
\newunicodechar{α}{\ensuremath{\alpha}}
\newunicodechar{β}{\ensuremath{\beta}}
\newunicodechar{γ}{\ensuremath{\gamma}}
\newunicodechar{δ}{\ensuremath{\delta}}
\newunicodechar{ε}{\ensuremath{\varepsilon}}
\newunicodechar{θ}{\ensuremath{\theta}}
\newunicodechar{λ}{\ensuremath{\lambda}}
\newunicodechar{μ}{\ensuremath{\mu}}
\newunicodechar{σ}{\ensuremath{\sigma}}
\newunicodechar{τ}{\ensuremath{\tau}}
\newunicodechar{φ}{\ensuremath{\varphi}}
\newunicodechar{ψ}{\ensuremath{\psi}}
\newunicodechar{ω}{\ensuremath{\omega}}
\newunicodechar{Σ}{\ensuremath{\Sigma}}
% majuscules produites par \MakeUppercase dans les titres (α-aminés -> Α)
\newunicodechar{Α}{\ensuremath{\alpha}}
\newunicodechar{Β}{\ensuremath{\beta}}
\newunicodechar{Γ}{\ensuremath{\Gamma}}
\newunicodechar{Δ}{\ensuremath{\Delta}}
\newunicodechar{Ε}{\ensuremath{\varepsilon}}
\newunicodechar{Θ}{\ensuremath{\Theta}}
\newunicodechar{Λ}{\ensuremath{\Lambda}}
\newunicodechar{Μ}{\ensuremath{\mu}}
\newunicodechar{Τ}{\ensuremath{\tau}}
\newunicodechar{Φ}{\ensuremath{\Phi}}
\newunicodechar{Ψ}{\ensuremath{\Psi}}
\newunicodechar{Ω}{\ensuremath{\Omega}}
\newunicodechar{↦}{\ensuremath{\mapsto}}
\newunicodechar{∈}{\ensuremath{\in}}
\newunicodechar{∉}{\ensuremath{\notin}}
\newunicodechar{∧}{\ensuremath{\wedge}}
\newunicodechar{⇔}{\ensuremath{\Leftrightarrow}}
\newunicodechar{⟺}{\ensuremath{\Longleftrightarrow}}
\newunicodechar{⇒}{\ensuremath{\Rightarrow}}
\newunicodechar{⟹}{\ensuremath{\Longrightarrow}}
\newunicodechar{∘}{\ensuremath{\circ}}
\newunicodechar{∪}{\ensuremath{\cup}}
\newunicodechar{∩}{\ensuremath{\cap}}
\newunicodechar{≡}{\ensuremath{\equiv}}
\newunicodechar{⊂}{\ensuremath{\subset}}
\newunicodechar{⊥}{\ensuremath{\perp}}
\newunicodechar{∃}{\ensuremath{\exists}}
\newunicodechar{∀}{\ensuremath{\forall}}
\newunicodechar{∛}{\ensuremath{\sqrt[3]{\,}}}
\newunicodechar{∜}{\ensuremath{\sqrt[4]{\,}}}
\newunicodechar{⃗}{\ensuremath{{}^{\to}}}
\newunicodechar{ⁿ}{\ifmmode^{n}\else\textsuperscript{\ensuremath{n}}\fi}
\newunicodechar{ˣ}{\ifmmode^{x}\else\textsuperscript{\ensuremath{x}}\fi}
\newunicodechar{ᵇ}{\ifmmode^{b}\else\textsuperscript{\ensuremath{b}}\fi}
\newunicodechar{ʳ}{\ifmmode^{r}\else\textsuperscript{\ensuremath{r}}\fi}
\newunicodechar{ᵘ}{\ifmmode^{u}\else\textsuperscript{\ensuremath{u}}\fi}
\newunicodechar{ʸ}{\ifmmode^{y}\else\textsuperscript{\ensuremath{y}}\fi}
\newunicodechar{ᵃ}{\ifmmode^{a}\else\textsuperscript{\ensuremath{a}}\fi}
\newunicodechar{ᵉ}{\ifmmode^{e}\else\textsuperscript{\ensuremath{e}}\fi}
\newunicodechar{ᵏ}{\ifmmode^{k}\else\textsuperscript{\ensuremath{k}}\fi}
\newunicodechar{ᵖ}{\ifmmode^{p}\else\textsuperscript{\ensuremath{p}}\fi}
\newunicodechar{ᴺ}{\ifmmode^{N}\else\textsuperscript{\ensuremath{N}}\fi}
\newunicodechar{ᶜ}{\ifmmode^{c}\else\textsuperscript{\ensuremath{c}}\fi}
\newunicodechar{ʰ}{\ifmmode^{h}\else\textsuperscript{\ensuremath{h}}\fi}
\newunicodechar{ᵠ}{\ifmmode^{q}\else\textsuperscript{\ensuremath{q}}\fi}
\newunicodechar{ₙ}{\ifmmode_{n}\else\textsubscript{\ensuremath{n}}\fi}
\newunicodechar{ₐ}{\ifmmode_{a}\else\textsubscript{\ensuremath{a}}\fi}
\newunicodechar{ᵢ}{\ifmmode_{i}\else\textsubscript{\ensuremath{i}}\fi}
\newunicodechar{ᵣ}{\ifmmode_{r}\else\textsubscript{\ensuremath{r}}\fi}
\newunicodechar{ₖ}{\ifmmode_{k}\else\textsubscript{\ensuremath{k}}\fi}
\newunicodechar{ᵦ}{\ifmmode_{\beta}\else\textsubscript{\ensuremath{\beta}}\fi}
\newunicodechar{ₓ}{\ifmmode_{x}\else\textsubscript{\ensuremath{x}}\fi}
\newfontfamily\popdisplay{ArchivoBlack-Regular.ttf}[Path=../../../fonts/]
\newfontfamily\poplogo{SpaceGrotesk-Bold.ttf}[Path=../../../fonts/]
\newfontfamily\popsemi{PlusJakartaSans-SemiBold.ttf}[Path=../../../fonts/]
\newfontfamily\popxbold{PlusJakartaSans-ExtraBold.ttf}[Path=../../../fonts/]
\newfontfamily\popxboldls{PlusJakartaSans-ExtraBold.ttf}[Path=../../../fonts/,LetterSpace=3.0]

\setlist[enumerate]{leftmargin=1.7em,itemsep=5pt,topsep=4pt}
\setlist[itemize]{leftmargin=1.7em,itemsep=4pt,topsep=4pt,label={\color{poporange}\textbullet}}
\setlength{\parindent}{0pt}
\setlength{\parskip}{5pt}
\renewcommand{\arraystretch}{1.25}

% pgfplots : style Pop par défaut
\pgfplotsset{
  every axis/.append style={axis line style={popink,line width=1pt},tick style={popink},
    grid style={popline,line width=0.5pt},label style={font=\small},tick label style={font=\footnotesize}},
  popcurve/.style={poporange,line width=1.8pt,no markers,smooth},
}
% Même style disponible dans un \draw TikZ simple (hors pgfplots)
\tikzset{popcurve/.style={poporange,line width=1.8pt}}
\tikzset{popgrid/.style={popline,very thin}}
\colorlet{popcurve}{poporange} % le nom du style est parfois utilisé comme couleur

% ============================================================
% Pill / squiggle
% ============================================================
\newcommand{\poppill}[3]{%
  \tikz[baseline=-0.55ex]\node[draw=popink,line width=1.2pt,rounded corners=2.6pt,%
    fill=#2,inner xsep=6.5pt,inner ysep=3pt]{%
    \popxboldls\fontsize{7}{7}\selectfont\color{#3}\MakeUppercase{#1}};%
}
\newcommand{\popsquiggle}[1]{%
  \tikz[baseline=-0.4ex]{
    \draw[#1,line width=2.6pt,line cap=round]
      plot [smooth] coordinates {(0,0.09) (0.34,0.01) (0.68,0.21) (1.02,0.01) (1.36,0.10)};
  }%
}
% Mot-clé surligné (marqueur orange sous le texte)
\newcommand{\cle}[1]{{\popxbold\color{popdark}#1}}

% ============================================================
% En-tête / pied
% ============================================================
\pagestyle{fancy}
\fancyhf{}
\renewcommand{\headrulewidth}{0pt}
\renewcommand{\footrulewidth}{0pt}
\lhead{\raisebox{-0.15ex}{\poplogo\fontsize{13}{13}\selectfont\color{popink}OnBuch\color{poporange}+}}
\rhead{\poppill{\DOCMATIERE\ \textbullet\ \DOCNIVEAU\ \textbullet\ Leçon \DOCLECON}{white}{popink}}
\lfoot{\popxbold\fontsize{7.6}{7.6}\selectfont\color{popmuted}\MakeUppercase{Cours signé OnBuch+}\ \textbullet\ {\popsemi\DOCTITRE}}
\rfoot{\tikz[baseline=-0.6ex]\node[rounded corners=8.5pt,fill=popink,minimum height=17pt,minimum width=17pt,inner xsep=5pt,inner sep=0]{\popxbold\fontsize{8}{8}\selectfont\color{popcream}\thepage\,/\,\pageref*{LastPage}};}

% ============================================================
% Titres
% ============================================================
\newcommand{\chaptitre}[2]{%
  \begin{tcolorbox}[enhanced,colback=poporange,colframe=popink,boxrule=2.6pt,arc=11pt,
      drop shadow=popink,left=15pt,right=15pt,top=12pt,bottom=11pt,before skip=3pt,after skip=18pt]
    \poppill{#2}{popink}{popcream}\par\vspace{9pt}
    {\raggedright\hyphenpenalty=10000\exhyphenpenalty=10000\popdisplay\fontsize{22}{24}\selectfont\color{popcream}\MakeUppercase{#1}\par}
  \end{tcolorbox}
}
% Section numérotée : pastille orange avec le numéro + titre Archivo
\newcounter{popsec}
\newcommand{\coursec}[1]{%
  \stepcounter{popsec}\setcounter{popsub}{0}%
  \par\vspace{16pt}\noindent
  \begin{minipage}{\linewidth}
  \tikz[baseline=-0.7ex]\node[draw=popink,line width=1.6pt,fill=poporange,rounded corners=4pt,minimum size=22pt,inner sep=0]{\popdisplay\fontsize{12}{12}\selectfont\color{popcream}\thepopsec};%
  \hspace{8pt}{\popdisplay\fontsize{15}{17}\selectfont\color{popink}\MakeUppercase{#1}}\par
  \vspace{4pt}\noindent{\color{poporange}\rule{36pt}{3.6pt}}
  \end{minipage}\par\nopagebreak\vspace{8pt}
}
\newcounter{popsub}
\newcommand{\courssub}[1]{%
  \stepcounter{popsub}%
  \par\vspace{9pt}\noindent{\popxbold\fontsize{11.5}{13}\selectfont\color{popink}{\color{poporange}\thepopsec.\thepopsub}\hspace{6pt}#1}\par\nopagebreak\vspace{4pt}
}

% ============================================================
% Boîtes pédagogiques — même recette : fond blanc, bordure épaisse colorée,
% ombre dure encre, badge plein. Titre optionnel [#1].
% ============================================================
\tcbset{popbase/.style={breakable,enhanced,colback=white,boxrule=2.3pt,arc=9pt,
  shadow={3pt}{-3pt}{0pt}{popink},left=14pt,right=14pt,top=11pt,bottom=11pt,before skip=11pt,after skip=15pt,
  fontupper=\popsemi\fontsize{10.6}{14.5}\selectfont\color{popink},
  fontlower=\popsemi\fontsize{10.6}{14.5}\selectfont\color{popink}}}
% Coche dessinée (le glyphe ✓ n'existe pas dans les polices du document)
\newcommand{\popcheck}[1]{\tikz[baseline=-0.5ex]\draw[#1,line width=1.6pt,line cap=round,line join=round](-0.13,0.0)--(-0.04,-0.09)--(0.13,0.1);}
\providecommand{\coche}{\popcheck{popgreen}}
\providecommand{\resultat}[1]{\par\smallskip\noindent\fcolorbox{popgreen}{popgreenL}{\parbox{\dimexpr\linewidth-2\fboxsep-2\fboxrule\relax}{\textbf{Résultat.} #1}}\par\smallskip}
\newcommand{\popboxhead}[3]{\poppill{#1}{#2}{white}\ifx\relax#3\relax\else\hspace{6pt}{\popxbold\fontsize{10.6}{12}\selectfont\color{popink}#3}\fi\par\vspace{7pt}}

\newtcolorbox{aretenir}[1][]{popbase,colframe=popgreen,before upper={\popboxhead{À retenir}{popgreen}{#1}}}
\newtcolorbox{exemplebox}[1][]{popbase,colframe=poporange,before upper={\popboxhead{Exemple}{poporange}{#1}}}
\newtcolorbox{definition}[1][]{popbase,colframe=popblue,before upper={\popboxhead{Définition}{popblue}{#1}}}
\newtcolorbox{propriete}[1][]{popbase,colframe=poppurple,before upper={\popboxhead{Propriété}{poppurple}{#1}}}
\newtcolorbox{methode}[1][]{popbase,colframe=popgold,before upper={\popboxhead{Méthode}{popgold}{#1}}}
\newtcolorbox{attention}[1][]{popbase,colframe=poppink,before upper={\popboxhead{Attention}{poppink}{#1}}}
\newtcolorbox{experience}[1][]{popbase,colframe=popblue,colback=popblueL,before upper={\popboxhead{Expérience}{popblue}{#1}}}
% « Le savais-tu ? » : carte violette pleine (contexte, culture, Cameroun)
\newtcolorbox{savaistu}[1][]{popbase,colback=poppurple,colframe=popink,
  fontupper=\popsemi\fontsize{10.4}{14.2}\selectfont\color{white},
  before upper={\poppill{Le savais-tu\,?}{white}{poppurple}\ifx\relax#1\relax\else\hspace{6pt}{\popxbold\color{white}#1}\fi\par\vspace{7pt}}}
% Exercice résolu : énoncé en haut, \tcblower puis la solution sur fond crème
\newcounter{popexo}\newcounter{popexr}
\newtcolorbox{exoresolu}[1][]{popbase,colframe=poporange,colbacklower=popcream,
  segmentation style={popink,line width=1.2pt,dashed},
  before upper={\stepcounter{popexr}\popboxhead{Exercice résolu \thepopexr}{poporange}{#1}},
  before lower={\poppill{Solution}{popink}{popcream}\par\vspace{6pt}}}
% Exercice (non résolu en place) — la correction est dans la partie Corrigés
\newtcolorbox{exercice}[1][]{popbase,colframe=popink,
  before upper={\stepcounter{popexo}\popboxhead{Exercice \thepopexo}{popink}{#1}}}
% Corrigé
\newtcolorbox{corrige}[1][]{popbase,colframe=popgreen,colback=popgreenL,
  before upper={\popboxhead{Corrigé}{popgreen}{#1}}}
% Objectifs / prérequis / bilan
\newtcolorbox{objectifs}{popbase,colframe=popink,colback=poporangeL,
  before upper={\popboxhead{Objectifs de la leçon}{popink}{}}}
\newtcolorbox{prerequis}{popbase,colframe=popink,colback=popblueL,
  before upper={\popboxhead{Prérequis}{popblue}{}}}
\newtcolorbox{bilan}{popbase,colframe=popink,colback=white,boxrule=2.8pt,
  before upper={\popboxhead{Fiche bilan}{poporange}{L'essentiel à connaître pour l'examen}}}
% Figure encadrée avec légende
\newtcolorbox{popfigure}[1][]{enhanced,breakable=false,colback=white,colframe=popline,boxrule=1.4pt,arc=8pt,
  left=10pt,right=10pt,top=10pt,bottom=8pt,before skip=10pt,after skip=12pt,halign=center,
  lower separated=false,halign lower=center,
  fontlower=\popsemi\fontsize{9}{11.5}\selectfont\color{popmuted},
  #1}
\newcommand{\legende}[1]{\par\vspace{6pt}{\popsemi\fontsize{9}{11.5}\selectfont\color{popmuted}#1}}

% ============================================================
% Signature OnBuch+
% ============================================================
\newcommand{\poplogoinline}[1]{{\poplogo\fontsize{#1}{#1}\selectfont\color{popink}OnBuch\color{poporange}+}}
\newcommand{\signatureonbuch}{%
  \begin{tcolorbox}[enhanced,colback=popink,colframe=popink,boxrule=2.2pt,arc=10pt,
      drop shadow=poporange,left=14pt,right=14pt,top=10pt,bottom=10pt,before skip=0pt,after skip=0pt]
    \tikz[baseline=-0.7ex]\node[circle,fill=popgreen,draw=popcream,line width=1.3pt,minimum size=22pt,inner sep=0]{\popcheck{white}};%
    \hspace{9pt}\begin{minipage}[c]{0.62\linewidth}
      {\popxbold\fontsize{12}{13}\selectfont\color{popcream}Signé\ \textbullet\ {\poplogo OnBuch\color{poporange}+}}\par\vspace{2pt}
      {\raggedright\popsemi\fontsize{8}{10.5}\selectfont\color{popline}\MakeUppercase{Équipe pédagogique OnBuch+}\\\MakeUppercase{Conforme au programme officiel MINESEC}\par}
    \end{minipage}\hfill
    {\poplogo\fontsize{22}{22}\selectfont\color{popcream}OnBuch\color{poporange}+}
  \end{tcolorbox}%
}

% ============================================================
% Page de garde
% #1 titre · #2 matière · #3 niveau · #4 module · #5 n° leçon · #6 durée ·
% #7 description / objectifs (texte ou itemize)
% ============================================================
\newcommand{\pagegarde}[7]{%
  \thispagestyle{empty}
  \newgeometry{a4paper,margin=1.7cm,top=1.6cm,bottom=1.4cm}%
  \begin{tikzpicture}[remember picture,overlay]
    \fill[poporange!18] ([xshift=-3.2cm,yshift=-3.6cm]current page.north east) circle (4.6cm);
    \fill[poppurple!14] ([xshift=2.4cm,yshift=7.6cm]current page.south west) circle (3.8cm);
  \end{tikzpicture}%
  {\poplogo\fontsize{30}{30}\selectfont\color{popink}OnBuch\color{poporange}+}\hfill
  \raisebox{0.6ex}{\poppill{Cours signé \textbullet\ Premium}{popink}{popcream}}\par
  \vspace{1pt}\popsquiggle{poppurple}\par
  \vspace{8pt}
  \begin{tcolorbox}[enhanced,colback=poporange,colframe=popink,boxrule=2.8pt,arc=13pt,
      drop shadow=popink,left=18pt,right=18pt,top=12pt,bottom=13pt,after skip=10pt]
    \poppill{#2 \textbullet\ #3}{popink}{popcream}\hspace{5pt}\poppill{#4}{white}{popink}\par\vspace{8pt}
    {\popdisplay\fontsize{14}{14}\selectfont\color{popink}LEÇON #5}\par\vspace{4pt}
    {\raggedright\hyphenpenalty=10000\exhyphenpenalty=10000\popdisplay\fontsize{25}{27}\selectfont\color{popcream}\MakeUppercase{#1}\par}
    \vspace{8pt}
    {\popxbold\fontsize{9.5}{9.5}\selectfont\color{popink}\MakeUppercase{Durée conseillée : #6}}
  \end{tcolorbox}
  \begin{tcolorbox}[enhanced,colback=white,colframe=popink,boxrule=2.2pt,arc=10pt,
      drop shadow=popink,left=16pt,right=16pt,top=10pt,bottom=10pt,after skip=8pt]
    \poppill{Dans cette leçon}{poppurple}{white}\par\vspace{5pt}
    \popsemi\fontsize{9.3}{12.6}\selectfont\color{popink}#7
  \end{tcolorbox}
  \noindent\begin{minipage}[t]{0.487\linewidth}
    \begin{tcolorbox}[enhanced,colback=popgreen,colframe=popink,boxrule=2.2pt,arc=10pt,
        drop shadow=popink,left=11pt,right=11pt,top=10pt,bottom=10pt,height=3.1cm,valign=center]
      \tcbox[colback=white,colframe=popink,boxrule=1.3pt,arc=5pt,boxsep=0pt,nobeforeafter,
        left=3pt,right=3pt,top=3pt,bottom=3pt,valign=center]{\qrcode[height=1.9cm]{https://play.google.com/store/apps/details?id=com.luvvix.onbuch&hl=fr}}%
      \hspace{8pt}\begin{minipage}[c]{0.5\linewidth}
        {\popdisplay\fontsize{11}{12.5}\selectfont\color{white}\MakeUppercase{Télécharge OnBuch}}\par\vspace{4pt}
        {\raggedright\popsemi\fontsize{8.4}{11}\selectfont\color{white}Gratuit \textbullet\ Google Play\\Tuteur IA, annales, concours, résultats\par}
      \end{minipage}
    \end{tcolorbox}
  \end{minipage}\hfill
  \begin{minipage}[t]{0.487\linewidth}
    \begin{tcolorbox}[enhanced,colback=poppurple,colframe=popink,boxrule=2.2pt,arc=10pt,
        drop shadow=popink,left=11pt,right=11pt,top=10pt,bottom=10pt,height=3.1cm,valign=center]
      \tikz[baseline=-0.7ex]\node[circle,fill=white,draw=popink,line width=1.3pt,minimum size=1.6cm,inner sep=0]{\popdisplay\fontsize{24}{24}\selectfont\color{poppurple}+};%
      \hspace{8pt}\begin{minipage}[c]{0.6\linewidth}
        {\popdisplay\fontsize{11}{12.5}\selectfont\color{white}\MakeUppercase{Passe à OnBuch+}}\par\vspace{4pt}
        {\raggedright\popsemi\fontsize{8.4}{11}\selectfont\color{white}Tous les cours signés, Tuteur IA illimité, fiches PDF — onglet OnBuch+ de l'app\par}
      \end{minipage}
    \end{tcolorbox}
  \end{minipage}
  \vspace{8pt}
  \popcontenu
  \vfill
  \signatureonbuch
  \restoregeometry
  \newpage
}

% Rangée « Ce cours contient » (page de garde)
\newcommand{\poptile}[3]{% #1 couleur #2 gros texte #3 légende
  \begin{tcolorbox}[enhanced,colback=#1,colframe=popink,boxrule=2pt,arc=9pt,shadow={2.5pt}{-2.5pt}{0pt}{popink},
      left=6pt,right=6pt,top=9pt,bottom=9pt,halign=center,valign=center,height=1.75cm,width=\linewidth,nobeforeafter]
    {\popdisplay\fontsize{12.5}{13}\selectfont\color{white}\MakeUppercase{#2}}\par\vspace{3pt}
    {\centering\popsemi\fontsize{7.8}{9.5}\selectfont\color{white}#3\par}
  \end{tcolorbox}}
\newcommand{\popcontenu}{%
  \par\noindent{\popxboldls\fontsize{7.5}{7.5}\selectfont\color{popmuted}CE COURS SIGNÉ CONTIENT}\par\vspace{7pt}
  \noindent\begin{minipage}[t]{0.235\linewidth}\poptile{popblue}{Cours}{complet, illustré, pas à pas}\end{minipage}\hfill
  \begin{minipage}[t]{0.235\linewidth}\poptile{poporange}{Méthodes}{et exercices résolus}\end{minipage}\hfill
  \begin{minipage}[t]{0.235\linewidth}\poptile{poppink}{Exercices}{type examen + corrigés}\end{minipage}\hfill
  \begin{minipage}[t]{0.235\linewidth}\poptile{popgreen}{Fiche bilan}{l'essentiel à retenir}\end{minipage}\par}

% Sommaire
\newtcolorbox{sommaire}{enhanced,colback=white,colframe=popink,boxrule=2.2pt,arc=10pt,
  drop shadow=popink,left=18pt,right=18pt,top=13pt,bottom=13pt,before skip=6pt,after skip=20pt,
  before upper={\poppill{Sommaire}{poppurple}{white}\par\vspace{9pt}\popsemi\fontsize{10.6}{14.5}\selectfont\color{popink}}}

% Signature de fin de cours — poussée tout en bas de la dernière page
\newcommand{\findecours}{%
  \par\vfill\nopagebreak
  \begin{center}\popsquiggle{poporange}\end{center}\vspace{4pt}
  \signatureonbuch
}
\AtBeginDocument{\def\llbracket{\mathopen{[\mkern-3.5mu[}}\def\rrbracket{\mathclose{]\mkern-3.5mu]}}} % intervalles d'entiers (glyphe ⟦ absent de la police)
% Glyphes absents de Fira Math : redéfinis à partir de glyphes disponibles
\AtBeginDocument{%
  \def\setminus{\mathbin{\backslash}}\let\smallsetminus\setminus
  \def\vdots{\mathord{\vbox{\baselineskip3.2pt\lineskiplimit0pt\kern4pt\hbox{.}\hbox{.}\hbox{.}}}}%
  \def\ddots{\mathinner{\mkern1mu\raise6.5pt\hbox{.}\mkern2mu\raise3.6pt\hbox{.}\mkern2mu\raise0.7pt\hbox{.}\mkern1mu}}%
  \def\triangle{\mathord{\scalebox{0.9}{$\symup{\Delta}$}}}%
  \def\nabla{\mathord{\raisebox{\depth}{\scalebox{1}[-1]{$\symup{\Delta}$}}}}%
  \def\checkmark{\popcheck{popdark}}%
  \def\star{\mathbin{\ast}}\def\bigstar{\mathord{\ast}}%
  \def\diamond{\mathbin{\scalebox{0.6}{$\Diamond$}}}%
  \def\bigcup{\mathop{\mathchoice{\raisebox{-0.3em}{\scalebox{1.7}{$\cup$}}}{\scalebox{1.25}{$\cup$}}{\cup}{\cup}}\displaylimits}%
  \def\bigcap{\mathop{\mathchoice{\raisebox{-0.3em}{\scalebox{1.7}{$\cap$}}}{\scalebox{1.25}{$\cap$}}{\cap}{\cap}}\displaylimits}%
  \def\lesseqqgtr{\mathrel{\lessgtr}}\def\bigoplus{\mathop{\oplus}\displaylimits}%
}
\providecommand{\ssi}{si et seulement si} % abréviation fréquente
\AtBeginDocument{\providecommand{\wideparen}{\overparen}} % arc de cercle

%% FIGFIX-BEGIN v5 — correctifs globaux des figures (docs/onbuch-generation/tools/figfix)
\usepackage{adjustbox}\usepackage{etoolbox}\tcbuselibrary{raster}
% toute figure TikZ/pgfplots plus large que la ligne est réduite à la largeur disponible
\BeforeBeginEnvironment{tikzpicture}{\begin{adjustbox}{max width=\linewidth}}
\AfterEndEnvironment{tikzpicture}{\end{adjustbox}}
% légendes pgfplots lisibles par-dessus les courbes
\pgfplotsset{every axis/.append style={legend style={fill=white,fill opacity=0.92,text opacity=1,draw=black!15,inner sep=3pt}}}
% étiquettes d'arêtes (syntaxe edge["..."]) avec fond blanc
\tikzset{every edge quotes/.append style={fill=white,inner sep=1.2pt}}
%% FIGFIX-END
\begin{document}

\pagegarde{Recherche de nouvelles sources d'énergie}{SVT}{Tle D}{Module IV : La Biotechnologie — Énergies renouvelables}{N19}{4 h}{Cette leçon explore comment la biomasse (végétaux, déchets organiques) peut être transformée en biocarburants par des procédés biotechnologiques. Elle compare ces énergies renouvelables aux énergies fossiles et évalue leur intérêt face aux déficits énergétiques et à la pollution au Cameroun.
\vspace{4pt}
\begin{multicols}{2}\small
\begin{itemize}
\item À la fin de cette leçon, je sais définir la biomasse et un biocarburant.
\item À la fin de cette leçon, je sais expliquer le principe biologique de la fermentation alcoolique à partir de résultats expérimentaux.
\item À la fin de cette leçon, je sais décrire les étapes de production du bioéthanol et du biodiesel.
\item À la fin de cette leçon, je sais citer les principales matières premières mobilisables au Cameroun (canne à sucre, manioc, huile de palme).
\item À la fin de cette leçon, je sais expliquer la production de biogaz à partir des déchets organiques.
\item À la fin de cette leçon, je sais comparer les avantages et les limites des biocarburants par rapport aux énergies fossiles.
\end{itemize}\end{multicols}}

\begin{sommaire}
\begin{multicols}{2}
\begin{enumerate}
\item Biomasse, biocarburants et contexte énergétique
\item La fermentation alcoolique : fondement expérimental du bioéthanol
\item Production du bioéthanol et matières premières camerounaises
\item Biodiesel et biogaz : les autres biocarburants
\item Biocarburants vs énergies fossiles : avantages et limites
\item Énergies renouvelables, éducation et citoyenneté
\item Activité d'intégration
\item Méthodes et exercices résolus
\item Exercices
\item Corrigés des exercices
\item Fiche bilan
\end{enumerate}
\end{multicols}
\end{sommaire}

\begin{objectifs}
\begin{itemize}
\item À la fin de cette leçon, je sais définir la biomasse et un biocarburant.
\item À la fin de cette leçon, je sais expliquer le principe biologique de la fermentation alcoolique à partir de résultats expérimentaux.
\item À la fin de cette leçon, je sais décrire les étapes de production du bioéthanol et du biodiesel.
\item À la fin de cette leçon, je sais citer les principales matières premières mobilisables au Cameroun (canne à sucre, manioc, huile de palme).
\item À la fin de cette leçon, je sais expliquer la production de biogaz à partir des déchets organiques.
\item À la fin de cette leçon, je sais comparer les avantages et les limites des biocarburants par rapport aux énergies fossiles.
\item À la fin de cette leçon, je sais interpréter des documents (tableaux, courbes, bilans) portant sur la production d'énergie renouvelable.
\item À la fin de cette leçon, je sais construire un argumentaire pour éduquer et informer sur l'utilisation des énergies renouvelables.
\end{itemize}
\end{objectifs}

\begin{prerequis}
\begin{itemize}
\item La photosynthèse et la fixation du carbone par les végétaux (Première)
\item La respiration cellulaire et la fermentation (notion de dégradation du glucose)
\item Les glucides : sucres simples, saccharose, amidon (structure et hydrolyse)
\item Les lipides : triglycérides et esters (notions de chimie organique)
\item La notion d'énergie et les combustibles fossiles (charbon, pétrole)
\end{itemize}
\end{prerequis}

\newpage

\coursec{Biomasse, biocarburants et contexte énergétique}

\textbf{Situation d'accroche.} À Bafoussam comme à Yaoundé, les délestages d'ENEO rythment la vie des quartiers : les élèves révisent leurs leçons de SVT à la lampe-tempête, les commerces font tourner des groupes électrogènes coûteux et polluants, et les ateliers soudent au ralenti. Pourtant, le Cameroun produit chaque année des millions de tonnes de canne à sucre, de manioc et d'huile de palme, ainsi que d'énormes quantités de déchets organiques : fanes de manioc, bagasse de canne, fumiers, ordures ménagères. Et si ces ressources végétales pouvaient remplacer une partie du pétrole importé ? Comment transformer une plante ou un déchet en carburant ? Ces « biocarburants » sont-ils vraiment une solution durable aux déficits énergétiques du pays ?

\textbf{Problématique.} Pour répondre à ces questions, nous devons d'abord définir rigoureusement ce qu'est la biomasse et ce qu'est un biocarburant, comprendre d'où vient le carbone qu'ils contiennent, et distinguer les énergies renouvelables des énergies fossiles. Nous situerons ensuite ces notions dans le contexte énergétique camerounais, avant de dresser la carte des principaux types de biocarburants que cette leçon étudiera.

\courssub{La biomasse : une énergie stockée dans la matière vivante}

\textbf{Observation.} Dans une concession agricole de la région de l'Ouest, on brûle après la récolte les tiges de maïs, les fanes de manioc et la paille de riz. Cette combustion dégage de la chaleur : la matière végétale contient donc de l'énergie utilisable. De même, le bois de chauffe, qui fournit l'essentiel de l'énergie de cuisson des ménages camerounais, n'est rien d'autre que de la matière organique végétale. D'où vient cette énergie stockée dans les végétaux ?

\begin{definition}[Biomasse]
La \cle{biomasse} est l'ensemble de la matière organique d'origine végétale ou animale, utilisable comme source d'énergie. Elle comprend notamment :
\begin{itemize}
\item le \textbf{bois} et ses dérivés (bois de chauffe, charbon de bois, sciures) ;
\item les \textbf{résidus de récolte} (pailles, fanes, tiges, bagasse de canne à sucre, balles de riz) ;
\item les \textbf{déchets organiques} (ordures ménagères fermentescibles, boues de stations d'épuration) ;
\item les \textbf{effluents d'élevage} (fumiers, lisiers).
\end{itemize}
\end{definition}

L'énergie contenue dans la biomasse est de l'\textbf{énergie chimique}, stockée dans les liaisons des molécules organiques (glucides, lipides, protéines). Or ces molécules ont été fabriquées par les êtres vivants à partir de matière minérale. Chez les végétaux chlorophylliens, c'est la \cle{photosynthèse} qui réalise cette conversion : l'énergie lumineuse du Soleil est captée par la chlorophylle et utilisée pour fabriquer de la matière organique à partir du dioxyde de carbone atmosphérique et de l'eau, selon le bilan global :

\[ \ce{6CO2 + 6H2O ->[$\text{lumière}$][\text{chlorophylle}] C6H12O6 + 6O2} \]

\begin{aretenir}[L'origine solaire de l'énergie de la biomasse]
La biomasse est une forme de \textbf{stockage de l'énergie solaire} : la photosynthèse transforme l'énergie lumineuse en énergie chimique, fixée dans la matière organique. Brûler ou transformer la biomasse, c'est restituer une énergie solaire récemment captée.
\end{aretenir}

\courssub{Les biocarburants : des carburants issus du vivant}

\textbf{Observation.} Dans certains pays (Brésil, États-Unis), les stations-service distribuent de l'essence mélangée à de l'éthanol produit à partir de canne à sucre ou de maïs. Des bus urbains européens roulent au « Diester », un carburant dérivé de l'huile de colza. Ces carburants ne proviennent pas du sous-sol : ils sont fabriqués à partir de végétaux cultivés.

\begin{definition}[Biocarburant]
Un \cle{biocarburant} (ou agrocarburant de 1\textsuperscript{re} génération lorsqu'il utilise des cultures alimentaires) est un carburant produit à partir de la biomasse, utilisable pur ou en mélange avec un carburant fossile (essence, gazole). Les trois grandes familles de biocarburants sont :
\begin{itemize}
\item le \textbf{bioéthanol}, alcool obtenu par fermentation de sucres ou d'amidon, mélangé à l'essence ;
\item le \textbf{biodiesel}, ester méthylique d'huile végétale (palme, soja, jatropha), mélangé au gazole ;
\item le \textbf{biogaz}, mélange gazeux riche en méthane \ce{CH4} obtenu par fermentation des déchets organiques en l'absence d'air (méthanisation).
\end{itemize}
\end{definition}

\begin{exemplebox}[Deux échelles de temps incomparables]
Comparons les durées de « fabrication » de deux carburants :
\begin{itemize}
\item le \textbf{pétrole} s'est formé par transformation de plancton et de matière organique enfouis dans les sédiments, sous l'effet de la pression et de la température, sur une durée de l'ordre de \textbf{100 millions d'années} ;
\item la \textbf{canne à sucre}, matière première du bioéthanol, boucle son cycle cultural en \textbf{12 à 18 mois} environ sous climat tropical.
\end{itemize}
Le rapport des durées est d'environ $10^{8}$ : ce que la nature a mis cent millions d'années à fabriquer, l'humanité le consomme en quelques siècles, alors qu'une plantation de canne se renouvelle à chaque campagne agricole.
\end{exemplebox}

\courssub{Énergies renouvelables et énergies fossiles : une question de temps}

L'exemple précédent éclaire la distinction fondamentale entre les deux grandes catégories de sources d'énergie.

\begin{definition}[Énergies fossiles et énergies renouvelables]
\begin{itemize}
\item Une \cle{énergie fossile} (pétrole, charbon, gaz naturel) provient de la transformation de matière organique enfouie depuis des \textbf{millions d'années}. Son stock sur Terre est limité et ne se reconstitue pas à l'échelle humaine : elle est \textbf{épuisable}.
\item Une \cle{énergie renouvelable} (solaire, hydraulique, éolien, biomasse) se régénère naturellement sur des durées courtes (quelques mois à quelques dizaines d'années pour la biomasse), c'est-à-dire à l'\textbf{échelle du temps humain}.
\end{itemize}
\end{definition}

\begin{attention}[Le critère décisif : l'échelle de temps]
Le critère de distinction n'est ni « propre/sale » ni « naturel/artificiel » : c'est le \textbf{temps de renouvellement} comparé au temps de consommation. Une ressource est renouvelable si sa vitesse de régénération est supérieure ou comparable à sa vitesse d'exploitation. Ainsi, même le bois, pourtant renouvelable, peut être surexploité : la déforestation autour des grandes villes camerounaises montre qu'une ressource renouvelable mal gérée s'épuise localement.
\end{attention}

\courssub{Le cycle du carbone : pourquoi dit-on que le bilan carbone d'un biocarburant est quasi neutre ?}

\begin{methode}[Démarche d'investigation : suivre un atome de carbone]
\textbf{Question.} D'où vient le carbone présent dans l'éthanol \ce{C2H5OH} produit à partir de la canne à sucre, et que devient-il lors de l'utilisation du biocarburant ?
\begin{enumerate}
\item \textbf{Observation / Hypothèse.} La canne à sucre est un végétal chlorophyllien : elle réalise la photosynthèse. Elle capte le \ce{CO2} atmosphérique pour synthétiser du saccharose (sucre de transport).
\item \textbf{Expérience / Document.} La transformation industrielle (broyage, fermentation par levures, distillation) convertit le saccharose \ce{C12H22O11} en éthanol \ce{C2H5OH}. Le carbone du sucre se retrouve dans l'éthanol.
\item \textbf{Interprétation.} La combustion de l'éthanol dans un moteur suit l'équation : $\ce{C2H5OH + 3O2 -> 2CO2 + 3H2O}$. Le carbone est rejeté sous forme de \ce{CO2}.
\item \textbf{Conclusion.} Le \ce{CO2} rejeté à l'étape 3 est \textbf{exactement celui qui a été capté} à l'étape 1, quelques mois plus tôt. Contrairement au pétrole, qui libère un carbone séquestré depuis des millions d'années (apport \emph{net} de \ce{CO2} à l'atmosphère), le biocarburant fait circuler un carbone déjà présent dans le cycle court de la biosphère.
\end{enumerate}
\end{methode}

\begin{popfigure}
\begin{center}
\begin{tikzpicture}[
  box/.style={draw, rounded corners, minimum width=3.2cm, minimum height=1.2cm, align=center, font=\small},
  arrow/.style={->, very thick, >=Stealth, font=\tiny, align=center}
]
\matrix[column sep=1.5cm, row sep=1.2cm] {
  \node[box, draw=popblue, fill=popblueL] (co2) {\ce{CO2} atmosphérique}; &
  \node[box, draw=popgreen, fill=popgreenL, align=center] (bio){Biomasse\\ (plantes, déchets)}; \\
  \node[box, draw=poppurple, fill=poppurpleL, align=center] (comb){Combustion\\ (moteur, chaudière)}; &
  \node[box, draw=poporange, fill=poporangeL, align=center] (carb){Biocarburant\\ (éthanol, biodiesel, biogaz)}; \\
};
\draw[arrow, popgreen] (co2) -- node[above, align=center]{1. Photosynthèse\\(fixation du \ce{CO2})} (bio);
\draw[arrow, poporange] (bio) -- node[right, align=center]{2. Transformation\\(fermentation, estérification,\\ méthanisation)} (carb);
\draw[arrow, poppurple] (carb) -- node[below, align=center]{3. Combustion\\(énergie libérée)} (comb);
\draw[arrow, popblue] (comb) -- node[left, align=center]{4. Rejet de \ce{CO2}\\(carbone récemment capté)} (co2);
\end{tikzpicture}
\end{center}
\legende{Cycle du carbone simplifié associé à un biocarburant. Le \ce{CO2} rejeté lors de la combustion (étape 4) correspond au \ce{CO2} fixé par la plante lors de la photosynthèse (étape 1) : le cycle se boucle en quelques mois, d'où un bilan carbone dit « quasi neutre ».}
\end{popfigure}

\begin{attention}[Erreur fréquente : un biocarburant n'est pas « sans \ce{CO2} »]
Il est faux d'écrire qu'un biocarburant « ne pollue pas » ou « n'émet pas de \ce{CO2} ». Toute combustion d'un composé organique \textbf{émet du \ce{CO2}}. La nuance essentielle est que ce \ce{CO2} a été \textbf{récemment capté} par la plante : il ne s'agit pas d'un carbone « nouveau » ajouté à l'atmosphère, contrairement au carbone fossile. On parle donc de \textbf{bilan carbone quasi neutre}, et non nul : il faut en effet ajouter les émissions liées aux engrais (protoxyde d'azote \ce{N2O}), aux machines agricoles, au transport et à la transformation industrielle, qui, elles, consomment souvent des énergies fossiles.
\end{attention}

\courssub{Le contexte énergétique camerounais : un paradoxe à résoudre}

\textbf{Document.} Le tableau ci-dessous rassemble des données indicatives sur la situation énergétique du Cameroun.

\begin{center}
\begin{tabularx}{\linewidth}{|l|X|}
\hline
\rowcolor{popblueL} \textbf{Constat} & \textbf{Donnée indicative} \\
\hline
Part de la biomasse traditionnelle (bois, charbon de bois) dans l'énergie consommée & environ 70 \% (cuisson, chauffage domestique) \\
\hline
Part des produits pétroliers & environ 20 \% (transport, groupes électrogènes) \\
\hline
Part de l'électricité (surtout hydraulique) & environ 10 \% \\
\hline
Accès à l'électricité en milieu rural & inférieur à 25 \% des ménages \\
\hline
Production agricole mobilisable pour biocarburants & canne à sucre (SOSUCAM), huile de palme (SOCAPALM, CDC), manioc (première racine cultivée du pays) \\
\hline
\end{tabularx}
\end{center}

\textbf{Interprétation.} Trois faits ressortent. D'une part, le pays connaît des \textbf{déficits énergétiques} chroniques : la demande croît plus vite que la production, d'où les délestages. D'autre part, le Cameroun, bien que producteur de pétrole brut, \textbf{importe une grande partie de ses carburants raffinés}, ce qui pèse lourdement sur la balance commerciale. Enfin, le pays dispose d'un \textbf{potentiel agricole considérable} et de masses importantes de déchets organiques urbains et ruraux, encore peu valorisés. La biomasse constitue donc déjà la première source d'énergie du pays, mais sous une forme traditionnelle (bois de chauffe) peu efficace et source de déforestation : la moderniser sous forme de biocarburants est un enjeu majeur.

\begin{savaistu}[Pétrole brut et carburants importés]
Le Cameroun extrait du pétrole brut au large de Kribi et dans le bassin du Logone, mais sa capacité de raffinage (la SONARA à Limbé, longtemps l'unique raffinerie du pays) est insuffisante et a été fortement réduite après l'incendie de 2019. Conséquence paradoxale : le pays exporte son brut et importe l'essence et le gazole raffinés à l'étranger, à des coûts élevés. Produire localement des biocarburants à partir de la canne, du manioc ou de l'huile de palme réduirait cette dépendance et créerait des emplois ruraux.
\end{savaistu}

\begin{popfigure}
\begin{center}
\begin{tikzpicture}
\begin{axis}[
  width=0.85\linewidth, height=7cm,
  ybar, bar width=1.8cm,
  ymin=0, ymax=80,
  ylabel={Part dans l'énergie consommée (\%)},
  symbolic x coords={Biomasse traditionnelle, Produits pétroliers, Électricité},
  xtick=data,
  nodes near coords={\pgfmathprintnumber\pgfplotspointmeta\,\%},
  every node near coord/.append style={font=\small, anchor=south},
  axis lines=left,
  enlarge x limits=0.4,
  tick label style={font=\small, align=center},
]
\addplot[fill=popgreenL, draw=popgreen] coordinates {(Biomasse traditionnelle,70)};
\addplot[fill=poporangeL, draw=poporange] coordinates {(Produits pétroliers,20)};
\addplot[fill=popblueL, draw=popblue] coordinates {(Électricité,10)};
\end{axis}
\end{tikzpicture}
\end{center}
\legende{Répartition indicative des sources d'énergie consommées au Cameroun : la biomasse traditionnelle (bois, charbon de bois) domine largement (environ 70 \%), devant les produits pétroliers (environ 20 \%) et l'électricité, essentiellement d'origine hydraulique (barrages de la Sanaga, environ 10 \%).}
\end{popfigure}

\courssub{Classification des biocarburants : la carte de la leçon}

Les biocarburants se classent selon la matière première utilisée et le procédé biologique ou chimique de transformation. Cette classification structure le reste de la leçon.

\begin{center}
\begin{tabularx}{\linewidth}{|l|X|X|X|}
\hline
\rowcolor{popblueL} \textbf{Biocarburant} & \textbf{Matière première principale} & \textbf{Procédé clé} & \textbf{Principales filières camerounaises} \\
\hline
Bioéthanol & Sucres (canne), amidon (manioc, maïs) & Fermentation alcoolique par levures (\emph{Saccharomyces}), puis distillation & Canne à sucre (SOSUCAM), Manioc, Maïs \\
\hline
Biodiesel & Huiles végétales (palme, jatropha, coton) & Transestérification de l'huile (réaction avec méthanol + catalyseur) & Huile de palme (SOCAPALM, CDC), Jatropha, Coton \\
\hline
Biogaz & Déchets organiques, effluents d'élevage, boues & Méthanisation (fermentation anaérobie par bactéries méthanogènes) & Fumiers, ordures ménagères, boues d'épuration, résidus agro-industriels \\
\hline
\end{tabularx}
\end{center}

\begin{exoresolu}[Renouvelable ou fossile ?]
\textbf{Énoncé.} Classer chacune des sources d'énergie suivantes en « fossile » ou « renouvelable », en justifiant par le temps de renouvellement : gaz naturel ; bois de chauffe ; bioéthanol de canne ; charbon ; hydroélectricité.
\tcblower
\textbf{Solution.}
\begin{itemize}
\item \textbf{Gaz naturel} : fossile. Il s'est formé, comme le pétrole, à partir de matière organique enfouie il y a des millions d'années ; son stock ne se reconstitue pas à l'échelle humaine.
\item \textbf{Bois de chauffe} : renouvelable. Un arbre de chauffe se régénère en quelques années à quelques dizaines d'années, à condition que le rythme de coupe n'excède pas le rythme de croissance (sinon déforestation).
\item \textbf{Bioéthanol de canne} : renouvelable. La canne à sucre se renouvelle à chaque cycle cultural, en 12 à 18 mois.
\item \textbf{Charbon} : fossile. Il résulte de la transformation de forêts enfouies depuis des centaines de millions d'années (ère primaire).
\item \textbf{Hydroélectricité} : renouvelable. Le cycle de l'eau, alimenté par l'énergie solaire, reconstitue en permanence la ressource (pluies, débits des fleuves comme la Sanaga).
\end{itemize}
\textbf{Point méthode.} Le seul critère à mobiliser est la comparaison entre le temps de renouvellement de la ressource et le temps de sa consommation par l'homme.
\end{exoresolu}

\begin{aretenir}[L'essentiel de la section]
\begin{itemize}
\item La \cle{biomasse} est la matière organique d'origine végétale ou animale utilisable comme source d'énergie ; elle stocke l'énergie solaire captée par la photosynthèse.
\item Un \cle{biocarburant} est un carburant produit à partir de la biomasse, pur ou en mélange avec un carburant fossile.
\item Une énergie est \cle{renouvelable} si elle se régénère à l'échelle du temps humain ; les énergies \cle{fossiles} demandent des millions d'années pour se former.
\item Le bilan carbone d'un biocarburant est \textbf{quasi neutre} : le \ce{CO2} émis à la combustion a été récemment capté par la plante — mais il n'est pas nul (intrants fossiles pour la culture et la transformation).
\item Le Cameroun, confronté aux déficits énergétiques et à l'importation de carburants raffinés, dispose d'un potentiel de biomasse considérable (canne, palme, manioc, déchets).
\item Les trois familles de biocarburants — \textbf{bioéthanol}, \textbf{biodiesel}, \textbf{biogaz} — feront l'objet des sections suivantes.
\end{itemize}
\end{aretenir}

\coursec{La fermentation alcoolique : fondement expérimental du bioéthanol}

Nous avons vu dans la section précédente que la biomasse — canne à sucre, manioc, résidus de palme — constitue un immense réservoir de matière organique que le Cameroun possède en abondance, et que les biocarburants proposent de convertir cette matière végétale en énergie liquide. Mais une question fondamentale demeure : \emph{comment} transforme-t-on concrètement le sucre d'une plante en un carburant capable de faire tourner un moteur ? La réponse ne tient ni à une machine, ni à un réacteur chimique sophistiqué : elle tient à un micro-organisme unicellulaire, la levure, et à un phénomène biologique que l'humanité exploite depuis des millénaires sans le comprendre : la \cle{fermentation alcoolique}. C'est elle qui constitue le fondement expérimental et industriel du bioéthanol. Étudions-la avec la rigueur d'une démarche scientifique complète.

\courssub{Observation historique : une transformation biologique vieille comme l'humanité}

Bien avant toute théorie microbienne, les hommes ont constaté qu'un jus de fruit, une bouillie de céréales ou une sève sucrée abandonnés à l'air se transforment spontanément en boisson alcoolisée. Le vin de palme, la bière de maïs ou de mil (le \og bil-bil \fg{} du Nord-Cameroun), l'alcool de canne (\og odontol \fg{}) en sont des illustrations locales bien connues. Dans tous les cas, le point de départ est identique : un milieu riche en \cle{sucre}. Et le point d'arrivée aussi : une boisson contenant de l'\cle{éthanol} et, souvent, un dégagement gazeux visible (les bulles de la bière en fermentation).

\begin{attention}[Ce que l'observation seule ne permet pas de conclure]
L'observation montre une \emph{corrélation} (sucre $\rightarrow$ alcool + gaz), mais n'identifie ni l'agent responsable, ni le mécanisme. Pendant des siècles, on a cru à une \og génération spontanée \fg{} ou à une simple réaction chimique. C'est Louis Pasteur, en 1857, qui démontre que la fermentation alcoolique est l'œuvre d'êtres vivants microscopiques : les \cle{levures} (champignons unicellulaires du genre \emph{Saccharomyces}). Pour établir rigoureusement le phénomène, il faut une expérience contrôlée.
\end{attention}

\courssub{Mise en évidence expérimentale de la fermentation alcoolique}

\begin{experience}[Production d'éthanol et de dioxyde de carbone par des levures]
\textbf{Matériel :} un ballon (ou flacon) muni d'un bouchon traversé par un tube à dégagement ; un tube à essais contenant de l'eau de chaux ; du jus de canne à sucre frais (ou une solution de glucose à \SI{100}{\gram\per\liter}) ; une suspension de levures de boulanger (\emph{Saccharomyces cerevisiae}) ; un bain-marie thermostaté à \SI{30}{\celsius}.

\textbf{Protocole :}
\begin{enumerate}
\item Verser le jus de canne dans le ballon et y ajouter la suspension de levures.
\item Boucher hermétiquement ; le tube à dégagement plonge dans l'eau de chaux du tube à essais.
\item Placer l'ensemble à \SI{30}{\celsius} et observer pendant 24 à 48 h.
\item Prélever un échantillon du liquide du ballon pour y rechercher l'éthanol (odeur caractéristique, dosage par distillation ou densimétrie).
\end{enumerate}

\textbf{Observations :} au bout de quelques heures, des bulles de gaz se dégagent du ballon et barbotent dans l'eau de chaux, qui se \emph{trouble} progressivement (dépôt blanc). Le liquide du ballon dégage une odeur d'alcool ; le dosage confirme la présence d'éthanol. Le taux de sucre, lui, diminue.

\textbf{Interprétation :} le trouble de l'eau de chaux prouve que le gaz dégagé est du \cle{dioxyde de carbone} (\ce{CO2}). Les levures ont donc transformé le sucre en éthanol et en \ce{CO2}.
\end{experience}

\begin{popfigure}
\begin{center}
\begin{tikzpicture}[scale=0.95, every node/.style={font=\small}]
% Ballon
\draw[thick, popblue] (-2.6,0) .. controls (-2.6,-1.8) and (2.6,-1.8) .. (2.6,0);
\draw[thick, popblue] (-2.6,0) -- (-2.6,1.1);
\draw[thick, popblue] (2.6,0) -- (2.6,1.1);
% Liquide dans le ballon
\fill[poporangeL] (-2.55,-0.05) .. controls (-2.55,-1.75) and (2.55,-1.75) .. (2.55,-0.05) -- cycle;
\draw[poporange, thick] (-2.55,-0.05) -- (2.55,-0.05);
% Bouchon
\fill[popmuted!60] (-0.55,1.1) rectangle (0.55,1.55);
% Tube à dégagement
\draw[thick, popdark] (0,1.1) -- (0,2.3) -- (4.6,2.3) -- (4.6,0.4);
% Bulles dans le ballon
\foreach \x/\y/\r in {-1.2/-0.9/0.09, -0.4/-1.2/0.07, 0.6/-0.8/0.10, 1.4/-1.1/0.07, 0.1/-0.5/0.06, -1.6/-0.5/0.06}
  \draw[popdark] (\x,\y) circle (\r);
% Tube à essais
\draw[thick, popblue] (3.9,-1.6) -- (3.9,0.9);
\draw[thick, popblue] (5.3,-1.6) -- (5.3,0.9);
\draw[thick, popblue] (3.9,-1.6) .. controls (3.9,-2.1) and (5.3,-2.1) .. (5.3,-1.6);
% Eau de chaux
\fill[popblueL] (3.95,-1.55) .. controls (3.95,-2.0) and (5.25,-2.0) .. (5.25,-1.55) -- (5.25,-0.3) -- (3.95,-0.3) -- cycle;
% Bulles dans l'eau de chaux
\foreach \y in {-1.2,-0.8,-0.45}
  \draw[popdark] (4.6,\y) circle (0.06);
% Annotations
\node[align=center, popdark] at (0,-2.6) {\textbf{Ballon} : jus de canne\\ + levures (\SI{30}{\celsius})};
\node[align=center, popdark] at (4.6,-2.6) {\textbf{Tube à essais} :\\ eau de chaux};
\node[popgreen, align=left] at (6.0,-0.9) {trouble blanc\\ $\Rightarrow$ présence de \ce{CO2}};
\draw[->, popgreen, thick] (5.5,-0.9) -- (5.35,-0.7);
\node[poporange, align=left] at (-4.9,-0.8) {dégagement\\ de bulles de gaz};
\draw[->, poporange, thick] (-3.6,-0.8) -- (-2.0,-0.9);
\end{tikzpicture}
\end{center}
\legende{Dispositif expérimental de mise en évidence de la fermentation alcoolique. Les levures additionnées au jus de canne dégagent un gaz qui, conduit dans l'eau de chaux, la trouble : ce gaz est du \ce{CO2}. Le liquide du ballon s'enrichit simultanément en éthanol.}
\end{popfigure}

\begin{methode}[Interpréter un test à l'eau de chaux]
L'eau de chaux est une solution aqueuse d'hydroxyde de calcium, \ce{Ca(OH)2}, limpide à l'état initial. Pour interpréter le test :
\begin{enumerate}
\item Observer l'aspect initial : solution \emph{limpide} (transparente).
\item Faire barboter le gaz à identifier dans la solution.
\item Si l'eau de chaux se \emph{trouble} (dépôt blanc de carbonate de calcium \ce{CaCO3}), conclure à la présence de \ce{CO2} : \ce{Ca(OH)2 + CO2 -> CaCO3 + H2O}.
\item Si elle reste limpide, le gaz testé n'est pas du \ce{CO2} (ou en contient trop peu).
\end{enumerate}
\textbf{Rédaction attendue au Bac :} \og l'eau de chaux se trouble en présence de dioxyde de carbone \fg{} — et non \og l'eau de chaux devient blanche \fg{}, formulation trop vague.
\end{methode}

\courssub{Expériences témoins : le facteur vivant est indispensable}

Une conclusion scientifique solide exige des \cle{témoins}. Reprenons le dispositif précédent en modifiant un seul paramètre à la fois :

\begin{center}
\begin{tabularx}{\linewidth}{|l|X|X|X|}
\hline
\rowcolor{popblueL}
\textbf{Lot} & \textbf{Contenu du ballon} & \textbf{Eau de chaux} & \textbf{Éthanol produit} \\
\hline
A (expérience) & Jus de canne + levures vivantes & Se trouble & Oui \\
\hline
B (témoin 1) & Jus de canne seul, sans levures & Reste limpide & Non \\
\hline
C (témoin 2) & Jus de canne + levures \emph{bouillies} & Reste limpide & Non \\
\hline
\end{tabularx}
\end{center}

\textbf{Interprétation.} Le lot B montre que le sucre ne se transforme pas spontanément : la levure est nécessaire. Le lot C est encore plus instructif : les levures bouillies sont bien présentes (leur matière est là), mais elles sont \emph{mortes} — la chaleur a détruit (dénaturé) leurs \cle{enzymes}. Or rien ne se produit. On en déduit que ce n'est pas la simple présence de la substance \og levure \fg{} qui fermente le sucre, mais l'activité d'enzymes contenues dans des cellules \emph{vivantes}.

\begin{propriete}[La fermentation alcoolique, propriété établie expérimentalement]
La fermentation alcoolique est une transformation du glucose en éthanol et en dioxyde de carbone, réalisée par des \textbf{levures vivantes} (grâce à leurs enzymes), en \textbf{absence de dioxygène} (anaérobiose). Elle s'accompagne d'un dégagement modéré d'énergie, stockée sous forme d'ATP utilisable par la cellule. Cette propriété, établie par les expériences de Pasteur et les témoins ci-dessus, est exigible au Baccalauréat.
\end{propriete}

\courssub{Interprétation biochimique : le bilan de la fermentation alcoolique}

Pourquoi les levures font-elles cela ? Pour survivre. En l'absence de dioxygène, la levure ne peut pas réaliser la respiration cellulaire complète ; elle dégrade alors le glucose de façon \emph{partielle} pour en extraire un peu d'énergie. Les molécules d'éthanol et de \ce{CO2} sont les déchets de cette dégradation incomplète — déchets précieux pour nous, puisque l'éthanol est un carburant.

\begin{aretenir}[Équation bilan de la fermentation alcoolique]
\[ \ce{C6H12O6 -> 2 C2H5OH + 2 CO2} + \text{énergie (ATP)} \]
Une mole de glucose (\SI{180}{\gram}) fournit deux moles d'éthanol ($2 \times \SI{46}{\gram} = \SI{92}{\gram}$) et deux moles de dioxyde de carbone. Ce bilan est à connaître parfaitement ; les étapes intermédiaires (glycolyse, formation d'acétyldéhyde) sont admises et \textbf{non exigibles}.
\end{aretenir}

\begin{attention}[Erreur fréquente : fermentation $\neq$ respiration]
Ne confonds jamais fermentation alcoolique et respiration cellulaire :
\begin{itemize}
\item la \textbf{respiration} (\ce{C6H12O6 + 6O2 -> 6CO2 + 6H2O}) est une dégradation \emph{complète} du glucose, qui \emph{consomme du dioxygène} et dégage \emph{beaucoup} d'énergie (environ 36 ATP par glucose) ;
\item la \textbf{fermentation} est une dégradation \emph{incomplète}, qui se déroule \emph{sans dioxygène} et ne dégage que \emph{peu} d'énergie (2 ATP par glucose) : l'éthanol produit contient encore l'essentiel de l'énergie chimique du glucose — c'est précisément pourquoi il peut servir de carburant !
\end{itemize}
Écrire \og la levure respire le sucre en produisant de l'alcool \fg{} est une erreur classique sanctionnée à l'examen.
\end{attention}

\courssub{Conditions nécessaires à la fermentation}

Les expériences et la pratique brassicole montrent que la fermentation alcoolique exige quatre conditions réunies simultanément :
\begin{itemize}
\item la présence d'un \textbf{sucre fermentescible} (glucose, saccharose hydrolysé, sucres issus de l'hydrolyse de l'amidon) ;
\item la présence de \textbf{levures vivantes} en quantité suffisante ;
\item l'\textbf{absence de dioxygène} (milieu fermé, anaérobiose) : en présence d'air, la levure préfère respirer et produit peu d'alcool ;
\item une \textbf{température modérée}, comprise entre 25 et \SI{35}{\celsius} : en dessous, les enzymes sont peu actives ; au-dessus de \SI{40}{\celsius} environ, elles se dénaturent et les levures meurent.
\end{itemize}

\courssub{Étude cinétique : évolution de la fermentation au cours du temps}

Suivons, dans une cuve de fermentation industrielle, la population de levures et la concentration en éthanol pendant 50 heures.

\begin{popfigure}
\begin{center}
\begin{tikzpicture}
% Premier axe : Éthanol (axe gauche)
\begin{axis}[
  width=0.92\linewidth, height=7cm,
  xlabel={Temps (h)},
  ylabel={Éthanol (\si{\gram\per\liter})},
  axis y line*=left,
  axis x line*=bottom,
  xmin=0, xmax=50, ymin=0, ymax=110,
  xtick={0,10,20,30,40,50},
  legend style={at={(0.02,0.98)}, anchor=north west, font=\small, cells={anchor=west}},
  every axis plot/.append style={very thick}
]
\addplot[poporange, domain=0:50, samples=200] {100/(1+exp(-(x-22)/5))};
\addlegendentry{Concentration en éthanol}
% Phases sur le premier axe
\draw[dashed, popmuted] (axis cs:10,0) -- (axis cs:10,110);
\draw[dashed, popmuted] (axis cs:30,0) -- (axis cs:30,110);
\node[font=\footnotesize, popdark, anchor=north] at (axis cs:5,105) {Phase latente};
\node[font=\footnotesize, popdark, anchor=north] at (axis cs:20,105) {Phase exponentielle};
\node[font=\footnotesize, popdark, anchor=north] at (axis cs:40,105) {Plateau};
\end{axis}
% Deuxième axe : Population levures (axe droit)
\begin{axis}[
  width=0.92\linewidth, height=7cm,
  ylabel={Population de levues (unités arbitraires)},
  axis y line*=right,
  axis x line=none,
  xmin=0, xmax=50, ymin=0, ymax=70,
  xtick=\empty,
  legend style={at={(0.02,0.88)}, anchor=north west, font=\small, cells={anchor=west}},
  every axis plot/.append style={very thick}
]
\addplot[popblue, domain=0:50, samples=200] {8 + 52/(1+exp(-(x-14)/3.5))};
\addlegendentry{Population de levures}
\end{axis}
\end{tikzpicture}
\end{center}
\legende{Évolution de la population de levures (axe droit, bleu) et de la concentration en éthanol (axe gauche, orange) dans une cuve de fermentation (jus de canne, \SI{30}{\celsius}, anaérobiose). Après une phase latente, la production d'éthanol s'accélère (phase exponentielle) puis plafonne : l'alcool accumulé devient toxique pour les levures.}
\end{popfigure}

L'exploitation de cette courbe fait apparaître trois phases :
\begin{enumerate}
\item \textbf{Phase latente} (0 à environ 10 h) : les levures se multiplient et s'adaptent au milieu ; la production d'éthanol est encore faible.
\item \textbf{Phase exponentielle} (10 à 30 h) : la population de levures est maximale et très active ; l'éthanol s'accumule rapidement, de façon quasi proportionnelle au temps.
\item \textbf{Plateau} (au-delà de 30 h) : la production ralentit puis s'arrête. Deux causes se combinent : l'\emph{épuisement du sucre} et surtout la \emph{toxicité de l'éthanol} accumulé, qui inhibe puis tue les levures.
\end{enumerate}

\begin{savaistu}[Pourquoi faut-il distiller ?]
Au-delà d'une teneur de 12 à 15 \% d'alcool en volume, l'éthanol devient toxique pour les levures elles-mêmes : elles meurent et la fermentation cesse spontanément. C'est pourquoi aucune boisson fermentée naturelle (vin, bil-bil) ne dépasse ce seuil. Pour obtenir le bioéthanol carburant, qui doit être presque pur (plus de 95 \%), il faut une étape supplémentaire : la \textbf{distillation}, qui sépare l'éthanol de l'eau grâce à la différence de leurs températures d'ébullition (\SI{78}{\celsius} pour l'éthanol, \SI{100}{\celsius} pour l'eau). C'est le principe des alambics traditionnels utilisés pour produire l'odontol.
\end{savaistu}

\begin{exemplebox}[Un exemple chiffré à retenir]
D'après l'équation bilan, \SI{180}{\gram} de glucose donnent théoriquement \SI{92}{\gram} d'éthanol. Comme la masse volumique de l'éthanol est d'environ \SI{0,79}{\kilogram\per\liter}, il faut environ \SI{1,55}{\kilogram} de glucose pour produire théoriquement \SI{1}{\liter} d'éthanol pur. En pratique, les rendements réels des distilleries atteignent 85 à 90 \% de cette valeur théorique.
\end{exemplebox}

\begin{exoresolu}[Rendement d'une fermentation de jus de canne]
\textbf{Énoncé.} Une coopérative de la région du Littoral fait fermenter une cuve contenant \SI{340}{\kilogram} de glucose issu de jus de canne, dans des conditions optimales (anaérobiose, \SI{30}{\celsius}). On admet que la fermentation est totale. 1) Écrire l'équation bilan de la fermentation alcoolique. 2) Calculer la masse théorique d'éthanol produite. 3) En déduire le volume théorique d'éthanol obtenu, sachant que sa masse volumique vaut \SI{0,79}{\kilogram\per\liter}. 4) Le rendement réel de la coopérative est de 85 \%. Quel volume d'éthanol est réellement obtenu ?
\tcblower
\textbf{Solution détaillée.}
1) Équation bilan :
\[ \ce{C6H12O6 -> 2 C2H5OH + 2 CO2} + \text{énergie (ATP)} \]
2) D'après la stœchiométrie, \SI{180}{\gram} de glucose ($M = \SI{180}{\gram\per\mol}$) donnent $2 \times \SI{46}{\gram} = \SI{92}{\gram}$ d'éthanol ($M = \SI{46}{\gram\per\mol}$). D'où :
\begin{align*}
m_{\text{éthanol}} &= 340 \times \dfrac{92}{180} \\
&\approx \SI{173,8}{\kilogram}
\end{align*}
3) Volume théorique :
\[ V = \dfrac{m}{\rho} = \dfrac{173,8}{0,79} \approx \SI{220}{\liter} \]
4) Volume réel avec un rendement de 85 \% :
\[ V_{\text{réel}} = 220 \times 0,85 \approx \SI{187}{\liter} \]
\textbf{Conclusion :} la cuve fournit environ \SI{187}{\liter} d'éthanol. Remarque de rédaction : toujours distinguer la quantité \emph{théorique} (issue de l'équation bilan) de la quantité \emph{réelle} (affectée par le rendement).
\end{exoresolu}

\begin{aretenir}[L'essentiel de la section]
\begin{itemize}
\item La fermentation alcoolique est réalisée par des \textbf{levures vivantes} (enzymes), en \textbf{anaérobiose}, entre 25 et \SI{35}{\celsius}, sur un milieu sucré.
\item Bilan exigible : \ce{C6H12O6 -> 2 C2H5OH + 2 CO2} + énergie (peu d'ATP).
\item Le \ce{CO2} dégagé est mis en évidence par le \textbf{trouble de l'eau de chaux}.
\item Les témoins (sans levure, levures bouillies) prouvent que le facteur vivant est indispensable.
\item La cinétique comporte trois phases (latente, exponentielle, plateau) ; le plateau traduit la toxicité de l'éthanol (mort des levures vers 12--15 \% d'alcool), d'où la nécessité de la distillation.
\item Ordre de grandeur : environ \SI{1,55}{\kilogram} de glucose pour \SI{1}{\liter} d'éthanol théorique.
\end{itemize}
\end{aretenir}

Le mécanisme biologique étant désormais établi, il reste à voir comment ce savoir est exploité à grande échelle : c'est l'objet de la section suivante, consacrée à la production industrielle du bioéthanol et aux matières premières que le Cameroun peut mobiliser.

\coursec{Production du bioéthanol et matières premières camerounaises}

Dans la section précédente, nous avons établi expérimentalement le fondement du bioéthanol : la \cle{fermentation alcoolique}, réalisée par les levures, transforme le \cle{glucose} en éthanol et en dioxyde de carbone selon le bilan $\ce{C6H12O6 -> 2 C2H5OH + 2 CO2}$. Mais ce résultat de laboratoire soulève immédiatement une question pratique : d'où vient le glucose en quantité suffisante pour produire du carburant à l'échelle d'un pays ? Et comment passe-t-on d'un jus sucré fermenté à un éthanol utilisable dans un moteur ? C'est toute la chaîne de production du bioéthanol que nous allons maintenant décrire, en l'illustrant avec les matières premières dont dispose le Cameroun : la canne à sucre, le manioc, le maïs.

\courssub{La chaîne de production du bioéthanol : de la plante au carburant}

\begin{definition}[Bioéthanol]
Le \cle{bioéthanol} est un éthanol ($\ce{C2H5OH}$) produit par \textbf{fermentation alcoolique} de sucres issus de la \cle{biomasse} végétale (matières sucrées ou amylacées), puis purifié par distillation. Il est utilisé comme carburant, pur ou en mélange avec l'essence.
\end{definition}

Observe une exploitation sucrière comme celle de la SOSUCAM à Mbandjock : des tonnes de canne y arrivent chaque jour. Intuitivement, on comprend que le sucre contenu dans les tiges peut nourrir des levures, exactement comme le jus de raisin nourrit les levures du vin. Mais entre la tige de canne coupée au champ et le litre d'éthanol qui sort de l'usine, il y a toute une succession d'étapes ordonnées, chacune ayant un rôle précis.

\begin{methode}[Décrire une filière de production en étapes ordonnées]
Pour décrire rigoureusement une filière de production d'un biocarburant (démarche exigible au Bac), procède toujours dans l'ordre suivant :
\begin{enumerate}
\item \textbf{Matière première} : identifier la biomasse utilisée (canne, manioc, maïs...) et la nature du sucre qu'elle contient (saccharose, amidon).
\item \textbf{Prétraitement} : extraction du sucre (broyage, pressage) et, si nécessaire, hydrolyse de l'amidon en glucose.
\item \textbf{Fermentation} : transformation du glucose en éthanol par les levures, en anaérobiose.
\item \textbf{Distillation} : séparation de l'éthanol du moût fermenté grâce à sa température d'ébullition ($\SI{78,4}{\celsius}$) inférieure à celle de l'eau ($\SI{100}{\celsius}$).
\item \textbf{Déshydratation} : élimination de l'eau résiduelle pour obtenir l'éthanol anhydre (carburant).
\end{enumerate}
\end{methode}

Appliquons cette méthode à la filière complète :

\begin{itemize}
\item \textbf{Récolte} : la canne à sucre est coupée à maturité (12 à 18 mois après plantation) ; le manioc est arraché après 9 à 12 mois de culture.
\item \textbf{Extraction et broyage} : la canne est broyée et pressée pour en extraire le jus sucré (le \emph{vesou}), riche en saccharose. Le manioc, lui, est pelé, râpé et délayé dans l'eau pour obtenir une suspension d'amidon.
\item \textbf{Hydrolyse (uniquement pour les matières amylacées)} : l'amidon est un polysaccharide que les levures ne savent pas utiliser directement. On le découpe en glucose grâce à des enzymes, les \cle{amylases}, selon le bilan : $\ce{(C6H10O5)_{n} + n H2O ->[amylases] n C6H12O6}$.
\item \textbf{Fermentation} : des levures (\emph{Saccharomyces cerevisiae}) sont ensemencées dans le jus ou le moût glucosé, en cuves fermées. En quelques dizaines d'heures, le glucose est converti en éthanol : $\ce{C6H12O6 -> 2 C2H5OH + 2 CO2}$. Précision utile : le saccharose du jus de canne est d'abord hydrolysé en glucose et fructose par l'\textbf{invertase}, une enzyme que la levure possède naturellement ; la filière canne n'a donc pas besoin d'une étape d'hydrolyse industrielle séparée.
\item \textbf{Distillation} : le moût fermenté contient environ 8 à 12 \% d'éthanol en volume. On le chauffe : l'éthanol, plus volatil que l'eau, s'évapore en premier ; sa vapeur est recueillie puis condensée.
\item \textbf{Déshydratation} : la distillation seule ne permet pas de dépasser environ 96 \% d'éthanol en volume (le mélange eau--éthanol forme un azéotrope). Une étape supplémentaire élimine l'eau restante pour obtenir l'\textbf{éthanol anhydre} (à plus de 99 \%), indispensable pour le mélange avec l'essence.
\end{itemize}

\begin{popfigure}
\begin{center}
\begin{tikzpicture}[
  font=\small,
  etape/.style={rectangle, rounded corners=3pt, draw=popblue, fill=popblueL, text width=2.9cm, align=center, minimum height=0.9cm},
  etapeb/.style={rectangle, rounded corners=3pt, draw=popgreen, fill=popgreenL, text width=2.9cm, align=center, minimum height=0.9cm},
  comm/.style={rectangle, rounded corners=3pt, draw=poporange, fill=poporangeL, text width=3.2cm, align=center, minimum height=0.9cm},
  fleche/.style={-{Stealth[length=2.5mm]}, thick, popdark}
]
% Filière canne
\node[etape, align=center] (canne) at (0,4.5){\textbf{Canne à sucre}\\ (saccharose)};
\node[etape] (broyage) at (4.2,4.5) {Broyage, pressage $\rightarrow$ jus sucré};
% Filière manioc
\node[etapeb, align=center] (manioc) at (0,1.5){\textbf{Manioc, maïs}\\ (amidon)};
\node[etapeb] (rasp) at (4.2,1.5) {Râpage, délayage $\rightarrow$ suspension d'amidon};
\node[etapeb] (hydro) at (8.4,1.5) {Hydrolyse par les amylases $\rightarrow$ glucose};
% Étapes communes
\node[comm, align=center] (ferm) at (8.4,4.5){Fermentation\\ (levures, anaérobiose)};
\node[comm, align=center] (dist) at (12.4,4.5){Distillation\\ (éthanol $\approx$ 96 \%)};
\node[comm, align=center] (desh) at (12.4,1.5){Déshydratation\\ (éthanol anhydre $>$ 99 \%)};
\node[comm, fill=poppinkL, draw=poppink, align=center] (carb) at (12.4,-1.2){\textbf{Bioéthanol carburant}\\ (E10, E85...)};
% Flèches
\draw[fleche] (canne) -- (broyage);
\draw[fleche] (broyage) -- (ferm);
\draw[fleche] (manioc) -- (rasp);
\draw[fleche] (rasp) -- (hydro);
\draw[fleche] (hydro) -- (ferm);
\draw[fleche] (ferm) -- (dist);
\draw[fleche] (dist) -- (desh);
\draw[fleche] (desh) -- (carb);
\end{tikzpicture}
\end{center}
\legende{Organigramme de la production de bioéthanol. Filière sucrière (bleu) : le saccharose du jus de canne est hydrolysé par l'invertase de la levure, il est donc directement utilisable. Filière amylacée (vert) : l'amidon du manioc ou du maïs exige une étape supplémentaire, l'hydrolyse enzymatique en glucose. Les étapes de fermentation, distillation et déshydratation (orange) sont communes aux deux filières.}
\end{popfigure}

\begin{attention}[L'amidon n'est pas directement fermentescible]
Erreur très fréquente au Bac : écrire que les levures fermentent l'amidon. C'est \textbf{faux}. Les levures ne possèdent pas les enzymes capables de découper l'amidon. Une \textbf{hydrolyse préalable} (par des amylases, ou par chauffage en milieu acide) est indispensable pour transformer l'amidon en glucose, seul substrat de la fermentation alcoolique. Conséquence pratique : la filière manioc/maïs comporte une étape de plus que la filière canne, ce qui la rend plus coûteuse en énergie et en temps.
\end{attention}

\courssub{Les matières premières mobilisables au Cameroun}

Le Cameroun, par la diversité de ses climats (zone forestière humide du Sud, savanes du Nord), dispose de plusieurs matières premières végétales riches en sucres ou en amidon. On les classe en deux grandes catégories, qui correspondent exactement aux deux filières technologiques décrites ci-dessus : les \textbf{matières sucrées} et les \textbf{matières amylacées}.

\courssub{Les matières sucrées : la canne à sucre, filière reine}

La \cle{canne à sucre} (\emph{Saccharum officinarum}) accumule dans ses tiges du \textbf{saccharose}, un disaccharide que l'invertase des levures hydrolyse facilement en glucose et fructose, tous deux fermentescibles. C'est la matière première la plus performante pour le bioéthanol.

Au Cameroun, la \textbf{SOSUCAM} (Société sucrière du Cameroun) exploite de vastes plantations à \textbf{Mbandjock} et \textbf{Nkoteng}, dans la région du Centre. Ces unités, qui produisent aujourd'hui du sucre de consommation, disposent de toute l'infrastructure (broyeurs, cuves, distilleries potentielles) pour développer une filière bioéthanol : les mélasses (résidus sirupeux de la cristallisation du sucre) et une partie du jus de canne pourraient alimenter des fermenteurs.

\begin{exemplebox}[Un rendement record : la canne au Brésil]
La canne à sucre produit environ \textbf{6 000 à 7 000 litres d'éthanol par hectare et par an} au Brésil, grâce à un climat chaud et humide favorable à une photosynthèse intense (la canne est une plante en $\mathrm{C}_4$, très efficace pour fixer le $\ce{CO2}$). C'est le meilleur rendement énergétique parmi les cultures à biocarburant de première génération.
\end{exemplebox}

\courssub{Les matières amylacées : manioc, maïs, patate douce}

Le \cle{manioc} (\emph{Manihot esculenta}) est l'une des cultures vivrières les plus importantes du Cameroun : ses racines tubérisées contiennent 25 à 30 \% d'amidon en masse fraîche. Le \textbf{maïs} et la \textbf{patate douce} sont également riches en amidon. Ces trois plantes peuvent alimenter une filière bioéthanol, à condition d'ajouter l'étape d'\textbf{hydrolyse enzymatique} décrite plus haut. Aux États-Unis, c'est le maïs qui fournit l'essentiel du bioéthanol national ; au Cameroun, le manioc, très rustique et largement cultivé dans les régions du Centre, du Sud et du Littoral, serait un candidat naturel.

\courssub{Comparaison des rendements : la canne devance les plantes à amidon}

Pour comparer objectivement les filières, on utilise le \textbf{rendement en éthanol par hectare cultivé et par an}, qui intègre à la fois la productivité agricole et la richesse en sucres fermentescibles.

\begin{popfigure}
\begin{center}
\begin{tikzpicture}
\begin{axis}[
  ybar, bar width=1.1cm,
  width=0.85\linewidth, height=6.5cm,
  ymin=0, ymax=7500,
  ylabel={Rendement en éthanol (L/ha/an)},
  symbolic x coords={Manioc, Maïs, Canne à sucre},
  xtick=data,
  nodes near coords,
  every node near coord/.append style={font=\small},
  ymajorgrids, grid style={popline!40},
  axis line style={popdark},
  tick label style={font=\small},
  label style={font=\small}
]
\addplot[fill=popblue!70, draw=popblue] coordinates {(Manioc,1500) (Maïs,3000) (Canne à sucre,6500)};
\end{axis}
\end{tikzpicture}
\end{center}
\legende{Rendements moyens en éthanol (ordres de grandeur, en litres par hectare et par an) des trois principales matières premières mobilisables au Cameroun. La canne à sucre (environ 6 500 L/ha/an) dépasse nettement le maïs (environ 3 000 L/ha/an) et le manioc (environ 1 500 L/ha/an).}
\end{popfigure}

\begin{center}
\begin{tabularx}{\linewidth}{|l|X|X|X|}
\hline
\rowcolor{popblueL} \textbf{Critère} & \textbf{Canne à sucre} & \textbf{Maïs} & \textbf{Manioc} \\
\hline
Sucre stocké & Saccharose (hydrolysé par l'invertase de la levure) & Amidon (hydrolyse industrielle requise) & Amidon (hydrolyse industrielle requise) \\
\hline
Rendement en éthanol (L/ha/an) & 6 000 à 7 000 & $\approx$ 3 000 & $\approx$ 1 500 \\
\hline
Étapes de production & 5 (sans hydrolyse industrielle) & 6 (avec hydrolyse) & 6 (avec hydrolyse) \\
\hline
Bilan énergétique & Très favorable & Moyen & Moyen \\
\hline
Conflit avec l'alimentation & Modéré (sucre) & Fort (céréale vivrière) & Fort (tubercule vivrier) \\
\hline
\end{tabularx}
\end{center}

\begin{aretenir}[L'essentiel sur les matières premières]
La canne à sucre donne le \textbf{meilleur rendement énergétique} car son saccharose est directement utilisable par les levures et sa productivité photosynthétique est élevée. Les plantes à amidon (manioc, maïs, patate douce) exigent une \textbf{hydrolyse préalable} par les amylases, ce qui alourdit la filière et réduit le bilan énergétique.
\end{aretenir}

\courssub{Les utilisations du bioéthanol}

L'éthanol produit peut être valorisé de plusieurs façons :

\begin{itemize}
\item \textbf{Carburant en mélange} : l'éthanol anhydre est mélangé à l'essence. La notation indique la proportion d'éthanol en volume : \cle{E10} contient 10 \% d'éthanol et 90 \% d'essence (utilisable dans la plupart des moteurs actuels sans modification) ; \cle{E85} contient jusqu'à 85 \% d'éthanol (réservé aux véhicules dits \emph{flex-fuel}, à moteur adapté).
\item \textbf{Alcool à brûler} : combustible domestique pour les réchauds, qui peut remplacer le bois de chauffe et ainsi limiter la déforestation des savanes et forêts camerounaises.
\item \textbf{Désinfection et pharmacie} : l'éthanol à 70 \% est un antiseptique de base dans les centres de santé ; il sert aussi de solvant dans l'industrie pharmaceutique et cosmétique.
\end{itemize}

\begin{exoresolu}[Exploiter des données de production]
\textbf{Énoncé.} Une usine envisage de produire du bioéthanol au Cameroun à partir de canne à sucre. On estime le rendement à 6 500 L d'éthanol par hectare et par an. 1) Quelle surface de plantation faut-il pour produire 13 millions de litres d'éthanol par an ? 2) Sachant qu'un litre d'essence est remplacé par environ 1,5 L d'éthanol (l'éthanol dégage moins d'énergie par litre), quel volume d'essence cette production permet-elle de remplacer ? 3) Pourquoi la même surface plantée en manioc ne suffirait-elle pas ?
\tcblower
\textbf{Solution.} 1) Surface nécessaire :
\[ S = \frac{13\,000\,000\ \text{L}}{6\,500\ \text{L/ha}} = 2\,000\ \text{ha}. \]
Il faut donc 2 000 hectares de canne à sucre, soit l'équivalent d'une plantation moyenne de la SOSUCAM.

2) Volume d'essence remplacé :
\[ V = \frac{13\,000\,000}{1,5} \approx 8\,700\,000\ \text{L}, \]
soit environ 8,7 millions de litres d'essence économisés par an.

3) Le manioc ne fournit qu'environ 1 500 L/ha/an, soit plus de quatre fois moins que la canne. Il faudrait donc environ 8 700 ha de manioc pour la même production. De plus, la filière manioc exige une hydrolyse préalable de l'amidon, qui consomme de l'énergie et dégrade encore le bilan. Enfin, détourner le manioc de l'alimentation humaine poserait un grave problème de sécurité alimentaire au Cameroun.
\end{exoresolu}

\courssub{Complément hors programme strict : les générations de biocarburants}

\begin{savaistu}[Le Brésil, pionnier de l'éthanol-carburant]
Dès les années 1970, après les chocs pétroliers, le Brésil lance le programme \emph{Proálcool} pour réduire sa dépendance au pétrole importé. Depuis, une grande partie des voitures brésiliennes roule à l'éthanol de canne, pur ou en mélange, et la quasi-totalité des véhicules neufs y sont \emph{flex-fuel}. Cet exemple montre qu'un pays tropical, au climat proche de celui du Cameroun, peut bâtir une filière biocarburant durable à grande échelle.
\end{savaistu}

Pour aller plus loin (notion non exigible mais valorisée dans une copie), on distingue :

\begin{itemize}
\item les \cle{biocarburants de première génération}, produits à partir de \textbf{cultures alimentaires} (canne, maïs, manioc, palme à huile). Leur limite majeure est la \textbf{compétition avec l'alimentation humaine} : chaque hectare de manioc destiné au carburant est un hectare qui ne nourrit plus les populations ;
\item les \cle{biocarburants de deuxième génération}, produits à partir de \textbf{résidus non alimentaires} : pailles, bagasse de canne, tiges de maïs, déchets de scierie. Ces matériaux \textbf{lignocellulosiques} (cellulose + lignine) demandent des prétraitements plus complexes, mais ils valorisent des déchets sans toucher aux cultures vivrières. Au Cameroun, la bagasse de la SOSUCAM et les résidus de scieries du Sud-Est constitueraient des ressources considérables.
\end{itemize}

\begin{attention}[Rédiger sans exagérer]
Ne conclus jamais que le bioéthanol est « la solution miracle » au déficit énergétique. Une copie de Terminale D doit rester nuancée : la filière consomme des terres, de l'eau et de l'énergie (distillation), et la première génération entre en concurrence avec l'alimentation. Formule toujours un avis argumenté : le bioéthanol est un \emph{complément} renouvelable aux énergies fossiles, dont l'intérêt dépend de la matière première choisie et de la gestion des terres.
\end{attention}

\begin{exercice}[À toi de jouer]
Une coopérative agricole de la région du Centre hésite entre deux projets : produire de l'éthanol à partir de jus de canne ou à partir de manioc.
\begin{enumerate}
\item Construis, pour chaque projet, la liste ordonnée des étapes de production, en justifiant l'étape supplémentaire de la filière manioc.
\item À l'aide des rendements du cours, calcule la surface nécessaire dans chaque cas pour produire 3 millions de litres d'éthanol par an.
\item Rédige un court texte (10 lignes maximum) conseillant la coopérative, en tenant compte du bilan énergétique et de la sécurité alimentaire.
\end{enumerate}
\end{exercice}

\begin{corrige}[Exercice]
1) Filière canne : récolte $\rightarrow$ broyage/pressage $\rightarrow$ fermentation $\rightarrow$ distillation $\rightarrow$ déshydratation. Filière manioc : récolte $\rightarrow$ râpage/délayage $\rightarrow$ \textbf{hydrolyse de l'amidon par les amylases} $\rightarrow$ fermentation $\rightarrow$ distillation $\rightarrow$ déshydratation. L'étape d'hydrolyse est indispensable car l'amidon, polysaccharide, n'est pas utilisable directement par les levures : seul le glucose est fermentescible.

2) Canne : $S = \dfrac{3\,000\,000}{6\,500} \approx 462$ ha. Manioc : $S = \dfrac{3\,000\,000}{1\,500} = 2\,000$ ha, soit plus de quatre fois plus de surface.

3) Éléments attendus : la filière canne est plus rentable (meilleur rendement, pas d'hydrolyse industrielle, bilan énergétique favorable) ; mais la canne est déjà mobilisée par l'industrie sucrière. Le manioc est une culture vivrière essentielle : le détourner vers le carburant menacerait l'alimentation des populations. Un bon conseil : privilégier la canne ou, mieux, valoriser les résidus (mélasses, bagasse) pour éviter tout conflit avec l'alimentation.
\end{corrige}

\coursec{Biodiesel et biogaz : les autres biocarburants}

Dans la section précédente, nous avons vu comment la fermentation alcoolique transforme les sucres de la canne, du maïs ou du manioc en bioéthanol, carburant des moteurs à essence. Mais le bioéthanol ne répond pas à tous les besoins : les camions, les bus et les groupes électrogènes qui sillonnent le Cameroun fonctionnent au \emph{diesel}, et des millions de tonnes de déchets organiques (fumier, résidus de cuisine, eaux usées) s'accumulent sans être valorisées. Deux autres biocarburants répondent à ces situations : le \cle{biodiesel}, obtenu à partir des huiles végétales, et le \cle{biogaz}, produit par la décomposition des déchets. Étudions-les tour à tour.

\courssub{Le biodiesel : transformer les huiles végétales en carburant}

\textbf{Observation et problème.} Une huile végétale brute (huile de palme, par exemple) brûle et dégage de l'énergie : on pourrait donc imaginer la verser directement dans le réservoir d'un moteur diesel. En pratique, cela échoue rapidement : l'huile brute est trop \emph{visqueuse}, elle s'écoule mal, encrasse les injecteurs et forme des dépôts dans les cylindres. \textbf{Problème scientifique :} comment modifier chimiquement une huile végétale pour obtenir un liquide fluide, proche du gazole, utilisable dans un moteur diesel classique ?

\textbf{Hypothèse.} La viscosité des huiles vient de la nature de leurs molécules, les \cle{triglycérides}, de grosses molécules formées d'une molécule de glycérol liée à trois acides gras. Si l'on « casse » ces grosses molécules en molécules plus petites, le liquide devrait devenir plus fluide.

\textbf{Vérification expérimentale.} Au laboratoire, on chauffe doucement de l'huile végétale avec du \cle{méthanol} en présence d'un catalyseur, la \cle{soude} (hydroxyde de sodium, \ce{NaOH}). Après réaction et repos, on observe la séparation de \textbf{deux couches liquides} : une couche supérieure claire et fluide (les esters méthyliques, c'est-à-dire le biodiesel) et une couche inférieure plus dense (la glycérine, aussi appelée glycérol). La couche supérieure, beaucoup moins visqueuse que l'huile de départ, brûle proprement dans un moteur diesel. L'hypothèse est validée.

\begin{definition}[Transestérification]
La \cle{transestérification} est une réaction chimique au cours de laquelle un \textbf{triglycéride} (huile ou graisse) réagit avec un \textbf{alcool} (généralement le méthanol \ce{CH3OH}), en présence d'un \textbf{catalyseur} (soude \ce{NaOH} ou potasse \ce{KOH}), pour donner des \textbf{esters méthyliques} (le biodiesel) et de la \textbf{glycérine} (coproduit). Bilan simplifié :
\[ \text{Triglycéride} + 3\ \ce{CH3OH} \xrightarrow{\text{NaOH}} 3\ \text{Esters méthyliques (biodiesel)} + \text{Glycérol} \]
\end{definition}

\begin{popfigure}
\begin{center}
\begin{tikzpicture}[scale=0.9, every node/.style={font=\small}]
% Triglycéride
\draw[thick, popblue, rounded corners] (-6.2,-0.9) rectangle (-3.4,0.9);
\node[popblue, font=\small\bfseries] at (-4.8,0.45) {Triglycéride};
\node[popmuted] at (-4.8,0.0) {(huile végétale)};
\node[popmuted] at (-4.8,-0.5) {glycérol + 3 acides gras};
% Méthanol
\draw[thick, poporange, rounded corners] (-2.6,-0.6) rectangle (-0.8,0.6);
\node[poporange, font=\small\bfseries] at (-1.7,0.2) {3 méthanol};
\node[popmuted] at (-1.7,-0.3) {\ce{CH3OH}};
% Plus
\node[font=\Large] at (-3.0,0) {$+$};
% Flèche réaction
\draw[-{Stealth[length=3mm]}, very thick, popdark] (-0.5,0) -- (1.6,0);
\node[above, popdark] at (0.55,0.05) {soude \ce{NaOH}};
\node[below, popmuted] at (0.55,-0.08) {(catalyseur)};
% Esters
\draw[thick, popgreen, rounded corners] (2.0,-0.9) rectangle (4.8,0.9);
\node[popgreen, font=\small\bfseries] at (3.4,0.45) {3 esters méthyliques};
\node[popmuted] at (3.4,0.0) {(biodiesel)};
\node[popmuted] at (3.4,-0.5) {carburant moteur diesel};
% Plus 2
\node[font=\Large] at (5.2,0) {$+$};
% Glycérine
\draw[thick, poppurple, rounded corners] (5.6,-0.6) rectangle (7.4,0.6);
\node[poppurple, font=\small\bfseries] at (6.5,0.2) {Glycérine};
\node[popmuted] at (6.5,-0.3) {(coproduit)};
\end{tikzpicture}
\end{center}
\legende{Schéma de la réaction de transestérification : une molécule de triglycéride réagit avec trois molécules de méthanol en présence de soude (catalyseur) pour donner trois molécules d'esters méthyliques (biodiesel) et une molécule de glycérine, valorisable en savonnerie et en cosmétique.}
\end{popfigure}

\textbf{Matières premières mobilisables au Cameroun.} Le Cameroun dispose d'atouts considérables pour le biodiesel :
\begin{itemize}
\item l'\textbf{huile de palme}, produite à grande échelle par les agro-industries (SOCAPALM, PHC -- Plantations du Haut Penja) et par d'innombrables artisans ruraux ;
\item l'\textbf{huile de coton}, coproduit de la filière cotonnière du Nord (SODECOTON) ;
\item le \textbf{jatropha} (\emph{Jatropha curcas}), plante \emph{non alimentaire} dont les graines très riches en huile poussent sur des sols pauvres et dégradés, sans entrer en compétition avec les cultures vivrières : c'est une piste prometteuse pour éviter le conflit « carburant contre nourriture ».
\end{itemize}

\begin{propriete}[Avantages techniques du biodiesel]
\begin{itemize}
\item Il est \textbf{utilisable dans les moteurs diesel classiques}, pur (B100) ou en mélange avec le gazole (B7, B30), sans modification majeure du moteur.
\item Il est \textbf{biodégradable} : en cas de déversement accidentel, il se décompose rapidement dans l'environnement.
\item Sa combustion dégage \textbf{beaucoup moins de fumées sulfurées} (oxydes de soufre) que le gazole, car les huiles végétales ne contiennent pratiquement pas de soufre : il contribue donc à réduire la pollution de l'air et les pluies acides.
\item Il possède un \textbf{pouvoir lubrifiant supérieur} au gazole, ce qui protège la pompe d'injection et les injecteurs.
\end{itemize}
\end{propriete}

\courssub{Le biogaz : l'énergie cachée dans les déchets}

\textbf{Observation et problème.} Tout le monde a remarqué des bulles de gaz qui s'échappent de la vase au fond d'un marécage, ou l'odeur et le gonflement d'une fosse à fumier. Ce gaz, appelé autrefois « gaz des marais », est \textbf{inflammable}. \textbf{Problème scientifique :} d'où vient ce gaz, et peut-on le produire volontairement à partir de nos déchets pour s'en servir comme combustible ?

\textbf{Hypothèse.} En l'absence d'air, des micro-organismes dégradent la matière organique et libèrent un gaz combustible.

\begin{experience}[Mise en évidence de la production de biogaz]
\textbf{Matériel :} fumier frais, eau, un digesteur (bidon ou cuve hermétiquement fermé), un tube de dégagement, un robinet, une veilleuse ou une allumette, un thermomètre.\par
\textbf{Protocole :}
\begin{enumerate}
\item Introduire dans le digesteur un mélange de fumier et d'eau (environ moitié-moitié).
\item Fermer hermétiquement le digesteur (conditions \textbf{anaérobies}, c'est-à-dire sans oxygène) et le maintenir à une température voisine de \SI{35}{\celsius} (méthanisation mésophile).
\item Recueillir le gaz dégagé par le tube dans une poche ou une cloche à gaz.
\item Après quelques jours, approcher une flamme de la sortie du gaz (avec les précautions d'usage : absence de fuite, aération).
\end{enumerate}
\textbf{Observations :} un gaz s'accumule progressivement ; il \textbf{s'enflamme} avec une flamme bleutée, semblable à celle du gaz naturel.\par
\textbf{Interprétation :} en l'absence d'oxygène, des \textbf{bactéries méthanogènes} (archées méthanogènes, terminologie admise « bactéries » au programme) présentes dans le fumier dégradent la matière organique et produisent un mélange gazeux riche en \cle{méthane} \ce{CH4}, gaz inflammable. C'est la \textbf{méthanisation}.\par
\textbf{Conclusion :} les déchets organiques constituent une source d'énergie renouvelable, exploitable à l'échelle d'une ferme ou d'un établissement scolaire.
\end{experience}

\begin{definition}[Méthanisation]
La \cle{méthanisation} (ou digestion anaérobie) est la \textbf{dégradation de la matière organique} (fumier, déchets de cuisine, eaux usées, résidus de récolte) par des \textbf{bactéries méthanogènes}, en \textbf{absence d'oxygène} (anaérobiose). Elle produit :
\begin{itemize}
\item le \cle{biogaz}, mélange gazeux combustible contenant typiquement \textbf{50 à 70 \% de méthane} \ce{CH4} et \textbf{30 à 50 \% de dioxyde de carbone} \ce{CO2}, avec des traces d'autres gaz (\ce{H2S}, \ce{H2O}) ;
\item le \cle{digestat}, résidu solide et liquide riche en éléments minéraux (azote, phosphore, potassium), utilisable comme \textbf{engrais} organique.
\end{itemize}
\emph{Remarque :} la composition exacte du biogaz et le rôle des différentes souches bactériennes (bactéries hydrolytiques, acétogènes puis méthanogènes, qui agissent en relais) sont admis sans démonstration au programme de Terminale D.
\end{definition}

La combustion du biogaz libère l'énergie stockée dans le méthane :
\[ \ce{CH4 + 2 O2 -> CO2 + 2 H2O} + \text{énergie} \]

\begin{popfigure}
\begin{center}
\begin{tikzpicture}[scale=1.0, every node/.style={font=\small}]
% Cuve du digesteur
\draw[very thick, popdark, fill=popcream] (0,0) ellipse (3.2 and 0.6);
\draw[very thick, popdark] (-3.2,0) -- (-3.2,-2.6);
\draw[very thick, popdark] (3.2,0) -- (3.2,-2.6);
\draw[very thick, popdark, fill=popblueL] (-3.2,-2.6) arc (180:360:3.2 and 0.7);
% Contenu liquide
\fill[popblue!20] (-3.05,-0.7) -- (-3.05,-2.55) arc (180:360:3.05 and 0.65) -- (3.05,-0.7) arc (0:180:3.05 and 0.55);
\draw[popblue] (-3.05,-0.7) arc (180:360:3.05 and 0.55);
\node[popblue] at (0,-1.7) {boue : fumier + eau};
\node[popblue] at (0,-2.25) {(bactéries méthanogènes, sans \ce{O2})};
% Zone gaz
\node[popgreen] at (0,-0.15) {biogaz (\ce{CH4} + \ce{CO2})};
% Dôme
\draw[very thick, popdark, fill=popgreenL] (-3.2,0) arc (180:0:3.2 and 1.5);
% Entrée
\draw[very thick, poporange] (-4.6,1.2) -- (-4.6,-0.9) -- (-3.2,-0.9);
\draw[-{Stealth[length=3mm]}, very thick, poporange] (-4.6,2.0) -- (-4.6,1.2);
\node[poporange, align=center] at (-4.6,2.5) {entrée des\\déchets organiques};
% Sortie gaz
\draw[very thick, popgreen] (0,1.5) -- (0,2.3) -- (2.2,2.3);
\draw[-{Stealth[length=3mm]}, very thick, popgreen] (2.2,2.3) -- (3.4,2.3);
\node[popgreen, align=left] at (4.9,2.3) {sortie du biogaz\\(cuisinière, lampe, groupe électrogène)};
% Sortie digestat
\draw[very thick, poppurple] (2.0,-2.9) -- (2.0,-3.8) -- (3.4,-3.8);
\draw[-{Stealth[length=3mm]}, very thick, poppurple] (3.4,-3.8) -- (4.6,-3.8);
\node[poppurple, align=left] at (6.1,-3.8) {sortie du digestat\\(engrais)};
% Label digesteur
\node[popdark, font=\small\bfseries] at (0,1.0) {chambre de fermentation};
\end{tikzpicture}
\end{center}
\legende{Schéma d'un digesteur à biogaz domestique. Les déchets organiques mélangés à l'eau entrent dans la chambre de fermentation hermétique (anaérobie) ; les bactéries méthanogènes y produisent le biogaz, recueilli sous le dôme et conduit vers une cuisinière, une lampe ou un groupe électrogène ; le résidu (digestat) est évacué et utilisé comme engrais.}
\end{popfigure}

\begin{propriete}[Le double intérêt du biogaz]
\begin{enumerate}
\item \textbf{Intérêt énergétique :} le biogaz sert à la \textbf{cuisson}, à l'\textbf{éclairage} et même à la production d'électricité (cogénération), remplaçant le bois de chauffe (ce qui lutte contre la déforestation) et les combustibles fossiles.
\item \textbf{Intérêt sanitaire et agronomique (assainissement) :} la méthanisation \textbf{réduit le volume des déchets} et détruit une grande partie des \textbf{agents pathogènes} (bactéries, œufs de parasites) présents dans les fumiers et les eaux usées. De plus, le \textbf{digestat} constitue un engrais organique de qualité qui améliore les rendements des cultures vivrières.
\end{enumerate}
\end{propriete}

\begin{exemplebox}[Un exemple chiffré]
Une tonne (\SI{1000}{kg}) de déchets organiques peut produire environ \textbf{\SIrange{80}{150}{\cubic\meter} de biogaz} selon sa nature. À raison de 60 \% de méthane, cela représente \SIrange{50}{90}{\cubic\meter} de méthane, de quoi alimenter la cuisinière d'une famille pendant plusieurs semaines, tout en produisant environ \SI{800}{kg} de digestat fertilisant.
\end{exemplebox}

\begin{methode}[Distinguer les trois biocarburants]
Pour ne jamais les confondre, raisonne en deux questions :
\begin{enumerate}
\item \textbf{Quelle est la matière première ?} Sucres ou amidon (canne, manioc, maïs) ; huiles végétales (palme, coton, jatropha) ; déchets organiques (fumier, eaux usées).
\item \textbf{Quel est le procédé ?} Fermentation alcoolique par des levures $\rightarrow$ \textbf{bioéthanol} ; transestérification chimique (méthanol + soude) $\rightarrow$ \textbf{biodiesel} ; méthanisation par des bactéries anaérobies $\rightarrow$ \textbf{biogaz}.
\end{enumerate}
Retiens le triptyque : \emph{sucre fermenté = éthanol ; huile transformée = biodiesel ; déchet digéré = biogaz}.
\end{methode}

\begin{attention}[Erreur fréquente : biodiesel ou biogaz ?]
Ne confonds pas \textbf{biodiesel} et \textbf{biogaz} : leurs noms se ressemblent mais tout les oppose. Le \textbf{biodiesel} est un \emph{liquide} obtenu par transformation chimique des \emph{huiles végétales}, utilisé dans les \emph{moteurs diesel}. Le \textbf{biogaz} est un \emph{gaz} (méthane) obtenu par dégradation biologique des \emph{déchets organiques}, utilisé pour la \emph{cuisson, l'éclairage ou la production d'électricité}. Confondre les deux dans une copie de Bac est une erreur éliminatoire pour l'exercice : vérifie toujours matière première et procédé avant de répondre.
\end{attention}

\begin{savaistu}[Le biogaz, déjà une réalité camerounaise]
Plusieurs fermes et lycées techniques du Cameroun expérimentent déjà des digesteurs à biogaz : le fumier des porcheries et des étables y est transformé en gaz de cuisson pour les cantines scolaires, réduisant la facture de bois et de gaz butane. Dans les zones rurales, des projets appuyés par des ONG et le programme PNDP (Programme National de Développement Participatif) équipent des exploitations familiales de petits digesteurs en maçonnerie ou en bâche souple : une technologie simple, locale, qui transforme un problème d'assainissement en solution énergétique.
\end{savaistu}

\begin{exoresolu}[Bilan d'une installation de biogaz]
\textbf{Énoncé.} Une ferme de la région de l'Ouest produit \SI{500}{kg} de fumier par semaine. On admet qu'une tonne de déchets organiques donne en moyenne \SI{100}{m^3} de biogaz contenant 60 \% de méthane. (a) Quel volume de biogaz la ferme produit-elle par semaine ? (b) Quel volume de méthane cela représente-t-il ? (c) Citer deux avantages de cette installation pour la ferme.
\tcblower
\textbf{Solution détaillée.}\par
(a) \SI{500}{kg} = \SI{0,5}{tonne}. Volume de biogaz : $V = 0,5 \times 100 = \SI{50}{m^3}$ de biogaz par semaine.\par
(b) Volume de méthane : $V_{\ce{CH4}} = 50 \times \dfrac{60}{100} = \SI{30}{m^3}$ de méthane par semaine.\par
(c) Avantages : production d'une \textbf{énergie de cuisson gratuite et renouvelable} (économie de bois et de gaz butane) ; \textbf{assainissement} de la ferme (réduction des déchets et des agents pathogènes) ; obtention d'un \textbf{digestat utilisé comme engrais} pour les cultures. (Deux réponses suffisent.)
\end{exoresolu}

\begin{aretenir}[L'essentiel de la section]
\begin{itemize}
\item Le \textbf{biodiesel} est obtenu par \textbf{transestérification} des huiles végétales (palme, coton, jatropha) avec du méthanol en présence de soude ; la glycérine est le coproduit. Il est utilisable dans les moteurs diesel classiques, biodégradable et pauvre en fumées sulfurées.
\item Le \textbf{biogaz} est produit par \textbf{méthanisation} : dégradation anaérobie des déchets organiques par des bactéries méthanogènes. Il contient 50 à 70 \% de méthane et sert à la cuisson, l'éclairage et la production d'électricité ; le digestat est un engrais.
\item Le biogaz a un \textbf{double intérêt} : énergétique et sanitaire (assainissement).
\item Triptyque à retenir : fermentation $\rightarrow$ bioéthanol ; transestérification $\rightarrow$ biodiesel ; méthanisation $\rightarrow$ biogaz.
\end{itemize}
\end{aretenir}

\begin{exercice}[Pour t'entraîner]
Une coopérative hésite entre installer une unité de biodiesel à base de jatropha et un digesteur à biogaz alimenté par les déchets du marché. Pour chaque option, indique la matière première, le procédé biologique ou chimique mis en jeu, le produit énergétique obtenu et un avantage spécifique. Laquelle des deux options contribue le plus à l'assainissement du quartier ? Justifie.
\end{exercice}

\begin{corrige}[Exercice]
\textbf{Unité biodiesel :} matière première = graines de jatropha (huile non alimentaire) ; procédé = transestérification (méthanol + soude) ; produit = esters méthyliques (biodiesel liquide) ; avantage = carburant pour moteurs diesel, biodégradable, sans compétition avec l'alimentation.\par
\textbf{Digesteur à biogaz :} matière première = déchets organiques du marché ; procédé = méthanisation par bactéries anaérobies ; produit = biogaz (\ce{CH4} + \ce{CO2}) ; avantage = énergie de cuisson et réduction des déchets.\par
C'est le \textbf{digesteur à biogaz} qui contribue le plus à l'assainissement : il élimine les déchets du marché, détruit une partie des agents pathogènes et fournit en outre un engrais (digestat).
\end{corrige}

\coursec{Biocarburants vs énergies fossiles : avantages et limites}

Dans la section précédente, nous avons vu que le biodiesel et le biogaz complètent le bioéthanol dans la famille des biocarburants : l'un est produit par transestérification des huiles végétales (palme, coton, jatropha), l'autre par méthanisation des déchets organiques. Ces trois filières partagent une même promesse : transformer la \cle{biomasse} en énergie utilisable. Mais une promesse n'est pas une preuve. Face au pétrole, qui alimente aujourd'hui l'essentiel des transports et des groupes électrogènes camerounais, les biocarburants sont-ils réellement avantageux ? Pour répondre, il faut comparer méthodiquement les deux familles d'énergie, critère par critère, sans enthousiasme naïf ni rejet systématique. C'est l'objet de cette section : construire un jugement scientifique équilibré, fondé sur des données chiffrées et sur l'analyse du \cle{cycle de vie} complet des carburants.

\courssub{Les avantages des biocarburants}

\textbf{Une ressource renouvelable.} Le pétrole s'est formé il y a des millions d'années à partir de plancton fossile : à l'échelle humaine, il est \cle{non renouvelable}. Chaque litre brûlé est un litre perdu définitivement, et les réserves mondiales s'épuisent. À l'inverse, la canne à sucre, le manioc ou le palmier à huile se ressèment chaque année : la biomasse est \cle{renouvelable} car sa production repose sur la photosynthèse, un processus alimenté en permanence par l'énergie solaire. Tant que le soleil brille et que des terres sont cultivées, la matière première peut être régénérée.

\textbf{Un bilan carbone allégé.} C'est l'argument le plus souvent avancé, et il mérite d'être compris précisément. Lorsqu'on brûle de l'essence fossile, on rejette dans l'atmosphère du \ce{CO2} qui était séquestré sous terre depuis des millions d'années : c'est un ajout net de carbone au cycle actuel. Lorsqu'on brûle du bioéthanol, le \ce{CO2} rejeté est celui que la plante a \emph{récemment} pompé dans l'atmosphère par photosynthèse pour fabriquer sa matière organique. En première approximation, le cycle est bouclé :
\[
\ce{6CO2 + 6H2O ->[$\text{photosynthèse}$] C6H12O6 + 6O2} \quad\text{puis}\quad \ce{C6H12O6 ->[$\text{fermentation}$] 2C2H5OH + 2CO2}
\]
et la combustion de l'éthanol restitue le \ce{CO2} absorbé. On parle de bilan carbone \emph{quasi neutre} — et non pas nul, nous verrons pourquoi dans les limites.

\textbf{Une réduction des importations et des emplois ruraux.} Le Cameroun, bien que producteur de pétrole, importe une partie de ses produits raffinés, ce qui pèse lourdement sur la balance commerciale et expose le pays aux fluctuations des cours mondiaux. Produire du bioéthanol à partir de la canne à sucre de SOSUCAM ou du manioc paysan, c'est substituer une ressource nationale à une importation. De plus, les filières biocarburants sont intensives en main-d'œuvre agricole : plantation, récolte, transport, distillation artisanale ou industrielle créent des emplois en zone rurale, là où l'exode vers Yaoundé et Douala s'accélère.

\textbf{Une valorisation des déchets.} Le biogaz illustre parfaitement cet avantage : les résidus de récolte, les effluents d'élevage, les déchets de marché et les eaux usées, qui posent de graves problèmes sanitaires dans nos villes, deviennent une matière première énergétique au lieu d'être une source de pollution. Valoriser un déchet, c'est résoudre deux problèmes à la fois.

\textbf{Une moindre émission de soufre.} Les biocarburants ne contiennent pratiquement pas de soufre, contrairement aux carburants fossiles dont la combustion dégage des oxydes de soufre (\ce{SO2}), responsables des pluies acides et de pathologies respiratoires. Leur combustion émet aussi moins de particules fines et de monoxyde de carbone incomplet.

\begin{attention}[Renouvelable ne signifie pas sans impact]
Erreur fréquente au Baccalauréat : écrire « les biocarburants sont écologiques car renouvelables ». C'est un raisonnement incomplet. Un carburant renouvelable peut avoir un impact environnemental et social désastreux s'il provoque déforestation, pénurie alimentaire ou pollution par les engrais. Il faut toujours raisonner sur \cle{tout le cycle de vie} : culture (engrais, pesticides, eau), transformation (énergie de distillation), transport, combustion. Le mot « renouvelable » qualifie la ressource, pas l'impact global.
\end{attention}

\courssub{Limite 1 : la compétition alimentaire, le débat « assiette contre réservoir »}

\textbf{Observation.} En 2007-2008, les prix mondiaux du maïs et du blé ont fortement augmenté, en partie parce qu'une fraction croissante de la récolte de maïs américaine était détournée vers la production d'éthanol. Des émeutes de la faim ont éclaté dans plusieurs pays.

\textbf{Interprétation.} Les biocarburants de première génération utilisent précisément les organes que l'homme consomme : le grain de maïs, le tubercule de manioc, le sucre de la canne, l'huile de palme. Or l'huile de palme est un aliment de base au Cameroun, et le manioc est l'aliment principal de millions de familles. Détourner ces ressources vers les réservoirs crée une \cle{compétition alimentaire} : l'offre alimentaire diminue, les prix montent, et ce sont les populations les plus pauvres qui en souffrent les premières. C'est le fameux dilemme « \emph{l'assiette contre le réservoir} ».

\courssub{Limite 2 : l'usage des terres et la déforestation}

Pour produire massivement des biocarburants, il faut de vastes surfaces cultivées. Deux scénarios posent problème :

\begin{itemize}
\item \textbf{La déforestation directe} : pour étendre les palmeraies industrielles, on abat des forêts tropicales. Or la forêt camerounaise (bassin du Congo, deuxième massif forestier mondial) est un immense réservoir de carbone et de biodiversité. La détruire libère d'un coup le carbone stocké dans les arbres et le sol : il faudrait alors des \emph{dizaines d'années} de production de biodiesel pour « rembourser » cette dette carbone. Le bénéfice climatique est annulé, voire inversé.
\item \textbf{Le changement d'affectation des terres} : même sans déforestation directe, convertir des terres vivrières en plantations énergétiques déplace les cultures alimentaires vers de nouvelles terres... qu'il faut défricher ailleurs. C'est l'effet de cascade, difficile à contrôler.
\end{itemize}

\courssub{Limite 3 : le rendement énergétique}

\textbf{Deux problèmes distincts} se cachent derrière cette limite.

\emph{Premier problème : l'énergie investie dans la production.} Cultiver la canne ou le maïs exige des engrais (fabriqués à partir de gaz fossile), des machines agricoles (qui consomment du gasoil), de l'irrigation, puis le transport et surtout la \cle{distillation}, étape très énergivore qui concentre l'éthanol après fermentation. Si l'on compte toute l'énergie fossile investie, le \emph{gain net} d'énergie peut être faible : pour certaines filières de maïs-éthanol, on investit presque autant d'énergie fossile qu'on en récupère dans le bioéthanol. La canne à sucre fait mieux, car ses résidus (la bagasse) peuvent être brûlés pour alimenter la distillerie.

\emph{Second problème : le pouvoir calorifique de l'éthanol.}

\begin{propriete}[Pouvoirs calorifiques volumiques comparés (valeurs admises, à connaître pour le Bac)]
Le \cle{pouvoir calorifique volumique} d'un carburant est l'énergie libérée par la combustion d'un litre de ce carburant. Le pouvoir calorifique volumique de l'éthanol ($\approx \SI{21}{MJ/L}$) est inférieur à celui de l'essence ($\approx \SI{32}{MJ/L}$). À volume égal, l'éthanol fournit donc environ un tiers d'énergie en moins : un véhicule roulant à l'éthanol pur consomme davantage de litres aux 100 km. Le biodiesel ($\approx \SI{33}{MJ/L}$) est proche du diesel fossile ($\approx \SI{36}{MJ/L}$).
\end{propriete}

\begin{popfigure}
\begin{center}
\begin{tikzpicture}
\begin{axis}[
  ybar, bar width=0.9cm,
  width=0.85\linewidth, height=6.5cm,
  ylabel={Pouvoir calorifique volumique (MJ/L)},
  symbolic x coords={Éthanol, Biodiesel, Essence, Diesel},
  xtick=data,
  ymin=0, ymax=40,
  ytick={0,10,20,30,40},
  nodes near coords,
  every node near coord/.append style={font=\small},
  axis line style={popdark},
  tick label style={font=\small},
]
\addplot[fill=popgreen!70, draw=popdark] coordinates {(Éthanol,21) (Biodiesel,33) (Essence,32) (Diesel,36)};
\end{axis}
\end{tikzpicture}
\end{center}
\legende{Pouvoirs calorifiques volumiques approximatifs de quatre carburants (en MJ/L). L'éthanol ($\approx 21$ MJ/L) est nettement moins énergétique que l'essence ($\approx 32$ MJ/L) et le diesel ($\approx 36$ MJ/L) ; le biodiesel ($\approx 33$ MJ/L) est proche du diesel fossile.}
\end{popfigure}

\courssub{Limite 4 : les coûts et les infrastructures}

Produire du bioéthanol à grande échelle exige des distilleries, des cuves de fermentation, des chaînes de déshydratation : des investissements lourds, difficiles à financer dans les économies fragiles. Côté utilisation, les moteurs à essence classiques tolèrent sans modification de faibles taux de mélange (5 à 10 \% d'éthanol, carburants E5/E10), mais les forts taux (E85) exigent des moteurs adaptés dits « flex-fuel », car l'éthanol est corrosif pour certains métaux et élastomères. Il faut aussi organiser toute une logistique de distribution (stations, stockage), car l'éthanol absorbe l'eau et ne se transporte pas par les oléoducs classiques.

\courssub{Tableau comparatif synthétique : biocarburants contre pétrole}

\begin{popfigure}
\begin{center}
\begin{tikzpicture}
\fill[popblue!15] (0,0) rectangle (10.5,1);
\fill[popgreen!12] (0,-1) rectangle (10.5,0);
\fill[poporange!12] (0,-2) rectangle (10.5,-1);
\fill[popgreen!12] (0,-3) rectangle (10.5,-2);
\fill[poporange!12] (0,-4) rectangle (10.5,-3);
\fill[popgreen!12] (0,-5) rectangle (10.5,-4);
\draw[popdark, thick] (0,1) rectangle (10.5,-5);
\draw[popdark] (3.2,1) -- (3.2,-5);
\draw[popdark] (6.9,1) -- (6.9,-5);
\foreach \y in {0,-1,-2,-3,-4} {\draw[popdark] (0,\y) -- (10.5,\y);}
\node[font=\bfseries\small] at (1.6,0.5) {Critère};
\node[font=\bfseries\small] at (5.05,0.5) {Biocarburants};
\node[font=\bfseries\small] at (8.7,0.5) {Énergies fossiles};
\node[font=\small, text width=3cm, align=center] at (1.6,-0.5) {Renouvelabilité};
\node[font=\small, text width=3.5cm, align=center] at (5.05,-0.5) {Oui (cultures annuelles)};
\node[font=\small, text width=3.4cm, align=center] at (8.7,-0.5) {Non (réserves épuisables)};
\node[font=\small, text width=3cm, align=center] at (1.6,-1.5) {Émissions de \ce{CO2}};
\node[font=\small, text width=3.5cm, align=center] at (5.05,-1.5) {Bilan quasi neutre si pas de déforestation};
\node[font=\small, text width=3.4cm, align=center] at (8.7,-1.5) {Ajout net de carbone fossile};
\node[font=\small, text width=3cm, align=center] at (1.6,-2.5) {Coût};
\node[font=\small, text width=3.5cm, align=center] at (5.05,-2.5) {Élevé (distillation, infrastructures)};
\node[font=\small, text width=3.4cm, align=center] at (8.7,-2.5) {Filière mature, mais prix volatils};
\node[font=\small, text width=3cm, align=center] at (1.6,-3.5) {Disponibilité};
\node[font=\small, text width=3.5cm, align=center] at (5.05,-3.5) {Locale, mais limitée par les terres};
\node[font=\small, text width=3.4cm, align=center] at (8.7,-3.5) {Gisements inégalement répartis, importations};
\node[font=\small, text width=3cm, align=center] at (1.6,-4.5) {Impact social};
\node[font=\small, text width=3.5cm, align=center] at (5.05,-4.5) {Emplois ruraux, mais compétition alimentaire};
\node[font=\small, text width=3.4cm, align=center] at (8.7,-4.5) {Rentabilité concentrée, pollution urbaine};
\end{tikzpicture}
\end{center}
\legende{Comparaison synthétique des biocarburants et des énergies fossiles selon cinq critères : renouvelabilité, émissions de \ce{CO2}, coût, disponibilité et impact social.}
\end{popfigure}

\begin{methode}[Construire un argumentaire équilibré]
Face à une question du type « Les biocarburants sont-ils une solution au déficit énergétique du Cameroun ? », procède en quatre étapes :
\begin{enumerate}
\item \textbf{Choisir au moins trois critères} croisés : économique (coût, emplois, importations), écologique (bilan carbone, déforestation, pollution), social (alimentation, accès à l'énergie).
\item \textbf{Présenter un argument et son contre-argument} pour chaque critère (ex. bilan carbone allégé \emph{mais} déforestation possible).
\item \textbf{Appuyer chaque argument sur une donnée chiffrée} (pouvoirs calorifiques, surfaces, prix) ou un exemple local (SOSUCAM, palmeraies, biogaz de décharge).
\item \textbf{Conclure par une position nuancée} : les biocarburants sont un complément pertinent, pas une solution miracle ; leur intérêt dépend des conditions de production.
\end{enumerate}
\end{methode}

\courssub{Les pistes de solution : vers les biocarburants de deuxième génération}

Les limites précédentes ne condamnent pas les biocarburants : elles désignent les conditions de leur durabilité. Deux pistes majeures émergent.

\textbf{Les biocarburants de deuxième génération.} Au lieu d'utiliser les organes comestibles (grain, tubercule, sucre), on valorise les \cle{résidus lignocellulosiques} non alimentaires : pailles, tiges de maïs, bagasse de canne, coques de cacao, sciures. Des traitements enzymatiques (cellulases) décomposent la cellulose en glucose, ensuite fermenté en éthanol. Ainsi, aucune terre supplémentaire n'est nécessaire et la compétition alimentaire disparaît. Au Cameroun, les immenses quantités de résidus de cacao, de café et de bananeraies constituent un gisement encore inexploité.

\textbf{Les cultures sur terres marginales.} Certaines plantes oléagineuses poussent sur des sols pauvres, impropres aux cultures vivrières, et ne concurrencent donc pas l'alimentation.

\begin{savaistu}[Le jatropha, la plante qui ne vole pas la terre des vivriers]
Le \cle{jatropha} (\emph{Jatropha curcas}) est un arbuste dont les graines, toxiques et non comestibles, contiennent une huile transformable en biodiesel. Il pousse sur des terres arides ou dégradées, sans irrigation ni engrais abondants, et sert même de haie vive protégeant les champs contre le bétail. Au Cameroun, des projets de plantations de jatropha ont été promus dans les années 2000-2010 précisément parce que cette culture \emph{n'entre pas en compétition avec les cultures vivrières} : elle illustre la stratégie des terres marginales. Les résultats économiques ont toutefois été mitigés, rappelant qu'une bonne idée biologique doit aussi être validée agronomiquement et commercialement.
\end{savaistu}

\begin{exoresolu}[Comparer deux carburants pour un parc automobile]
Une coopérative de transport de l'Ouest camerounais hésite entre l'essence ($\SI{32}{MJ/L}$) et un mélange E85 contenant 85 \% d'éthanol ($\SI{21}{MJ/L}$) et 15 \% d'essence. Un véhicule parcourt 600 km avec 48 L d'essence. (a) Calculer l'énergie fournie par ces 48 L d'essence. (b) Calculer le pouvoir calorifique du mélange E85, puis le volume d'E85 nécessaire pour fournir la même énergie. (c) Citer un avantage et une limite de ce choix pour la coopérative.
\tcblower
(a) Énergie fournie par l'essence :
\[ E = 48 \times 32 = \SI{1536}{MJ}. \]
(b) Pouvoir calorifique du mélange E85 (moyenne pondérée en volume, hypothèse de volumes additifs) :
\[ P_{E85} = 0{,}85 \times 21 + 0{,}15 \times 32 = 17{,}85 + 4{,}8 = \SI{22,65}{MJ/L}. \]
Volume nécessaire pour fournir $\SI{1536}{MJ}$ :
\[ V = \frac{1536}{22{,}65} \approx \SI{67,8}{L}. \]
Il faut donc environ 68 L d'E85 contre 48 L d'essence, soit une surconsommation d'environ 41 \% en volume, conséquence directe du pouvoir calorifique plus faible de l'éthanol.
(c) \emph{Avantage} : le mélange utilise une ressource renouvelable locale (éthanol de canne ou de manioc), réduisant les importations et les émissions nettes de \ce{CO2}. \emph{Limite} : la surconsommation en volume peut annuler le gain économique, et l'usage du E85 exige des moteurs adaptés (flex-fuel), donc un investissement dans le renouvellement du parc.
\end{exoresolu}

\begin{aretenir}[L'essentiel de la section]
\begin{itemize}
\item Les biocarburants présentent quatre avantages majeurs : ressource \cle{renouvelable}, bilan carbone allégé (le \ce{CO2} rejeté a été récemment absorbé par la plante), réduction des importations et création d'emplois ruraux, valorisation des déchets, faible émission de soufre.
\item Leurs quatre limites : la \cle{compétition alimentaire} (« assiette contre réservoir »), la déforestation et l'usage des terres qui peuvent annuler le bénéfice carbone, un rendement énergétique réduit (énergie investie en culture et distillation ; éthanol à $\SI{21}{MJ/L}$ contre $\SI{32}{MJ/L}$ pour l'essence), et des coûts d'infrastructure élevés (distilleries, moteurs adaptés).
\item « Renouvelable » ne veut pas dire « sans impact » : tout jugement doit porter sur le \cle{cycle de vie} complet et croiser les critères économique, écologique et social.
\item Les solutions d'avenir : biocarburants de \cle{deuxième génération} à partir de résidus non alimentaires (lignocellulose) et cultures énergétiques sur terres marginales (jatropha).
\end{itemize}
\end{aretenir}

\coursec{Énergies renouvelables, éducation et citoyenneté}

Nous avons vu, dans la section précédente, que les biocarburants présentent des avantages réels (caractère renouvelable, bilan carbone atténué, valorisation des déchets) mais aussi des limites sérieuses (compétition avec les cultures vivrières, rendements énergétiques modestes, coûts de transformation). La conclusion qui s'impose n'est donc ni l'enthousiasme aveugle ni le rejet : les biocarburants sont \emph{une} pièce d'un puzzle plus vaste, le \cle{mix énergétique}. Cette dernière section replace les biocarburants dans le bouquet des énergies renouvelables, montre leur intérêt concret face aux déficits énergétiques et à la pollution au Cameroun, et te donne les outils pour devenir, toi-même, un acteur de la transition énergétique : car la biotechnologie n'est pas seulement une affaire de laboratoires, c'est aussi une affaire de citoyens.

\courssub{Les biocarburants dans le bouquet des énergies renouvelables}

\begin{definition}[Énergie renouvelable]
Une \cle{énergie renouvelable} est une énergie issue de sources naturelles dont le renouvellement est assez rapide pour être considéré comme inépuisable à l'échelle humaine : rayonnement solaire, vent, chute d'eau, chaleur interne de la Terre (géothermie) et \cle{biomasse} (matière organique végétale ou animale). Elle s'oppose aux énergies \cle{fossiles} (pétrole, charbon, gaz naturel), formées sur des millions d'années et dont les réserves s'épuisent.
\end{definition}

L'intuition est simple : aucune source renouvelable n'est parfaite, mais leurs défauts ne se produisent pas au même moment. Le soleil ne brille pas la nuit ; le vent est irrégulier ; les barrages hydroélectriques dépendent de la pluviométrie (les saisons sèches prolongées réduisent la production, ce que le Cameroun connaît bien) ; la biomasse demande des terres et du temps de culture. En revanche, la biomasse possède un atout unique parmi les renouvelables : elle est \cle{stockable}. Le bois, l'éthanol, le biodiesel et le biogaz peuvent être conservés et utilisés \emph{à la demande}, là où l'électricité solaire ou éolienne doit être consommée immédiatement ou stockée dans des batteries coûteuses.

\begin{propriete}[Complémentarité des énergies renouvelables]
Un \cle{mix énergétique} performant combine des sources aux profils complémentaires :
\begin{itemize}
\item l'\cle{hydroélectricité} fournit une production de base puissante et pilotable (on ouvre ou ferme les vannes) ;
\item le \cle{solaire photovoltaïque} produit en journée, de façon décentralisée, idéal pour l'éclairage et les petits usages ruraux ;
\item l'\cle{éolien} produit selon les régimes de vents (pertinent dans certaines zones du Nord-Cameroun et du littoral) ;
\item les \cle{biocarburants} fournissent des combustibles stockables pour la cuisson (biogaz, alcool à brûler), le transport (bioéthanol, biodiesel) et l'électrification d'appoint.
\end{itemize}
La complémentarité lisse les intermittences : quand l'une faiblit, l'autre prend le relais.
\end{propriete}

\begin{attention}[Erreur fréquente : le biocarburant « solution miracle »]
Il est faux d'affirmer que les biocarburants \emph{remplaceront} à eux seuls les énergies fossiles. Les surfaces cultivables sont limitées et entrent en compétition avec l'alimentation ; les rendements de conversion de la biomasse sont faibles (la photosynthèse ne convertit qu'environ 1 à 2 \% de l'énergie solaire reçue en biomasse). Au Baccalauréat, la formulation attendue est : les biocarburants \textbf{s'inscrivent dans un mix énergétique} aux côtés de l'hydroélectricité, du solaire et de l'éolien, et contribuent à \emph{réduire} — non à supprimer — la dépendance aux fossiles.
\end{attention}

\begin{savaistu}[Le géant hydroélectrique camerounais]
Le Cameroun dispose d'un potentiel hydroélectrique estimé à environ \SI{20}{GW}, l'un des plus importants d'Afrique (souvent cité au deuxième ou troisième rang continental), dont moins de 10 \% est actuellement exploité (barrages de Songloulou et d'Édéa sur la Sanaga, Lagdo sur la Bénoué, Memve'ele, Nachtigal). Paradoxe : malgré ce potentiel, une grande partie des ménages ruraux n'est pas raccordée au réseau et les \cle{délestages} (coupures tournantes) restent fréquents. C'est précisément là que les solutions décentralisées — biogaz villageois, petits systèmes solaires, alcool à brûler local — prennent tout leur sens.
\end{savaistu}

\courssub{Intérêt face aux déficits énergétiques : la décentralisation}

Observons une situation concrète : dans un village de la région de l'Ouest, non raccordé au réseau national, les ménages dépensent une part importante de leurs revenus en bois de chauffe, en pétrole lampant et en transport pour s'approvisionner. Or ce même village produit chaque jour des tonnes de déjections animales, de résidus de récolte et de déchets de cuisine — autant de \emph{matière première énergétique gratuite} laissée à l'abandon.

\begin{definition}[Production décentralisée d'énergie]
La \cle{production décentralisée} (ou distribuée) consiste à produire l'énergie au plus près du lieu de consommation, à petite échelle : digesteur à biogaz familial ou villageois, distillerie artisanale d'alcool à brûler, kits solaires domestiques. Elle s'oppose à la production centralisée (grands barrages, centrales thermiques) qui exige un réseau de transport coûteux.
\end{definition}

La démarche d'investigation est ici économique autant que biologique. Prenons le document suivant : une famille rurale consomme environ \SI{3}{m^3} de biogaz par jour pour la cuisson. Un digesteur familial bien dimensionné, alimenté par les déjections de 4 à 6 vaches (environ \SI{40}{kg} à \SI{60}{kg} de bouses fraîches par jour, diluées dans un volume équivalent d'eau), produit justement de l'ordre de \SI{2}{m^3} à \SI{4}{m^3} de biogaz par jour. La fermentation anaérobie (méthanisation) réalisée par le consortium bactérien du digesteur convertit la matière organique en un mélange d'environ 60 \% de méthane \ce{CH4} et 40 \% de \ce{CO2}, directement utilisable dans une cuisinière à gaz. Le \cle{digestat} résiduel, débarrassé d'une grande partie des germes pathogènes et des odeurs, constitue un excellent engrais organique pour les champs : rien ne se perd.

\begin{exemplebox}[Exemple chiffré : le digesteur familial]
Un digesteur familial alimenté par les déjections de \textbf{4 à 6 vaches} produit environ \SI{3}{m^3} de biogaz par jour, ce qui \textbf{couvre les besoins de cuisson d'une famille} de 5 à 8 personnes (soit l'équivalent d'environ \SI{1,5}{kg} de gaz butane ou de \SI{10}{kg} de bois par jour). Effets en cascade : économie d'achat de combustible, réduction de la corvée de bois (qui pèse surtout sur les femmes et les filles), diminution de la pression sur les forêts, et engrais gratuit pour le potager.
\end{exemplebox}

Les bénéfices collectifs de la décentralisation sont directs :
\begin{itemize}
\item \textbf{réduction des délestages} : chaque ménage ou village autonome en cuisson et en éclairage soulage d'autant le réseau national ;
\item \textbf{baisse des coûts pour les ménages} : la matière première (déjections, résidus de manioc, tiges de maïs) est locale et quasi gratuite ;
\item \textbf{création d'emplois ruraux} : fabrication et maintenance des digesteurs, petites distilleries d'éthanol à partir de canne ou de manioc ;
\item \textbf{sécurisation de l'approvisionnement} : moins de dépendance aux importations de carburants, dont les prix fluctuent.
\end{itemize}

\courssub{Intérêt face à la pollution et à la gestion des déchets}

La combustion du bois et du charbon de bois en milieu clos dégage des \cle{particules fines} et du monoxyde de carbone \ce{CO}, responsables d'infections respiratoires — un problème majeur de santé publique dans les cuisines traditionnelles. Les carburants fossiles, eux, émettent du dioxyde de soufre \ce{SO2} (pluies acides) et des oxydes d'azote. Les biocarburants brûlent plus proprement : le biogaz ne dégage pratiquement \textbf{pas de particules ni de \ce{SO2}}, et l'éthanol, molécule oxygénée, permet une combustion plus complète que l'essence :
\[ \ce{C2H5OH + 3O2 -> 2CO2 + 3H2O} \]

Le deuxième bénéfice est sanitaire et urbain. À Yaoundé ou Douala, les déchets organiques (épluchures, marchés, abattoirs) représentent souvent plus de la moitié des ordures ménagères ; abandonnés, ils fermentent à l'air libre, dégagent du méthane (un gaz à effet de serre environ 25 fois plus puissant que le \ce{CO2} sur 100 ans) et attirent les vecteurs de maladies. La méthanisation \emph{capture} ce méthane pour en faire un combustible utile au lieu de le laisser s'échapper : on transforme un problème d'assainissement en ressource énergétique. De même, à la campagne, les résidus agricoles (pailles, fanes, rafles) trouvent une valorisation au lieu d'être brûlés à l'air libre.

\begin{attention}[Nuance indispensable]
« Renouvelable » ne signifie pas « zéro impact ». La combustion des biocarburants émet du \ce{CO2} (partiellement compensé par la photosynthèse des cultures), et une méthanisation mal conduite laisse fuir du méthane. Le raisonnement correct au Bac compare toujours les \emph{bilans globaux} (du champ ou du déchet jusqu'à la combustion), jamais la seule étape finale.
\end{attention}

\courssub{Le rôle du citoyen et de l'élève : agir et informer}

La transition énergétique commence par des gestes et des choix collectifs. En tant qu'élève de Terminale D, tu disposes déjà des connaissances scientifiques pour être un relais crédible dans ta famille et ton quartier.

\begin{itemize}
\item \textbf{Trier et valoriser les déchets organiques} : composter les épluchures ou alimenter un digesteur plutôt que brûler ou enfouir.
\item \textbf{Promouvoir le biogaz domestique} : expliquer autour de soi le principe de la méthanisation, identifier les exploitations (élevages, lycées agricoles, marchés) où un digesteur serait rentable.
\item \textbf{Consommer de façon responsable} : éteindre les appareils, préférer les foyers améliorés, soutenir les filières locales (alcool à brûler, charbon durable) plutôt que la coupe anarchique du bois.
\item \textbf{S'informer et informer} : vérifier les chiffres, dénoncer les idées reçues (« les biocarburants vont nous affamer » est un raccourci ; tout dépend des filières choisies).
\end{itemize}

\begin{methode}[Concevoir une affiche ou un dépliant de sensibilisation]
Pour élaborer un message efficace, adapté à un public donné (élèves, parents, agriculteurs), suis quatre étapes :
\begin{enumerate}
\item \textbf{Un problème} clairement identifié et proche du public visé. Exemple pour des parents : « Le bois de chauffe coûte de plus en plus cher et la fumée rend les enfants malades. »
\item \textbf{Une solution} concrète et réaliste localement : « Un digesteur à biogaz transforme les bouses de vos vaches en gaz de cuisson gratuit. »
\item \textbf{Un chiffre clé}, simple et mémorable : « 4 vaches = le gaz de cuisson de toute la famille. »
\item \textbf{Un appel à l'action} précis : « Contactez le lycée agricole ou l'antenne MINEPIA de votre arrondissement pour un montage de digesteur. »
\end{enumerate}
Adapte le vocabulaire au public : images et phrases courtes pour les élèves, arguments économiques pour les parents, rendements et engrais pour les agriculteurs. Évite le jargon scientifique non expliqué.
\end{methode}

\begin{exoresolu}[Exercice résolu : dimensionner et argumenter]
\textbf{Énoncé.} Un village de 50 familles non électrifiées élève en moyenne 5 vaches par famille. On estime qu'un digesteur familial alimenté par les déjections de 4 à 6 vaches produit \SI{3}{m^3} de biogaz par jour, suffisant pour la cuisson d'une famille, et que \SI{1}{m^3} de biogaz remplace environ \SI{3,5}{kg} de bois de chauffe. (a) Le cheptel du village peut-il couvrir les besoins de cuisson de toutes les familles ? (b) Quelle masse de bois le village économiserait-il par an ? (c) Rédige le « chiffre clé » d'une affiche destinée aux éleveurs.
\tcblower
\textbf{Solution.} (a) Chaque famille possède 5 vaches, soit un effectif compris dans la fourchette 4 à 6 vaches requise par digesteur. Chaque famille peut donc alimenter son propre digesteur et couvrir ses besoins de cuisson : \textbf{le cheptel suffit} pour les 50 familles.

(b) Production villageoise : $50 \times \SI{3}{m^3/jour} = \SI{150}{m^3/jour}$. Bois remplacé par jour : $150 \times 3,5 = \SI{525}{kg/jour}$. Sur un an : $525 \times 365 \approx \SI{191625}{kg}$, soit environ \textbf{192 tonnes de bois économisées par an}, autant de forêt épargnée et de dépenses évitées.

(c) Chiffre clé possible : « \textbf{Vos 5 vaches cuisinent pour vous : zéro franc de bois, 10 kg de bois économisés chaque jour !} » — un problème (coût du bois), une solution (digesteur), un chiffre, un appel implicite à l'action.
\end{exoresolu}

\courssub{Synthèse de la leçon : un village énergétiquement autonome}

Imaginons le village idéal qui combine les trois filières étudiées : un champ de canne à sucre (ou de manioc) alimente une petite unité de production d'éthanol ; un digesteur collectif méthanise les déjections animales et les déchets du marché ; des panneaux solaires équipent l'école et le centre de santé. Chaque filière compense les faiblesses des autres : le solaire produit le jour, le biogaz cuit matin et soir, l'éthanol stocke l'énergie pour les mobylettes et la motopompe.

\begin{popfigure}
\begin{center}
\begin{tikzpicture}[scale=0.95, every node/.style={transform shape},
  maison/.style={draw=popdark, fill=popcream, rounded corners=2pt, minimum width=1.5cm, minimum height=0.9cm, font=\scriptsize, align=center},
  source/.style={draw=popdark, rounded corners=3pt, minimum width=2.6cm, minimum height=1.1cm, font=\scriptsize, align=center}]
  % Soleil
  \node[circle, draw=popgold, fill=popgoldL, minimum size=1.1cm, font=\scriptsize] (soleil) at (-6,3.4) {Soleil};
  \foreach \a in {0,45,...,315} \draw[popgold, thick] (-6,3.4) ++(\a:0.62) -- ++(\a:0.3);
  % Panneaux solaires
  \node[source, fill=popblueL, align=center] (solar) at (-6,1.4){Panneaux\\solaires};
  \draw[->, thick, popblue] (soleil) -- (solar);
  % Champ de canne
  \node[source, fill=popgreenL, align=center] (canne) at (0,3.2){Champ de canne\\(ou manioc)};
  \node[source, fill=popgreenL!70, align=center] (ethanol) at (0,1.2){Distillerie\\bioéthanol};
  \draw[->, thick, popgreen] (canne) -- (ethanol);
  % Digesteur
  \node[source, fill=poporangeL, align=center] (dechets) at (6,3.2){Déjections animales\\+ déchets organiques};
  \node[source, fill=poporangeL!80, align=center] (digesteur) at (6,1.2){Digesteur\\à biogaz};
  \draw[->, thick, poporange] (dechets) -- (digesteur);
  % Usages
  \node[maison, align=center] (ecole) at (-6,-1.2){École et centre\\de santé};
  \node[maison, align=center] (cuisine) at (0,-1.2){Cuisines\\des familles};
  \node[maison, align=center] (moteurs) at (3,-1.2){Mobylettes,\\motopompe};
  \node[maison, align=center] (champs) at (6,-1.2){Engrais\\(digestat)};
  \draw[->, thick, popblue] (solar) -- (ecole);
  \draw[->, thick, popblue] (solar.south) .. controls (-3,0) .. (cuisine.north);
  \draw[->, thick, popgreen] (ethanol) -- (moteurs);
  \draw[->, thick, poporange] (digesteur.south) -- (cuisine.north east);
  \draw[->, thick, poporange] (digesteur) -- (champs);
  % Boucle engrais -> champ
  \draw[->, thick, poporange, dashed] (champs.north) .. controls (8.6,0.4) and (8.6,3.2) .. (canne.east);
\end{tikzpicture}
\end{center}
\legende{Schéma fonctionnel d'un village énergétiquement autonome. Trois filières complémentaires : solaire photovoltaïque (éclairage, petits usages), fermentation alcoolique de la canne ou du manioc (bioéthanol stockable pour les moteurs), méthanisation des déjections et déchets (biogaz de cuisson). Le digestat, engrais organique, retourne aux champs (flèche pointillée) : la boucle matière est bouclée.}
\end{popfigure}

Pour fixer l'ensemble de la leçon, voici la carte mentale qui relie la ressource (biomasse), les procédés biotechnologiques, les produits obtenus et le double regard avantages/limites que tu dois savoir mobiliser au Baccalauréat.

\begin{popfigure}
\begin{center}
\begin{center}\tcbox[colback=popink,colframe=popink,coltext=white,arc=9pt,boxsep=4pt,left=6pt,right=6pt]{\bfseries\small Biomasse et biocarburants}\end{center}
\vspace{-2pt}
\begin{tcbraster}[raster columns=2,raster equal height=rows,raster column skip=6pt,raster row skip=6pt,enhanced,blankest]
\begin{tcolorbox}[enhanced,colback=white,colframe=popgreen,boxrule=0.9pt,arc=3pt,title={Matières premières},fonttitle=\bfseries\scriptsize,coltitle=white,colbacktitle=popgreen,left=4pt,right=4pt,top=2pt,bottom=2pt,boxsep=1pt,toptitle=1.5pt,bottomtitle=1.5pt,halign=left]
\begin{itemize}[nosep,leftmargin=*,itemsep=1pt,font=\scriptsize]
\item Canne à sucre, manioc (sucres, amidon)
\item Huile de palme (lipides)
\item Déchets organiques, déjections
\end{itemize}
\end{tcolorbox}
\begin{tcolorbox}[enhanced,colback=white,colframe=poporange,boxrule=0.9pt,arc=3pt,title={Procédés},fonttitle=\bfseries\scriptsize,coltitle=white,colbacktitle=poporange,left=4pt,right=4pt,top=2pt,bottom=2pt,boxsep=1pt,toptitle=1.5pt,bottomtitle=1.5pt,halign=left]
\begin{itemize}[nosep,leftmargin=*,itemsep=1pt,font=\scriptsize]
\item Fermentation alcoolique (levures)
\item Transestérification des huiles
\item Méthanisation (anaérobie)
\end{itemize}
\end{tcolorbox}
\begin{tcolorbox}[enhanced,colback=white,colframe=popblue,boxrule=0.9pt,arc=3pt,title={Produits},fonttitle=\bfseries\scriptsize,coltitle=white,colbacktitle=popblue,left=4pt,right=4pt,top=2pt,bottom=2pt,boxsep=1pt,toptitle=1.5pt,bottomtitle=1.5pt,halign=left]
\begin{itemize}[nosep,leftmargin=*,itemsep=1pt,font=\scriptsize]
\item Bioéthanol
\item Biodiesel
\item Biogaz (\ce{CH4} + \ce{CO2})
\end{itemize}
\end{tcolorbox}
\begin{tcolorbox}[enhanced,colback=white,colframe=poppurple,boxrule=0.9pt,arc=3pt,title={Avantages et limites},fonttitle=\bfseries\scriptsize,coltitle=white,colbacktitle=poppurple,left=4pt,right=4pt,top=2pt,bottom=2pt,boxsep=1pt,toptitle=1.5pt,bottomtitle=1.5pt,halign=left]
\begin{itemize}[nosep,leftmargin=*,itemsep=1pt,font=\scriptsize]
\item Renouvelable, stockable, moins de particules
\item Valorisation des déchets, emplois ruraux
\item Compétition avec vivrier, faibles rendements
\item S'intègre dans un mix énergétique
\end{itemize}
\end{tcolorbox}
\end{tcbraster}

\end{center}
\legende{Carte mentale de synthèse de la leçon. De la biomasse (matières premières) aux biocarburants (produits) par les procédés biotechnologiques, avec le regard critique (avantages et limites) exigé dans toute argumentation au Baccalauréat.}
\end{popfigure}

\begin{aretenir}[L'essentiel de la section]
\begin{itemize}
\item Les biocarburants sont des énergies renouvelables \textbf{stockables}, complémentaires de l'hydroélectricité, du solaire et de l'éolien ; ils ne remplacent pas seuls les fossiles mais s'intègrent dans un \textbf{mix énergétique}.
\item La production \textbf{décentralisée} (digesteur familial, alcool à brûler local) réduit les délestages, les coûts des ménages et la dépendance aux importations : les déjections de 4 à 6 vaches couvrent les besoins de cuisson d'une famille.
\item Face à la pollution : moins de particules et de \ce{SO2}, capture du méthane des déchets, engrais organique en coproduit.
\item Le citoyen et l'élève agissent : tri et valorisation des déchets, promotion du biogaz, consommation responsable, et savent construire un message de sensibilisation (un problème, une solution, un chiffre clé, un appel à l'action).
\end{itemize}
\end{aretenir}

\begin{exercice}[Pour t'entraîner]
Ton quartier organise une journée « énergie propre ». (1) Choisis un public (élèves du primaire, mamans du marché, ou éleveurs) et rédige une affiche complète selon la méthode en quatre étapes. (2) Explique en cinq lignes pourquoi il serait imprudent de convertir toutes les terres vivrières du village en culture de canne pour le bioéthanol. (3) Calcule le nombre de vaches nécessaires pour alimenter les digesteurs d'un quartier de 30 familles (une famille = un digesteur = 4 à 6 vaches).
\end{exercice}

\begin{corrige}[Exercice « Pour t'entraîner »]
(1) Exemple pour les mamans du marché — problème : « La fumée du feu de bois fatigue vos yeux et fait tousser les enfants » ; solution : « Le biogaz des déchets du marché brûle sans fumée » ; chiffre clé : « 4 vaches = tout le gaz de cuisson d'une famille » ; appel à l'action : « Inscrivez-vous à la démonstration de digesteur samedi à la mairie ». (2) Convertir les terres vivrières créerait une compétition avec l'alimentation : hausse des prix des denrées, insécurité alimentaire, appauvrissement des sols ; il faut réserver les filières biocarburants aux excédents, aux résidus et aux terres non alimentaires. (3) Entre $30 \times 4 = 120$ et $30 \times 6 = 180$ vaches, soit un cheptel de \textbf{120 à 180 vaches} pour le quartier.
\end{corrige}

\newpage

\coursec{Activité d'intégration}

\coursec{Leçon 19 : Recherche de nouvelles sources d'énergie}

\courssub{Biomasse, biocarburants et contexte énergétique}

\begin{definition}[Biomasse et biocarburants]
On appelle \cle{biomasse} l'ensemble des matières d'origine biologique (végétale, animale, bactérienne, fongique) disponibles sur une base renouvelable. Les \cle{biocarburants} (ou agrocarburants) sont des carburants produits à partir de cette biomasse. On distingue trois générations :
\begin{itemize}
\item \textbf{1\textsuperscript{re} génération :} cultures alimentaires (sucre, amidon, huile) -- \emph{ex. : canne à sucre, maïs, manioc, palmier à huile}.
\item \textbf{2\textsuperscript{e} génération :} biomasses lignocellulosiques non alimentaires (résidus agricoles, bois, paille, déchets ménagers).
\item \textbf{3\textsuperscript{e} génération :} micro-algues et cyanobactéries (rendement surfacique élevé, pas de compétition foncière avec l'alimentation).
\end{itemize}
\end{definition}

\begin{savaistu}[Contexte camerounais]
Le Cameroun importe la totalité de ses produits pétroliers raffinés (SONARA ne raffine que le brut local, insuffisant). Le déficit énergétique se manifeste par des délestages électriques fréquents et un coût élevé du gasoil et de l'essence à la pompe. Le pays dispose pourtant d'atouts majeurs : climat tropical favorable à la biomasse (canne à sucre SOSUCAM, palmier à huile SOCAPAM/SPFS, manioc, maïs, jatropha), importante production de déchets agricoles (bagasse, rafles, écorces) et d'élevage (fumier). La \cle{Stratégie Nationale de Développement des Énergies Renouvelables (SNDER)} vise à porter la part des EnR dans le mix électrique à 25 \% à l'horizon 2030.
\end{savaistu}

\begin{propriete}[Principaux types de biocarburants]
\begin{itemize}
\item \textbf{Bioéthanol :} alcool (éthanol) issu de la fermentation alcoolique de sucres (C6H12O6) par des levures. Utilisé en mélange avec l'essence (E10, E85) ou dans moteurs flex-fuel.
\item \textbf{Biodiesel (EMHV -- Esters Méthyliques d'Huiles Végétales) :} issu de la \cle{transestérification} d'huiles végétales (palme, coton, jatropha, arachide) ou animales avec un alcool (méthanol) en présence d'un catalyseur (base forte type KOH/NaOH). Utilisé en mélange avec le gasoil (B7, B30, B100).
\item \textbf{Biogaz :} mélange gazeux (50--70 \% CH4, 30--50 \% CO2, traces H2S, H2O) produit par \cle{méthanisation} (digestion anaérobie) de matières organiques (fumier, lisier, déchets agro-alimentaires, boues de station d'épuration). Utilisable en cogénération (chaleur + électricité), comme carburant (bioGNV) ou injecté dans le réseau de gaz.
\end{itemize}
\end{propriete}

\courssub{La fermentation alcoolique : fondement expérimental du bioéthanol}

\begin{experience}[Mise en évidence de la fermentation alcoolique (Document 3 de l'activité)]
\textbf{Matériel :} 2 tubes à essai, solution de glucose 10 \%, levure de bière (Saccharomyces cerevisiae), eau de chaux, bains thermostatés à 35 °C et 5 °C, tuyaux de recueil de gaz.
\textbf{Protocole :} Introduire 10 mL de solution glucosée + levure dans chaque tube. Connecter à l'eau de chaux. Placer tube 1 à 35 °C, tube 2 à 5 °C. Observer après 30 min.
\textbf{Observations :}
\begin{center}
\begin{tabularx}{\linewidth}{|l|X|X|}\hline
\rowcolor{popblueL} \textbf{Tube} & \textbf{Dégagement gazeux} & \textbf{Eau de chaux} \\ \hline
Tube 1 (35 °C) & Abondant, bulles régulières & Trouble blanchâtre rapide (précipité de CaCO3) \\ \hline
Tube 2 (5 °C) & Très faible, quasi nul & Limpide (pas de trouble) \\ \hline
\end{tabularx}
\end{center}
\textbf{Interprétation :} Le trouble de l'eau de chaux (Ca(OH)2) prouve la production de \ce{CO2} (\ce{Ca(OH)2 + CO2 -> CaCO3 v + H2O}). L'odeur d'alcool et le contexte anaérobie confirment la formation d'éthanol. La forte différence de vitesse entre 35 °C (optimum enzymatique des levures) et 5 °C (ralentissement cinétique, enzymes inactives) démontre le caractère \textbf{enzymatique} de la réaction (catalysée par le complexe \cle{zymase}).
\end{experience}

\begin{propriete}[Équation bilan et rendement théorique]
La fermentation alcoolique est une voie métabolique anaérobie (glycolyse + réduction du pyruvate) assurant la régénération du \cle{NAD+} nécessaire à la poursuite de la glycolyse.
\[ \ce{C6H12O6 ->[zymase][\text{levure, anaérobiose, 30-35°C}] 2 C2H5OH + 2 CO2 + 2 ATP} \]
\textbf{Rendement massique théorique :} 180 g glucose $\rightarrow$ 92 g éthanol + 88 g CO2. Soit 0,511 g éthanol / g glucose.
\textbf{Rendement volumique pratique :} 1 t de céréales (amidon $\approx$ 65-70 \%) $\rightarrow$ 350-400 L éthanol (après hydrolyse amidon $\rightarrow$ glucose).
\end{propriete}

\begin{attention}[Pièges fréquents]
\begin{itemize}
\item Confondre \emph{fermentation} (anaérobie, 2 ATP/glucose) et \emph{respiration cellulaire} (aérobie, ~30-32 ATP/glucose).
\item Oublier que la levure a besoin d'éléments azotés, vitamines, minéraux (milieu complet) pour une fermentation industrielle efficace.
\item Croire que l'éthanol est le produit final unique : le \ce{CO2} est un co-produit valorisable (boissons gazeuses, industrie alimentaire, capture pour usage industriel).
\end{itemize}
\end{attention}

\courssub{Production du bioéthanol et matières premières camerounaises}

\begin{methode}[Filière bioéthanol 1\textsuperscript{re} génération (voie amidonnière : maïs, manioc, pomme de terre)]
\begin{enumerate}
\item \textbf{Préparation / Broyage :} réduction en farine/pulpe pour augmenter la surface de contact.
\item \textbf{Hydrolyse (liquéfaction + saccharification) :} cuisson sous pression (gélatinisation amidon) puis action d'enzymes (\cle{$\alpha$-amylase} puis \cle{glucoamylase}) pour hydrolyser l'amidon en glucose. \emph{Pour la canne à sucre (voie sucrine) : simple extraction du jus par broyage/pressage, pas d'hydrolyse nécessaire.}
\item \textbf{Fermentation :} inoculation de levures (S. cerevisiae), 30-35 °C, pH 4,5-5,5, anaérobiose stricte, 48-72 h. Titrage alcoolique final $\approx$ 8-12 \% vol.
\item \textbf{Distillation :} séparation éthanol/eau par distillation fractionnée (point d'ébullition éthanol 78,4 °C, eau 100 °C). On obtient de l'alcool vinique $\approx$ 92-96 \% vol.
\item \textbf{Déshydratation :} élimination de l'eau résiduelle (tamis moléculaires, azéotrope cyclopentane) pour obtenir de l'éthanol anhydre (> 99,5 \% vol) compatible essence.
\item \textbf{Valorisation des coproduits :} drêches (résidus solides riches en protéines) $\rightarrow$ alimentation animale (DDGS) ; vinasses (effluents liquides) $\rightarrow$ fertirrigation ou méthanisation ; \ce{CO2} $\rightarrow$ capture.
\end{enumerate}
\end{methode}

\begin{exemplebox}[Rendements au Cameroun (Document 1)]
\begin{center}
\begin{tabularx}{\linewidth}{|l|X|X|X|}\hline
\rowcolor{popblueL} \textbf{Culture} & \textbf{Rendement agricole (t/ha)} & \textbf{Part valorisable} & \textbf{Éthanol produit (L/ha)} \\ \hline
Canne à sucre (SOSUCAM) & 70 (tiges) & Jus sucré (saccharose) & \num{6000} \\ \hline
Maïs (grains) & 2,5 & Amidon (grains) & \num{2500} \\ \hline
Manioc (racines) & 12 & Amidon (racines) & \num{2160} \\ \hline
Pomme de terre (tubercules) & 15 & Amidon (tubercules) & \num{1500} \\ \hline
Palmier à huile & 20 (régimes) & Huile (biodiesel) & -- \\ \hline
Jatropha curcas & 3-5 (graines) & Huile (biodiesel, non alimentaire) & -- \\ \hline
\end{tabularx}
\end{center}
\emph{Source : données moyennes observées au Cameroun (MINADER, IRAD, SOSUCAM). La canne à sucre reste la plus productive en L/ha pour l'éthanol.}
\end{exemplebox}

\begin{aretenir}[Limites des agrocarburants 1\textsuperscript{re} génération]
La compétition \textbf{alimentation vs carburant} (surfaces arables, eau, engrais) est l'objection éthique et économique majeure. L'indice de \cle{retour énergétique (EROI -- Energy Return On Investment)} est souvent faible pour le maïs (1,3-1,6) mais favorable pour la canne à sucre (8-10) grâce à la valorisation énergétique de la bagasse (cogénération usine).
\end{aretenir}

\courssub{Biodiesel et biogaz : les autres biocarburants}

\begin{methode}[Production du biodiesel par transestérification]
\begin{enumerate}
\item \textbf{Prétraitement de l'huile :} dégommage, neutralisation (acides gras libres), séchage. Si taux d'acidité > 2 \%, estérification acide préalable.
\item \textbf{Transestérification basique :} Triglycéride $+ 3\ \ce{CH3OH} \xrightarrow{\text{KOH/NaOH, 60°C}} 3$ EMHV (biodiesel) $+$ Glycérine.
\item \textbf{Séparation :} décantation (glycérine dense en bas, biodiesel en haut).
\item \textbf{Lavage / Séchage :} eau (élimination savons, catalyseur, glycérine résiduelle) puis séchage sous vide.
\item \textbf{Contrôle qualité :} norme EN 14214 / ASTM D6751 (indice d'acidité, teneur en esters, glycérine libre/totale, point de trouble, teneur en soufre).
\end{enumerate}
\emph{Coproduit principal : glycérine brute (valorisation cosmétique, pharmaceutique, chimie verte après purification).}
\end{methode}

\begin{experience}[Méthanisation (production de biogaz) -- Principe du digesteur]
\textbf{Substrats :} fumier bovin/porcin, lisier, résidus de cultures, déchets ménagers organiques, boues d'épuration.
\textbf{Procédé :} digestion anaérobie en 4 étapes successives par consortium microbien :
\begin{enumerate}
\item \textbf{Hydrolyse :} bactéries hydrolytiques $\rightarrow$ polymères (protéines, lipides, glucides) $\rightarrow$ monomères (acides aminés, acides gras, sucres).
\item \textbf{Acidogenèse :} bactéries acidogènes $\rightarrow$ monomères $\rightarrow$ acides gras volatils (AGV), alcools, \ce{CO2}, \ce{H2}, \ce{NH3}.
\item \textbf{Acétogenèse :} bactéries acétogènes $\rightarrow$ AGV/alcools $\rightarrow$ acétate, \ce{H2}, \ce{CO2}.
\item \textbf{Méthanogenèse :} archées méthanogènes (strictement anaérobies) $\rightarrow$ 2 voies : \ce{CH3COO- + H+ -> CH4 + CO2} (acétoclastique, ~70 \%) et \ce{4 H2 + CO2 -> CH4 + 2 H2O} (hydrogénotrophe, ~30 \%).
\end{enumerate}
\textbf{Paramètres clés :} T° mésophile (35-38 °C) ou thermophile (50-55 °C), pH 6,8-7,5, C/N $\approx$ 20-30, temps de séjour hydraulique (TSH) 20-40 j, agitation lente, absence d'inhibiteurs (NH3 > 3 g/L, H2S, antibiotiques, métaux lourds).
\end{experience}

\begin{propriete}[Caractéristiques du biogaz et du digestat]
\begin{itemize}
\item \textbf{Biogaz :} 50-70 \% \ce{CH4}, PCI $\approx$ 6-7 kWh/m3 (vs 10 kWh/m3 gaz naturel). 1 m3 biogaz $\approx$ 0,6 L gasoil en énergie (donnée programme).
\item \textbf{Digestat :} résidu liquide/solide riche en \ce{NH4+} (azote minéral directement assimilable), \ce{PO4^3-}, \ce{K+}. Excellent \cle{engrais organique} (substitution engrais chimiques), amendement humique. Attention : hygénisation requise si pathogènes (70 °C, 1 h).
\end{itemize}
\end{propriete}

\begin{savaistu}[Biogaz au Cameroun : expériences locales]
Des digesteurs à dôme fixe (type chinois) ou à ballon (type indien) ont été installés dans des centres de formation agricole (ex. CFP de Mbalmayo, fermes pilotes à l'Ouest) et quelques ménages ruraux. Le programme \emph{Programme National Biogaz (PNB)} appuyé par la GIZ/SNV a démontré la faisabilité technique. Les freins : coût d'investissement initial (500 000 - 2 000 000 FCFA selon volume), disponibilité en eau (1:1 fumier/eau), gestion du digestat, maintenance (désulfuration, condensation).
\end{savaistu}

\courssub{Biocarburants vs énergies fossiles : avantages et limites}

\begin{aretenir}[Bilan comparatif multicritère]
\begin{center}
\begin{tabularx}{\linewidth}{|l|X|X|}\hline
\rowcolor{popblueL} \textbf{Critère} & \textbf{Biocarburants (1\textsuperscript{re} gén.)} & \textbf{Énergies fossiles (pétrole)} \\ \hline
\textbf{Renouvelabilité} & Renouvelable (cycle court C, 1 an) & Non renouvelable (stocks finis, cycle géologique) \\ \hline
\textbf{Bilan carbone (cycle de vie)} & Quasi neutre \emph{si} pas de changement d'affectation des sols (CAS/ILUC). Émissions : culture, engrais (N2O), transformation, transport. & Fortement émetteur (C fossile $\rightarrow$ atmosphère). \\ \hline
\textbf{Pollution locale} & Moins de SOx, particules, CO, HAP. Mais NOx similaires ou supérieurs (biodiesel). & SOx, NOx, CO, HAP, particules fines (PM2,5), benzène. \\ \hline
\textbf{Sécurité énergétique} & Production locale, décentralisée possible. Réduit importations. & Dépendance importations, volatilité prix, géopolitique. \\ \hline
\textbf{Compétition alimentaire} & \textbf{Majeure} (terres, eau, engrais). Risque hausse prix denrées. & Aucune (sauf raffinage). \\ \hline
\textbf{Coût de production} & Généralement > pétrole (subventions nécessaires). Variable selon matière 1\textsuperscript{re}. & Historiquement bas, mais volatile. \\ \hline
\textbf{Adaptation moteurs} & Éthanol : matériaux compatibles (joints, réservoir), injection adaptée. Biodiesel : lubrifiant, nettoyant (risque colmatage filtres au début). & Standard. \\ \hline
\end{tabularx}
\end{center}
\end{aretenir}

\begin{attention}[Analyse du cycle de vie (ACV) et ILUC]
Le "zéro carbone" à la combustion est un leurre. L'\cle{ACV} (Well-to-Wheel) inclut : fabrication engrais (N2O, puissant GES), travail du sol (perte C organique), transformation (énergie process), transport. Le \cle{changement d'affectation des sols indirect (ILUC)} : conversion forêts/pâturages $\rightarrow$ cultures énergétiques $\rightarrow$ émission massive C stocké + perte biodiversité. \textbf{Directive européenne RED II :} plafonnement agrocarburants 1\textsuperscript{re} gén. à 7 \% conso transport, promotion 2\textsuperscript{e}/3\textsuperscript{e} gén.
\end{attention}

\courssub{Énergies renouvelables, éducation et citoyenneté}

\begin{methode}[Démarche d'éducation et de sensibilisation aux EnR]
\begin{enumerate}
\item \textbf{Diagnostic local :} identifier les ressources (biomasse, soleil, vent, hydro), les besoins (électricité, cuisson, transport, engrais), les acteurs (collectivités, ONG, privé, État).
\item \textbf{Analyse multicritère :} technique (ressource, technologie), économique (CAPEX, OPEX, revenu, subventions), environnemental (ACV, biodiversité, eau, sols), social (emploi, santé, genre, acceptabilité, sécurité alimentaire).
\item \textbf{Choix argumenté :} transparence sur avantages \textbf{ET} limites (maintenance, formation, durée de vie, dépendance intrants).
\item \textbf{Communication adaptée :} langage accessible, supports visuels (schémas, photos), réunions participatives, langues locales, leaders d'opinion (chefs, relais communautaires, radios locales).
\item \textbf{Suivi-évaluation :} indicateurs techniques (production, taux de panne), économiques (économies, revenus), sociaux (adoption, satisfaction), environnementaux (bois économisé, engrais chimique remplacé).
\end{enumerate}
\end{methode}

\courssub{Activité d'intégration : Un biocarburant pour le village de Nkoné}

\courssub{Objectifs de l'activité}
À l'issue de cette activité d'intégration, tu seras capable de :
\begin{itemize}
\item interpréter les résultats d'une expérience de fermentation alcoolique et d'en écrire l'équation bilan ;
\item exploiter des documents (tableaux, courbes) pour dimensionner une installation de production de biocarburant ;
\item comparer deux filières de biocarburants (bioéthanol et biogaz) selon des critères économiques, écologiques et sociaux ;
\item rédiger une note de sensibilisation argumentée à destination d'une communauté, en vue d'éduquer et d'informer sur l'utilisation des énergies renouvelables.
\end{itemize}

\courssub{Situation}
\textbf{Un biocarburant pour le village de Nkoné.}

Le village de Nkoné (fictif), situé dans la région du Nord-Ouest du Cameroun, compte \num{2400} habitants répartis en 400 ménages. Le village n'est pas raccordé au réseau électrique national : l'éclairage du centre de santé, de l'école et de la place publique est assuré par un groupe électrogène communal fonctionnant au gasoil. En raison des délestages fréquents et du coût du carburant, le groupe ne fonctionne que 4 heures par jour. Le budget carburant de la communauté s'élève à \num{1200000} FCFA par an pour une consommation de \num{1800} litres de gasoil.

L'activité agricole repose sur la culture du maïs (rendement moyen : 2,5 tonnes de grains par hectare) et de la pomme de terre. L'élevage bovin occupe 60 familles qui possèdent au total 300 têtes de bétail ; chaque bovin produit en moyenne \SI{10}{kg} de fumier frais par jour, dont seulement la moitié est collectable (les animaux paissant une partie de la journée).

Le conseil municipal hésite entre deux projets :
\begin{itemize}
\item \textbf{Option A} : produire du bioéthanol à partir du maïs pour alimenter un moteur adapté ;
\item \textbf{Option B} : produire du biogaz à partir du fumier collecté dans un digesteur.
\end{itemize}

\textbf{Document 1 : rendements en éthanol de quelques cultures}

\begin{center}
\begin{tabularx}{\linewidth}{|l|X|X|}\hline
\rowcolor{popblueL} \textbf{Culture} & \textbf{Rendement agricole (t/ha)} & \textbf{Éthanol produit (L/ha)} \\ \hline
Canne à sucre & 70 & \num{6000} \\ \hline
Maïs & 2,5 & \num{2500} \\ \hline
Manioc & 12 & \num{2160} \\ \hline
Pomme de terre & 15 & \num{1500} \\ \hline
\end{tabularx}
\end{center}

\textbf{Document 2 : production de biogaz en fonction de la quantité de fumier introduite quotidiennement dans le digesteur}

\begin{popfigure}
\begin{center}
\begin{tikzpicture}
\begin{axis}[
width=0.85\linewidth, height=6cm,
xlabel={Fumier introduit (kg/jour)},
ylabel={Biogaz produit (m\textsuperscript{3}/jour)},
xmin=0, xmax=2000, ymin=0, ymax=45,
grid=both, grid style={popline!40},
xtick={0,500,1000,1500,2000},
ytick={0,10,20,30,40},
legend style={at={(0.03,0.97)}, anchor=north west, font=\small}
]
\addplot[popcurve, domain=0:2000, samples=100] {30*(1-exp(-x/600))};
\addlegendentry{Production de biogaz}
\addplot[poporange, dashed, domain=0:2000, samples=2] {30};
\addlegendentry{Saturation $\approx$ 30 m\textsuperscript{3}/jour}
\draw[popgreen, dashed] (axis cs:1500,0) -- (axis cs:1500,27.6) -- (axis cs:0,27.6);
\end{axis}
\end{tikzpicture}
\end{center}
\legende{Courbe de production de biogaz d'un digesteur villageois en fonction de la charge quotidienne en fumier. En pointillés verts : lecture pour \SI{1500}{kg/jour} de fumier, soit environ \SI{27,6}{m^3/jour} de biogaz.}
\end{popfigure}

\textbf{Document 3 : expérience de fermentation alcoolique réalisée au laboratoire du lycée}

Deux tubes à essai contiennent chacun \SI{10}{mL} d'une solution de glucose à 10 \% et de la levure de bière. Le tube 1 est placé à \SI{35}{\celsius}, le tube 2 à \SI{5}{\celsius}. Le gaz dégagé est recueilli et barbote dans de l'eau de chaux. Résultats après 30 minutes :

\begin{center}
\begin{tabularx}{\linewidth}{|l|X|X|}\hline
\rowcolor{popblueL} \textbf{Tube} & \textbf{Dégagement gazeux} & \textbf{Eau de chaux} \\ \hline
Tube 1 (\SI{35}{\celsius}) & Abondant & Se trouble rapidement \\ \hline
Tube 2 (\SI{5}{\celsius}) & Très faible & Reste limpide \\ \hline
\end{tabularx}
\end{center}

\textbf{Données complémentaires :} 1 L d'éthanol fournit environ 60 \% de l'énergie d'1 L de gasoil ; \SI{1}{m^3} de biogaz fournit une énergie comparable à 0,6 L de gasoil ; le digestat (résidu du digesteur) est un excellent engrais organique ; le maïs est la principale nourriture des habitants de Nkoné.

\courssub{Consignes}
\begin{enumerate}
\item À partir du document 3, dégage le problème scientifique posé, formule une hypothèse, interprète les résultats et écris l'équation bilan de la fermentation alcoolique.
\item En exploitant le document 2, détermine la quantité de fumier collectable chaque jour à Nkoné, puis estime la production quotidienne de biogaz et l'économie de gasoil correspondante sur une année.
\item À l'aide du document 1, calcule la surface de maïs nécessaire pour produire l'équivalent énergétique des \num{1800} litres de gasoil annuels, puis compare les deux options selon trois critères : économique, écologique et social.
\item Rédige une note de sensibilisation de 10 à 15 lignes à l'attention du chef du village, justifiant le choix de ton groupe et en précisant honnêtement les limites du projet retenu.
\end{enumerate}

\courssub{Ressources à mobiliser}
\begin{itemize}
\item L'équation bilan de la fermentation alcoolique : \ce{C6H12O6 -> 2 C2H5OH + 2 CO2} (dégagement de \ce{CO2} mis en évidence par le trouble de l'eau de chaux).
\item La fermentation est une réaction enzymatique : sa vitesse dépend de la température (optimum vers \SI{35}{\celsius}).
\item Lecture graphique : repérage d'un point sur une courbe, notion de saturation (palier).
\item Conversions énergétiques : 1 L éthanol $\approx$ 0,6 L gasoil ; \SI{1}{m^3} biogaz $\approx$ 0,6 L gasoil.
\item Démarche comparative multicritère : avantages et limites des biocarburants face aux énergies fossiles (renouvelabilité, bilan carbone, compétition avec l'alimentation, coûts).
\end{itemize}

\courssub{Démarche et corrigé commenté}

\begin{exoresolu}[Un biocarburant pour Nkoné : corrigé commenté]
Énoncé : répondre aux quatre consignes de l'activité à partir des documents 1, 2 et 3 et des données complémentaires.
\tcblower
\textbf{Étape 1 : interprétation de l'expérience de fermentation.}

\emph{Problème} : le glucose peut-il être transformé par les levures en l'absence de dioxygène, et quels produits cette transformation dégage-t-elle ?

\emph{Hypothèse} : les levures dégradent le glucose en produisant un gaz carbonique et de l'éthanol.

\emph{Interprétation} : le trouble de l'eau de chaux dans le tube 1 révèle la présence de dioxyde de carbone \ce{CO2} (formation d'un précipité blanc de carbonate de calcium \ce{CaCO3}). Le dégagement gazeux abondant à \SI{35}{\celsius} mais très faible à \SI{5}{\celsius} montre que la fermentation est une réaction catalysée par des enzymes (complexe zymase), dont l'activité est optimale aux températures modérées (30-35 °C) et fortement ralentie au froid. Le milieu dégage en outre une odeur caractéristique d'alcool : de l'éthanol s'est formé.

\emph{Conclusion et équation bilan} :
\[ \ce{C6H12O6 ->[levure][\text{anaérobiose, 35°C}] 2 C2H5OH + 2 CO2} \]
Une molécule de glucose donne deux molécules d'éthanol et deux molécules de dioxyde de carbone.

\textbf{Étape 2 : dimensionnement du digesteur (option B).}

Fumier collectable par jour :
\[ 300 \text{ bovins} \times \SI{10}{kg/jour} \times 0{,}5 = \SI{1500}{kg/jour}. \]
Lecture graphique sur le document 2 : pour \SI{1500}{kg/jour}, la production est d'environ \SI{27,6}{m^3/jour} de biogaz (on est proche du palier de saturation à 30 m3/jour : augmenter la charge au-delà n'augmenterait presque plus la production).

Équivalent gasoil annuel :
\begin{align*}
\text{Biogaz annuel} &= 27{,}6 \times 365 \approx \SI{10074}{m^3/an} \\
\text{Équivalent gasoil} &= 10074 \times 0{,}6 \approx \SI{6044}{L}
\end{align*}
Le biogaz couvrirait donc largement les \num{1800} litres de gasoil annuels du groupe électrogène (après adaptation du moteur au biogaz ou dual-fuel), avec un excédent utilisable pour la cuisson des ménages. L'économie directe sur le budget carburant est de l'ordre de \num{1200000} FCFA par an, et le digestat remplace partiellement les engrais chimiques pour les cultures.

\textbf{Étape 3 : comparaison des deux options.}

Besoin énergétique en équivalent éthanol :
L'éthanol ne fournissant que 60 \% de l'énergie du gasoil, il faut $1800 / 0{,}6 = \num{3000}$ litres d'éthanol pour remplacer \num{1800} L de gasoil.
\emph{(Piège à éviter : ne pas calculer $1800 \times 0,6 = 1080$ L, ce qui correspondrait à l'énergie de 1080 L d'éthanol exprimée en équivalent gasoil, et non l'inverse).}

Surface de maïs nécessaire (document 1 : \num{2500} L/ha) :
\[ \frac{3000}{2500} = 1{,}2 \text{ hectare}, \]
soit 3 tonnes de grains de maïs soustraites chaque année à l'alimentation du village (rendement 2,5 t/ha).

\begin{center}
\begin{tabularx}{\linewidth}{|l|X|X|}\hline
\rowcolor{popblueL} \textbf{Critère} & \textbf{Option A : bioéthanol (maïs)} & \textbf{Option B : biogaz (fumier)} \\ \hline
Économique & Nécessite 1,2 ha de maïs/an, une unité de fermentation et de distillation coûteuse (cuves, chauffage, colonne distillation, énergie), main-d'œuvre qualifiée. & Matière première gratuite (déchet), digesteur simple à construire (béton, briques), faible coût opératoire, économie de \num{1200000} FCFA/an + valeur engrais digestat. \\ \hline
Écologique & Renouvelable mais mobilise terres arables, eau, engrais azotés (émissions N2O), pesticides. Risque érosion. & Valorise un déchet polluant (fumier), réduit le méthane émis naturellement à l'air libre (GES puissant), produit un engrais organique (substitution engrais de synthèse), boucle C/N locale. \\ \hline
Social & \textbf{Compétition directe avec l'alimentation} (le maïs est la nourriture de base). Risque de hausse prix local, insécurité alimentaire. & Aucune compétition alimentaire ; améliore l'hygiène (gestion déjections) ; cuisson propre (moins de fumée intérieure) pour les ménages ; création emplois locaux (collecte, maintenance). \\ \hline
\end{tabularx}
\end{center}

\textbf{Étape 4 : note de sensibilisation (exemple de production attendue).}

\emph{« Au Chef du village de Nkoné. Notre groupe a étudié deux solutions pour réduire la facture de gasoil du village. Nous recommandons le biogaz : le fumier de nos 300 bovins, aujourd'hui perdu, peut produire chaque jour près de \SI{28}{m^3} de gaz, de quoi remplacer tout le gasoil du groupe électrogène et offrir du gaz de cuisson aux familles. Le résidu de digestion est un engrais gratuit pour nos champs. Le projet coûte peu, ne touche pas à nos récoltes de maïs et assainit le village. Ses limites : il faut construire le digesteur, adapter le moteur, collecter le fumier chaque jour et entretenir l'installation (désulfuration, étanchéité) ; la production baisse si le troupeau diminue. Le bioéthanol à base de maïs, lui, priverait le village de nourriture et coûterait plus cher. Nous demandons l'accord du conseil pour lancer le projet biogaz. »}
\end{exoresolu}

\courssub{Bilan de la leçon}

Cette leçon a permis de mobiliser l'ensemble des savoir-faire du module dans une situation authentique de déficit énergétique. Elle a d'abord montré que la fermentation alcoolique est un phénomène biologique vérifiable expérimentalement : le trouble de l'eau de chaux atteste la production de \ce{CO2}, et la comparaison des températures révèle le caractère enzymatique de la réaction (zymase). Elle a ensuite illustré l'exploitation quantitative de documents : lecture d'une courbe de production de biogaz avec la notion de saturation du digesteur, calculs de conversion énergétique et dimensionnement d'une installation à partir de données réelles de cheptel. La filière biodiesel (transestérification huiles végétales) a été présentée comme alternative pour les zones oléagineuses (palmier, jatropha). Enfin, la comparaison multicritère a mis en évidence une leçon essentielle de citoyenneté scientifique : le « meilleur » biocarburant n'est pas celui qui a le meilleur rendement théorique, mais celui qui s'adapte au contexte local. À Nkoné, le biogaz l'emporte car il valorise un déchet gratuit, évite la compétition avec l'alimentation humaine — limite majeure des agrocarburants de première génération — et apporte des bénéfices sanitaires et agronomiques annexes. Tu es désormais capable d'éduquer et d'informer ta communauté sur les énergies renouvelables, en argumentant avec rigueur scientifique et en reconnaissant honnêtement les limites de chaque solution.

\newpage

\coursec{Méthodes et exercices résolus}

\coursec{Leçon 19 : Recherche de nouvelles sources d'énergie}

\courssub{Biomasse, biocarburants et contexte énergétique}

\begin{definition}[Biomasse et biocarburants]
On appelle \cle{biomasse} l'ensemble des matières organiques d'origine végétale (algues, plantes, résidus agricoles, forestiers) ou animale (déjections, effluents d'élevage) pouvant devenir sources d'énergie. Un \cle{biocarburant} est un carburant produit à partir de cette biomasse par des procédés biologiques (fermentation, méthanisation) ou chimiques (transestérification). Les trois filières principales sont le \textbf{bioéthanol}, le \textbf{biodiesel} et le \cle{biogaz}.
\end{definition}

\begin{savaistu}[Contexte énergétique camerounais]
Le Cameroun importe une grande partie de ses produits pétroliers raffinés (super, gasoil, kérosène) malgré sa production de brut. Les déficits énergétiques se manifestent par des délestages électriques récurrents et le coût élevé du carburant pour le transport (moto-taxis, taxis, poids lourds). La valorisation de la biomasse locale (canne à sucre SOSUCAM, huile de palme SCDP/SOCAPALM, manioc, déchets urbains de Yaoundé et Douala, résidus de cacao/café) est un levier stratégique pour la \cle{souveraineté énergétique} et la réduction de la facture d'importation, tout en luttant contre la déforestation liée au bois de chauffe.
\end{savaistu}

\begin{aretenir}[Les trois grandes filières de biocarburants]
\begin{center}
\begin{tabularx}{\linewidth}{|l|X|X|X|}
\hline
\rowcolor{popblueL} Filière & Biocarburant & Matières premières camerounaises & Procédé clé \\
\hline
Alcoolique & Bioéthanol & Canne à sucre, manioc, maïs, patate douce & Fermentation alcoolique (levures) + Distillation \\
\hline
Estérification & Biodiesel & Huile de palme, coton, jatropha, ricin & Transestérification (méthanol + soude) \\
\hline
Méthanisation & Biogaz & Fumier, résidus de récolte, ordures ménagères, boues de stations d'épuration & Digestion anaérobie (consortium bactérien) \\
\hline
\end{tabularx}
\end{center}
\end{aretenir}

\courssub{La fermentation alcoolique : fondement expérimental du bioéthanol}

\begin{definition}[Fermentation alcoolique]
La \cle{fermentation alcoolique} est une voie métabolique anaérobie (sans $\mathrm{O_2}$) réalisée par des levures (\emph{Saccharomyces cerevisiae}) qui dégrade incomplètement le glucose en éthanol et dioxyde de carbone, avec un faible rendement énergétique (2 ATP par glucose).
\[ \ce{C6H12O6 -> 2 C2H5OH + 2 CO2} \quad (\text{+ 2 ATP}) \]
\end{definition}

\begin{methode}[Analyser un dispositif expérimental de fermentation]
Face à un schéma ou un compte rendu de TP sur la fermentation :
\begin{enumerate}
\item \textbf{Repère les conditions} : milieu glucidique + levures, température optimale (25 à \SI{35}{\celsius}), dispositif clos (tube à dégagement ou bouchon percé) garantissant l'\textbf{anaérobiose}.
\item \textbf{Identifie les tests de mise en évidence} :
\begin{itemize}
\item $\mathrm{CO_2}$ : trouble l'\cle{eau de chaux} (formation d'un précipité blanc de $\mathrm{CaCO_3}$) ;
\item Éthanol : odeur caractéristique, test d'inflammabilité, dosage chromatique ;
\item Disparition du glucose : test à la \emph{liqueur de Fehling} (précipité rouge brique $\mathrm{Cu_2O}$ à chaud) qui devient négatif.
\end{itemize}
\item \textbf{Exploite les courbes cinétiques} : une courbe de dégagement de $\mathrm{CO_2} = f(t)$ montre une phase latente, une phase exponentielle (vitesse max), un ralentissement et un plateau.
\item \textbf{Conclus rigoureusement} : la fermentation est une dégradation incomplète du glucose, sans oxygène, produisant éthanol, $\mathrm{CO_2}$ et 2 ATP (contre $\approx$ 36 ATP en respiration aérobie).
\end{enumerate}
\end{methode}

\begin{popfigure}
\begin{center}
\begin{tikzpicture}[node distance=1.2cm, >=Stealth]
% Styles
\tikzset{
    vessel/.style={draw=popdark, thick, fill=popblueL!20, rounded corners, minimum width=2.5cm, minimum height=3cm},
    tube/.style={draw=popdark, thick, fill=white},
    label/.style={font=\small, text=popdark},
    arrow/.style={->, thick, poporange}
}
% Erlenmeyer
\node[vessel] (erlen) {};
\node[label, above=0.2cm of erlen.north] {Milieu glucidé + Levures};
\node[label, left=0.2cm of erlen.west, align=center] {$30\,^\circ\mathrm{C}$\protect\\Anaérobiose};
% Tube sortie
\draw[tube] (erlen.north east) -- ++(0.5,0) |- ++(1,1.5) node[above, label] {$\mathrm{CO_2}$};
% Eau de chaux
\node[vessel, right=2.5cm of erlen] (chaux) {};
\node[label, above=0.2cm of chaux.north] {Eau de chaux $\mathrm{Ca(OH)_2}$};
\draw[arrow] (erlen.north east) -- ++(0.5,0) -| (chaux.north west) node[midway, above, label] {Gaz};
% Flèche trouble
\draw[arrow] (chaux.south) -- ++(0,-0.8) node[below, label, align=center]{Trouble (blanc)\protect\\$\mathrm{CaCO_3}\downarrow$};
% Légendes internes
\node[label, fill=popcream, rounded corners, inner sep=2pt] at (erlen.center) {Glucose $\xrightarrow{\text{levures}}$ $\mathrm{CO_2}$ + Éthanol};
\end{tikzpicture}
\end{center}
\legende{Schéma du dispositif classique de mise en évidence de la fermentation alcoolique : dégagement de $\mathrm{CO_2}$ trouble l'eau de chaux.}
\end{popfigure}

\begin{exoresolu}[Suivi cinétique d'une fermentation en TP]
Au laboratoire du lycée, on introduit $\SI{20}{g}$ de glucose et $\SI{2}{g}$ de levure de boulanger dans $\SI{500}{mL}$ d'eau à $\SI{30}{\celsius}$. On mesure le volume de $\mathrm{CO_2}$ dégagé :

\begin{center}
\begin{tabular}{|l|c|c|c|c|c|c|}
\hline
\rowcolor{popblueL} Temps (h) & 0 & 2 & 4 & 6 & 8 & 10 \\
\hline Volume de $\mathrm{CO_2}$ (mL) & 0 & 15 & 90 & 210 & 260 & 265 \\
\hline
\end{tabular}
\end{center}

1. Trace l'allure de la courbe $V_{\mathrm{CO_2}} = f(t)$ et décris ses phases.\\
2. Pourquoi le dégagement de $\mathrm{CO_2}$ s'arrête-t-il ? Propose deux hypothèses vérifiables.\\
3. Calcule la vitesse moyenne de dégagement de $\mathrm{CO_2}$ entre $t = \SI{2}{h}$ et $t = \SI{6}{h}$.
\tcblower
\textbf{1. Allure de la courbe et phases}

\begin{popfigure}
\begin{center}
\begin{tikzpicture}
\begin{axis}[width=0.85\linewidth, height=6cm,
xlabel={Temps (h)}, ylabel={Volume de $\mathrm{CO_2}$ (mL)},
xmin=0, xmax=11, ymin=0, ymax=300,
xtick={0,2,4,6,8,10}, ytick={0,50,100,150,200,250,300},
grid=both, grid style={popline!40},
legend pos=north west]
\addplot[popcurve, mark=*, mark size=2pt] coordinates {(0,0) (2,15) (4,90) (6,210) (8,260) (10,265)};
\addlegendentry{Expérimental}
\end{axis}
\end{tikzpicture}
\end{center}
\legende{Courbe de dégagement de $\mathrm{CO_2}$ : phase latente (0--2 h), phase active (2--6 h), plateau (8--10 h).}
\end{popfigure}

On distingue trois phases : une \textbf{phase latente} (0 à 2 h) où les levures s'adaptent et se multiplient ; une \textbf{phase de fermentation active} (2 à 6 h) où la pente est maximale (vitesse max) ; un \textbf{ralentissement puis plateau} (6 à 10 h) où la production cesse.

\textbf{2. Hypothèses explicatives du plateau}

\begin{itemize}
\item \emph{Hypothèse 1 (Épuisement du substrat)} : le glucose initial ($\SI{20}{g}$ dans $\SI{500}{mL}$, soit $\SI{40}{g.L^{-1}}$) a été entièrement consommé. \textbf{Vérification} : test à la liqueur de Fehling sur le milieu final ; s'il est négatif (absence de précipité rouge brique), le glucose a disparu.
\item \emph{Hypothèse 2 (Inhibition par l'éthanol)} : l'accumulation d'éthanol devient toxique pour les levures (au-delà de $\approx$ 12--15 \% vol.). \textbf{Vérification} : dosage de l'éthanol (densimétrie ou chromatographie).
\end{itemize}
Calcul théorique : $\SI{20}{g}$ glucose $\rightarrow$ $\approx \SI{9,2}{g}$ éthanol (rendement 90 \%) dans $\SI{0,5}{L}$ $\rightarrow$ $\approx \SI{1,8}{\%}$ vol. Bien en dessous du seuil toxique. \textbf{L'épuisement du glucose est l'hypothèse la plus vraisemblable.}

\textbf{3. Vitesse moyenne entre 2 h et 6 h}
\[ v = \dfrac{\Delta V}{\Delta t} = \dfrac{210 - 15}{6 - 2} = \dfrac{195}{4} = \SI{48,75}{mL.h^{-1}} \approx \SI{48,8}{mL.h^{-1}} \]

\textbf{Conclusion} : le dégagement de $\mathrm{CO_2}$, témoin direct de l'activité fermentaire, atteint un maximum de $\approx \SI{48,8}{mL.h^{-1}}$ entre 2 h et 6 h, puis s'annule, traduisant l'arrêt de la fermentation par épuisement du substrat glucidique.
\end{exoresolu}

\courssub{Production du bioéthanol et matières premières camerounaises}

\begin{methode}[Décrire la filière de production du bioéthanol]
Pour expliquer la production d'un biocarburant à partir de la biomasse, suis la démarche en 5 étapes :
\begin{enumerate}
\item \textbf{Identifier la matière première (biomasse)} : plante sucrière (canne à sucre $\rightarrow$ \cle{saccharose} $\mathrm{C_{12}H_{22}O_{11}}$), plante amylacée (manioc, maïs $\rightarrow$ \cle{amidon}) ou déchets organiques.
\item \textbf{Décrire la préparation du substrat} :
\begin{itemize}
\item Canne à sucre : broyage + extraction du jus (vesou) riche en saccharose.
\item Plantes amylacées (manioc) : cuisson + \cle{hydrolyse} enzymatique de l'amidon en glucose par des \emph{amylases} (maltage ou enzymes industrielles).
\end{itemize}
\item \textbf{Écrire la réaction clé} : la \cle{fermentation alcoolique} par \emph{Saccharomyces cerevisiae} en \textbf{anaérobiose} :
\[ \ce{C6H12O6 -> 2 C2H5OH + 2 CO2} \quad (\text{+ 2 ATP}) \]
\item \textbf{Décrire la purification} : distillation du moût fermenté (titre alcoométrique $\le$ 15 \% vol., limite de tolérance des levures) pour obtenir de l'éthanol concentré (96 \% vol. par distillation simple, 99,5 \% par déshydratation pour carburant).
\item \textbf{Conclure sur l'utilisation} : éthanol carburant pur ou mélangé à l'essence (E10 = 10 \% éthanol, E85 = 85 \% éthanol).
\end{enumerate}
\textbf{Réflexe rédaction} : cite toujours l'organisme (levure), la condition (absence de $\mathrm{O_2}$) et l'équation bilan. C'est le trio qui rapporte tous les points au Bac.
\end{methode}

\begin{exoresolu}[Production d'éthanol à partir de la canne à sucre (SOSUCAM)]
Une sucrerie camerounaise traite du jus de canne contenant du saccharose. Le saccharose est hydrolysé en glucose et fructose (inversion), puis le glucose est fermenté par des levures. On traite une cuve contenant $\SI{360}{kg}$ de glucose.\\
1. Écris l'équation de la fermentation alcoolique du glucose.\\
2. Calcule la masse d'éthanol théorique obtenue. On donne $M(\mathrm{C_6H_{12}O_6}) = \SI{180}{g.mol^{-1}}$ et $M(\mathrm{C_2H_5OH}) = \SI{46}{g.mol^{-1}}$.\\
3. Le rendement réel de la fermentation est de 90 \%. Calcule la masse d'éthanol réellement produite, puis son volume (masse volumique de l'éthanol : $\rho = \SI{0,79}{kg.L^{-1}}$).
\tcblower
\textbf{1. Équation de la fermentation alcoolique}

En anaérobiose, les levures dégradent le glucose en éthanol et dioxyde de carbone :
\[ \ce{C6H12O6 -> 2 C2H5OH + 2 CO2} \]
Vérification : 6 C, 12 H, 6 O de chaque côté. L'équation est équilibrée.

\textbf{2. Masse théorique d'éthanol}

Quantité de matière de glucose :
\[ n(\text{glucose}) = \dfrac{m}{M} = \dfrac{\SI{360000}{g}}{\SI{180}{g.mol^{-1}}} = \SI{2000}{mol} \]
D'après l'équation, 1 mol glucose $\rightarrow$ 2 mol éthanol :
\[ n(\text{éthanol}) = 2 \times 2000 = \SI{4000}{mol} \]
Masse théorique :
\[ m_{\text{th}} = n \times M = 4000 \times 46 = \SI{184000}{g} = \SI{184}{kg} \]

\textbf{3. Masse et volume réels (rendement 90 \%)}

\[ m_{\text{réel}} = m_{\text{th}} \times 0,90 = 184 \times 0,90 = \SI{165,6}{kg} \]
\[ V = \dfrac{m}{\rho} = \dfrac{165,6}{0,79} \approx \SI{209,6}{L} \]

\textbf{Conclusion} : la fermentation de $\SI{360}{kg}$ de glucose fournit théoriquement $\SI{184}{kg}$ d'éthanol ; avec un rendement de 90 \%, on obtient $\SI{165,6}{kg}$, soit environ $\SI{210}{L}$ de bioéthanol carburant.
\end{exoresolu}

\begin{methode}[Calculer un rendement ou un bilan énergétique de filière]
\begin{enumerate}
\item \textbf{Identifie les données} : masse de biomasse, rendement de conversion (massique ou volumique), pouvoir calorifique inférieur (PCI) du produit final (en $\SI{}{kWh.kg^{-1}}$ ou $\SI{}{MJ.kg^{-1}}$).
\item \textbf{Applique les relations de base} :
\[ m_{\text{produit}} = m_{\text{biomasse}} \times \text{rendement massique} \qquad ; \qquad E = m_{\text{produit}} \times \text{PCI} \]
\item \textbf{Respecte les unités} : convertis les tonnes en kilogrammes ($\SI{1}{t} = \SI{1000}{kg}$), les pourcentages en fractions (diviser par 100).
\item \textbf{Interprète le résultat} : compare l'énergie obtenue à un besoin concret (consommation d'un foyer, autonomie d'un véhicule) pour conclure sur l'intérêt de la filière.
\end{enumerate}
\textbf{Réflexe Bac} : écris systématiquement la formule littérale \emph{avant} l'application numérique, et encadre le résultat final avec son unité.
\end{methode}

\begin{exoresolu}[Bilan énergétique d'une unité de bioéthanol à base de canne (SOSUCAM)]
Un hectare de canne à sucre donne $\SI{80}{t}$ de canne. Une tonne de canne fournit environ $\SI{70}{kg}$ de sucres fermentescibles. La fermentation + distillation convertit ce sucre en éthanol avec un rendement massique global de 50 \% (1 kg sucre $\rightarrow$ 0,5 kg éthanol). Le PCI de l'éthanol est de $\SI{7,4}{kWh.kg^{-1}}$.\\
1. Calcule la masse d'éthanol produite par hectare.\\
2. Calcule l'énergie contenue dans cet éthanol.\\
3. Une moto-taxi consomme en moyenne $\SI{3}{kWh}$ d'essence par $\SI{100}{km}$. Combien de kilomètres cette énergie permet-elle théoriquement de parcourir ? Commente.
\tcblower
\textbf{1. Masse d'éthanol par hectare}

Masse de sucre par hectare :
\[ m_{\text{sucre}} = 80\,\text{t} \times 70\,\text{kg/t} = \SI{5600}{kg} \]
Masse d'éthanol (rendement 50 \%) :
\[ m_{\text{éthanol}} = 5600 \times 0,50 = \SI{2800}{kg} \]

\textbf{2. Énergie contenue (PCI)}
\[ E = m \times \text{PCI} = 2800 \times 7,4 = \SI{20720}{kWh} \]

\textbf{3. Distance théorique parcourue}
\[ d = \dfrac{E}{\text{Conso}} \times 100 = \dfrac{20720}{3} \times 100 \approx \SI{690700}{km} \]

\textbf{Commentaire} : ce potentiel théorique ($\approx \SI{700000}{km}$) équivaut à plus de 20 ans d'utilisation d'une moto-taxi ($\approx \SI{30000}{km/an}$). Il faut toutefois nuancer : ce calcul ne tient pas compte de l'énergie dépensée pour la culture (engrais, fuel tracteurs), le transport, la distillation (énergie thermique) — c'est le \textbf{bilan énergétique net} (Energy Return on Investment) qui détermine la viabilité réelle. Néanmoins, l'ordre de grandeur confirme l'intérêt stratégique du bioéthanol pour réduire les importations de carburant au Cameroun.

\textbf{Conclusion} : un hectare de canne fournit $\SI{2800}{kg}$ d'éthanol ($\approx \SI{20720}{kWh}$), de quoi faire rouler théoriquement une moto-taxi pendant plus de vingt ans.
\end{exoresolu}

\courssub{Biodiesel et biogaz : les autres biocarburants}

\begin{methode}[Distinguer et décrire les trois grandes filières de biocarburants]
\begin{enumerate}
\item \textbf{Bioéthanol} : fermentation alcoolique de sucres (canne, manioc) par levures, puis distillation. Matière première type Cameroun : canne à sucre (SOSUCAM), manioc.
\item \textbf{Biodiesel} : \cle{transestérification} d'une huile végétale (triglycérides) avec un alcool (méthanol $\mathrm{CH_3OH}$) en présence d'un catalyseur basique (soude $\mathrm{NaOH}$ ou potasse $\mathrm{KOH}$) :
\[ \text{Triglycéride} + 3\,\mathrm{CH_3OH} \xrightarrow{\mathrm{NaOH}} 3\,\text{Esters méthyliques d'acides gras (EMAG = Biodiesel)} + \text{Glycérol} \]
Le glycérol est un sous-produit valorisable (savonnerie, cosmétique, pharmaceutique).
\item \textbf{Biogaz} : \cle{méthanisation} (digestion anaérobie) de déchets organiques (fumier, résidus de récolte, ordures ménagères, boues) par un consortium de bactéries (hydrolytiques, acétogènes, méthanogènes) dans un \emph{digesteur} étanche. Produit : mélange $\approx$ 60 \% $\mathrm{CH_4}$ (méthane) + 40 \% $\mathrm{CO_2}$ (+ traces $\mathrm{H_2S}$), utilisable pour cuisson, éclairage, cogénération électricité/chaleur. Le résidu (\emph{digestat}) est un engrais organique riche en azote minéral.
\end{enumerate}
\textbf{Réflexe} : associe toujours chaque biocarburant à son \emph{procédé}, sa \emph{matière première} et son \emph{usage final}.
\end{methode}

\begin{popfigure}
\begin{center}
\begin{tikzpicture}[node distance=1.5cm, >=Stealth, font=\small]
\tikzset{
    tank/.style={draw=popdark, thick, fill=popgreenL!30, rounded corners, minimum width=2.5cm, minimum height=2cm},
    pipe/.style={draw=popdark, thick, ->},
    label/.style={align=center, text=popdark},
    biogas/.style={draw=poporange, thick, ->, dashed},
    digestat/.style={draw=popgreen, thick, ->}
}
% Entrées
\node[label, align=center] (entree){Déchets organiques\\(fumier, résidus, O.M.)};
\node[tank, below=0.5cm of entree, align=center] (digesteur){DIGESTEUR\\(Anaérobiose, 35--38°C,\\pH 7--8, brassage)};
\draw[pipe] (entree) -- (digesteur);
% Sorties
\node[label, right=2cm of digesteur, align=center] (gaz){Biogaz\\$\approx$ 60\% $\mathrm{CH_4}$ + 40\% $\mathrm{CO_2}$};
\draw[biogas] (digesteur.east) -- (gaz.west);
\node[label, below=1.5cm of gaz, align=center] (usage){Cuisson, Électricité,\\Chaleur, Carburant (GNV)};
\draw[pipe] (gaz.south) -- (usage.north);
\node[label, below=2cm of digesteur, align=center] (digest){Digestat\\Engrais organique riche (N, P, K)};
\draw[digestat] (digesteur.south) -- (digest.north);
% Bactéries
\node[label, fill=popcream, rounded corners, inner sep=3pt, left=1.5cm of digesteur, align=center]{Consortium bactérien\\(Hydrolyse $\rightarrow$ Acidogenèse $\rightarrow$ Acétogenèse $\rightarrow$ Méthanogenèse)};
\end{tikzpicture}
\end{center}
\legende{Schéma fonctionnel d'un digesteur de méthanisation (biogaz) : entrées, processus bactérien en 4 étapes, sorties énergétiques et agronomiques.}
\end{popfigure}

\begin{exoresolu}[Choix d'une filière adaptée à une exploitation agricole camerounaise]
Une coopérative agricole de la région du Littoral dispose de : (a) $\SI{10}{t}$ de manioc par an ; (b) une palmeraie produisant des régimes de fruits ; (c) les déjections de 50 porcs et les résidus de récolte.\\
1. Pour chaque ressource, indique le biocarburant le mieux adapté et le procédé de production.\\
2. La méthanisation des déchets produit $\SI{120}{m^3}$ de biogaz par an, contenant 60 \% de méthane. Sachant que le pouvoir énergétique du méthane est de $\SI{10}{kWh.m^{-3}}$, calcule l'énergie annuelle récupérable.\\
3. Cite deux avantages annexes de la méthanisation pour la coopérative.
\tcblower
\textbf{1. Adéquation ressource--filière}

\begin{center}
\begin{tabularx}{\linewidth}{|l|X|X|}
\hline
\rowcolor{popblueL} Ressource & Biocarburant adapté & Procédé \\
\hline
Manioc (amidon) & Bioéthanol & Hydrolyse de l'amidon (amylases) $\rightarrow$ Fermentation alcoolique (levures) $\rightarrow$ Distillation \\
\hline
Fruits de palme (huile) & Biodiesel & Transestérification de l'huile de palme avec du méthanol (catalyseur $\mathrm{NaOH}$) \\
\hline
Déjections porcines + résidus & Biogaz & Méthanisation (digestion anaérobie bactérienne) en digesteur étanche \\
\hline
\end{tabularx}
\end{center}

\textbf{2. Énergie récupérable}

Volume de méthane dans le biogaz :
\[ V(\mathrm{CH_4}) = 120 \times \dfrac{60}{100} = \SI{72}{m^3} \]
Énergie annuelle :
\[ E = 72 \times 10 = \SI{720}{kWh} \]
Soit de quoi alimenter la cuisson d'une dizaine de foyers ruraux ou produire de l'électricité via un petit groupe électrogène adapté (cogénération).

\textbf{3. Avantages annexes de la méthanisation}

\begin{itemize}
\item Le \emph{digestat} résiduel est un engrais organique riche en azote ammoniacal, directement assimilable, qui améliore les rendements des cultures vivrières (manioc, maïs) et réduit l'achat d'engrais chimiques.
\item Assainissement de l'exploitation : suppression des tas de fumier en putréfaction à l'air libre $\rightarrow$ réduction des odeurs, des mouches, des germes pathogènes et de la pollution des nappes/eaux de surface par lessivage.
\end{itemize}

\textbf{Conclusion} : la coopérative valorise l'ensemble de ses biomasses : bioéthanol (manioc), biodiesel (palme), biogaz + engrais (déchets). C'est un modèle d'économie circulaire agricole.
\end{exoresolu}

\courssub{Biocarburants vs énergies fossiles : avantages et limites}

\begin{methode}[Construire une comparaison argumentée (avantages / limites)]
Pour comparer rigoureusement, utilise toujours les mêmes critères et présente-les en tableau :
\begin{enumerate}
\item \textbf{Origine et renouvelabilité} : fossile = stock non renouvelable (pétrole, charbon, millions d'années) ; biocarburant = flux renouvelable (biomasse cultivée chaque saison).
\item \textbf{Bilan carbone (cycle de vie)} : le $\mathrm{CO_2}$ émis à la combustion du biocarburant a été capté par la plante lors de sa croissance (photosynthèse) $\rightarrow$ bilan \emph{quasi neutre} (hors énergie fossile utilisée pour la culture/transformation). Les fossiles libèrent du carbone géologique stocké $\rightarrow$ aggravation nette de l'effet de serre.
\item \textbf{Pollution locale (combustion)} : biocarburants = moins de soufre ($\mathrm{SO_2}$), moins de particules fines (PM), mais émettent $\mathrm{NO_x}$ et $\mathrm{CO_2}$ ; fossiles = $\mathrm{SO_2}$, particules, $\mathrm{NO_x}$, HAP, métaux lourds.
\item \textbf{Limites structurelles des biocarburants} : compétition sols alimentaires / sols énergétiques (« manger ou rouler »), risque de déforestation (palmier à huile), forte consommation d'eau et d'intrants (engrais, pesticides), rendement surfacique parfois faible (ex. soja vs canne).
\item \textbf{Conclusion nuancée attendue au Bac} : les biocarburants sont un \emph{complément} indispensable aux fossiles dans la transition, à condition d'une production durable (déchets, terres marginales, jatropha sur sols dégradés, 2e/3e génération) et d'un mix énergétique diversifié (hydro, solaire, éolien).
\end{enumerate}
\end{methode}

\begin{exoresolu}[Biocarburants contre énergies fossiles : débat argumenté]
Au Cameroun, les coupures d'électricité et le coût des importations de carburant pèsent sur l'économie. Un débat oppose deux élèves : Awa affirme que « les biocarburants sont la solution miracle », tandis que Bertrand soutient qu'« ils sont pires que le pétrole ».\\
1. Construis un tableau comparatif biocarburants / énergies fossiles selon quatre critères.\\
2. En t'appuyant sur ce tableau, rédige une conclusion argumentée qui dépasse les deux positions extrêmes.
\tcblower
\textbf{1. Tableau comparatif}

\begin{center}
\begin{tabularx}{\linewidth}{|l|X|X|}
\hline
\rowcolor{popblueL} Critère & Biocarburants & Énergies fossiles \\
\hline
Renouvelabilité & Renouvelables : biomasse régénérée à l'échelle de la saison ou de quelques années & Non renouvelables : stocks formés en millions d'années, en épuisement inexorable \\
\hline
Bilan carbone (cycle de vie) & Quasi neutre : $\mathrm{CO_2}$ émis $\approx$ $\mathrm{CO_2}$ capté par photosynthèse (hors intrants fossiles) & Très négatif : libération de carbone fossile géologique, cause majeure du réchauffement climatique \\
\hline
Pollution locale (combustion) & Faible teneur en soufre, moins de particules ; émettent $\mathrm{NO_x}$, $\mathrm{CO}$, $\mathrm{CO_2}$ & Émissions importantes de $\mathrm{SO_2}$ (pluies acides), particules fines (PM2,5), $\mathrm{NO_x}$, HAP cancérigènes \\
\hline
Impacts annexes & Compétition cultures vivrières/énergétiques, risque déforestation, demande en eau/engrais, ILUC (changement d'affectation des sols indirect) & Marées noires, pollution forage/transport, dépendance géopolitique/importations, volatilité des prix \\
\hline
\end{tabularx}
\end{center}

\textbf{2. Conclusion argumentée}

Awa et Bertrand ont chacun partiellement raison, mais leurs positions sont caricaturales. Les biocarburants ne sont pas une « solution miracle » : produits à grande échelle sans régulation foncière, ils peuvent concurrencer les cultures vivrières (manioc, maïs, riz) et menacer la sécurité alimentaire d'un pays comme le Cameroun, tout en incitant à la déforestation pour planter du palmier à huile. Mais ils ne sont pas « pires que le pétrole » : leur bilan carbone est structurellement meilleur, ils permettent de valoriser des déchets agricoles (méthanisation) ou des terres marginales (jatropha sur sols dégradés), réduisent la facture d'importation de carburants (devise) et créent des emplois ruraux non délocalisables. La position scientifiquement correcte et citoyennement responsable est donc nuancée : les biocarburants constituent un \emph{complément renouvelable} aux énergies fossiles, à intégrer dans un \textbf{mix énergétique diversifié} (hydroélectricité, solaire photovoltaïque, biomasse) et à condition d'être produits \textbf{durablement} (certification, priorité aux déchets et résidus, protection des forêts et de la sécurité alimentaire).

\textbf{Conclusion} : les biocarburants sont un outil précieux de la transition énergétique camerounaise, à condition de ne jamais sacrifier la sécurité alimentaire ni la biodiversité.
\end{exoresolu}

\courssub{Énergies renouvelables, éducation et citoyenneté}

\begin{methode}[Construire un message de sensibilisation efficace (Savoir-être)]
Pour rédiger un texte d'éducation ou d'information (affiche, discours, article, spot radio) sur les énergies renouvelables :
\begin{enumerate}
\item \textbf{Accroche} : pars d'un problème concret vécu par le public (délestages, coût du bois/charbon, fumées des cuisines traditionnelles, déchets dans les rues).
\item \textbf{Argumente avec des faits scientifiques} : une donnée chiffrée locale, un mécanisme simple (photosynthèse, fermentation, méthanisation), jamais de slogans vides.
\item \textbf{Propose des solutions adaptées au contexte} : digesteur familial pour éleveurs, foyers améliorés à bois/charbon, valorisation des déchets de marché (biogaz), solaire pour éclairage/pompage.
\item \textbf{Anticipe les objections} : coût initial, maintenance, disponibilité de la matière première ; réponds avec des arguments économiques (amortissement, emplois locaux, subventions possibles).
\item \textbf{Conclus par un appel à l'action} clair, réaliste et valorisant la citoyenneté (protection environnement, santé, économie).
\end{enumerate}
\end{methode}

\begin{exoresolu}[Rédaction d'un message de sensibilisation : Biogaz familial à Yaoundé]
Dans le cadre de la semaine de l'environnement, le club Sciences de ton lycée te charge de rédiger un court texte (10 à 15 lignes) destiné aux habitants d'un quartier de Yaoundé, pour promouvoir l'installation de digesteurs familiaux produisant du biogaz à partir des déchets de cuisine. Rédige ce texte en respectant la méthode de sensibilisation.
\tcblower
\textbf{Texte de sensibilisation proposé :}

« Chaque jour, nos quartiers de Yaoundé jettent des tonnes de déchets de cuisine : épluchures de manioc, de plantain, restes de repas, fumier de poulaillers. Pourtant, ces déchets cachent un trésor : le \emph{biogaz}. Grâce à un digesteur familial enterré, des bactéries transforment, sans oxygène, ces matières organiques en méthane, un gaz combustible qui peut cuire vos repas et éclairer votre maison. Une famille produisant $\SI{2}{kg}$ de déchets par jour peut obtenir $\approx \SI{0,5}{m^3}$ de biogaz, soit de quoi cuire un repas quotidien, sans acheter de bois ni de charbon. Résultat : moins de dépenses (économie de $\approx \SI{15000}{FCFA}$/mois), moins de fumées toxiques dans la cuisine — qui provoquent toux, conjonctivites et maladies respiratoires chez les femmes et enfants —, moins de déchets dans les caniveaux et les ruisseaux, et en prime un excellent engrais naturel, le digestat, pour vos jardins et pépinières. Le coût d'installation (environ $\SI{150000}{FCFA}$ pour un modèle familial) est amorti en deux à trois ans d'économies de charbon. Alors, agissons ensemble : transformons nos déchets en énergie propre, pour la santé de nos familles, la propreté de notre ville et la protection de nos forêts. Le biogaz, c'est l'énergie du quartier, par le quartier, pour le quartier ! »

\textbf{Analyse de la démarche} : le texte part d'un problème concret (déchets, coût du charbon, santé), explique le mécanisme scientifique simplement (méthanisation anaérobie), donne des chiffres parlants ($\SI{2}{kg}$/jour, $\SI{0,5}{m^3}$, $\SI{15000}{FCFA}$/mois, amortissement 2-3 ans), anticipe l'objection du coût initial, et se termine par un appel à l'action citoyen et collectif.

\textbf{Conclusion} : un bon message de sensibilisation combine rigueur scientifique, données chiffrées locales, solutions concrètes adaptées au contexte et appel à l'action mobilisateur.
\end{exoresolu}

\begin{savaistu}[Initiatives camerounaises et perspectives]
\begin{itemize}
\item \textbf{SOSUCAM (Mbandjock/Nkoteng)} : principale usine sucrière, potentiel énorme de bagasse (cogénération électricité/vapeur) et de melasse (bioéthanol). Projet d'unité de distillation d'éthanol carburant à l'étude.
\item \textbf{SCDP / SOCAPALM} : huileries palme. Valorisation des effluents (POME) par méthanisation pour l'autonomie énergétique des usines. Biodiesel (EMAG) testé sur flottes de camions.
\item \textbf{Programme PNDP / MINEE / GIZ} : déploiement de digesteurs familiaux et institutionnels (lycées, prisons, hôpitaux) dans les régions du Nord, de l'Adamaoua et de l'Est (élevage bovin important).
\item \textbf{Jatropha curcas} : culture sur terres dégradées/marginales (Nord) pour biodiesel sans compétition alimentaire. Rendement encore perfectible.
\item \textbf{Mix électrique} : Le Cameroun vise 25 \% d'énergies renouvelables (hors hydroélectricité) dans le mix à l'horizon 2035 (solaire PV, biomasse, petit hydro).
\end{itemize}
\end{savaistu}

\begin{attention}[Les 5 erreurs qui coûtent des points au Bac]
\begin{enumerate}
\item \textbf{Confondre fermentation alcoolique et respiration cellulaire.} La fermentation se déroule \emph{sans oxygène} (anaérobiose) et produit éthanol + $\mathrm{CO_2}$ + 2 ATP ; la respiration consomme $\mathrm{O_2}$ et produit $\mathrm{CO_2}$ + $\mathrm{H_2O}$ + $\approx$ 36 ATP. Écrire $\ce{C6H12O6 + 6O2 -> ...}$ pour une fermentation est une erreur éliminatoire.
\item \textbf{Oublier l'étape d'hydrolyse pour les plantes amylacées.} Les levures ne fermentent pas directement l'amidon : le manioc, le maïs ou la patate douce doivent d'abord être hydrolysés en glucose (amylases, maltage, cuisson). La canne à sucre, elle, fournit directement des sucres fermentescibles (saccharose $\rightarrow$ glucose + fructose).
\item \textbf{Se tromper dans les calculs de rendement.} Erreurs classiques : appliquer le pourcentage dans le mauvais sens (diviser au lieu de multiplier par 0,9), oublier de convertir les tonnes en kilogrammes ($\times 1000$), ou confondre rendement massique (kg/kg) et rendement volumique (L/t). \textbf{Écris toujours la formule littérale avant l'application numérique.}
\item \textbf{Affirmer que les biocarburants sont « zéro pollution » ou « zéro $\mathrm{CO_2}$ ».} Leur combustion émet bien du $\mathrm{CO_2}$ et des $\mathrm{NO_x}$ ; c'est le \emph{bilan global} (photosynthèse incluse) qui est \emph{quasi neutre}. Rédige : « bilan carbone proche de la neutralité » ou « recyclage du carbone atmosphérique », jamais « n'émet pas de $\mathrm{CO_2}$ ».
\item \textbf{Présenter les biocarburants comme LA solution unique, sans limites.} Le correcteur attend une conclusion nuancée mentionnant au moins une limite majeure : compétition avec les cultures vivrières (sécurité alimentaire), risque de déforestation (palmier à huile), consommation d'eau et d'engrais (eutrophisation), ou faible rendement surfacique (ex. soja). Une copie qui ignore les limites perd les points d'argumentation et d'esprit critique (Savoir-être).
\end{enumerate}
\end{attention}

\newpage

\coursec{Exercices}

\courssub{Je vérifie mes connaissances}

\begin{exercice}[QCM : Biomasse et biocarburants]
Parmi les affirmations suivantes, une seule est exacte. La recopier en la complétant.
\begin{enumerate}
\item La biomasse utilisée pour les biocarburants de première génération provient exclusivement de déchets agricoles.
\item Le bioéthanol est produit par fermentation alcoolique de sucres (saccharose, glucose, amidon hydrolysé) par des levures en conditions \emph{aérobies}.
\item Le biodiesel (EMHV) est obtenu par transestérification d'huiles végétales ou de graisses animales avec un alcool (généralement le méthanol) en présence d'un catalyseur basique.
\item Le biogaz est un mélange majoritairement composé de dioxyde de carbone ($\ce{CO2}$) et d'hydrogène ($\ce{H2}$).
\end{enumerate}
\end{exercice}

\begin{exercice}[Vrai / Faux justifié : Contexte énergétique camerounais]
Dire si les affirmations sont vraies ou fausses et justifier brièvement votre réponse en une phrase.
\begin{enumerate}
\item Le Cameroun importe la totalité de sa consommation de pétrole raffiné malgré sa production de brut.
\item La canne à sucre, le manioc et le palmier à huile sont les trois principales filières identifiées pour la production de biocarburants de première génération au Cameroun.
\item La fermentation alcoolique libère de l'énergie sous forme d'ATP avec un rendement énergétique supérieur à celui de la respiration cellulaire.
\item L'utilisation de biocarburants permet de réduire les émissions nettes de gaz à effet de serre car le $\ce{CO2}$ émis à la combustion a été fixé par la photosynthèse lors de la croissance de la plante.
\end{enumerate}
\end{exercice}

\begin{exercice}[Définitions et équations bilan]
\begin{enumerate}
\item Définir le terme \cle{biomasse} dans le contexte de la production d'énergie.
\item Écrire l'équation bilan de la fermentation alcoolique réalisée par les levures (\emph{Saccharomyces cerevisiae}) à partir du glucose. Préciser les conditions de milieu.
\item Écrire l'équation bilan simplifiée de la méthanisation (production de biogaz) à partir de la matière organique (représentée par $\ce{C6H12O6}$). Citer les deux principaux gaz composant le biogaz.
\end{enumerate}
\end{exercice}

\begin{exercice}[Calcul direct : Rendement surfacique]
Un hectare de canne à sucre produit en moyenne \SI{80}{\tonne} de cannes. Le taux de saccharose est de \SI{14}{\%} de la masse fraîche. Le rendement de conversion saccharose $\rightarrow$ éthanol est de \SI{90}{\%} du rendement théorique (théorique : \SI{180}{\gram} glucose $\rightarrow$ \SI{92}{\gram} éthanol + \SI{88}{\gram} $\ce{CO2}$ ; masse molaire $\ce{C2H5OH}$ = \SI{46}{\gram\per\mol}).
Calculer la masse d'éthanol (en \SI{}{\kg}) potentiellement producible par hectare. (On néglige la différence de masse molaire entre glucose et saccharose hydrolysé pour ce calcul simplifié).
\end{exercice}

\courssub{Je m'entraîne}

\begin{exercice}[Expérience de fermentation : identification et suivi]
On réalise une expérience de fermentation alcoolique. On introduit dans un ballon \SI{100}{\mL} de jus de canne à sucre stérilisé (teneur en sucre initiale : \SI{180}{\gram\per\liter}) et une suspension de levures. Le ballon est muni d'un tube plongeant dans un tube à essai contenant de l'eau de chaux. On suit la masse du ballon et le trouble de l'eau de chaux pendant \SI{48}{\hour} à \SI{30}{\celsius}.
\begin{enumerate}
\item Quel gaz est responsable du trouble de l'eau de chaux ? Écrire l'équation de la réaction de détection.
\item Expliquer la diminution de la masse du ballon au cours du temps.
\item Au bout de \SI{48}{\hour}, la concentration massique en éthanol mesurée est de \SI{85}{\gram\per\liter}. Calculer le rendement de conversion du sucre en éthanol (en \%) par rapport au rendement théorique maximal (\SI{180}{\gram} sucre $\rightarrow$ \SI{92}{\gram} éthanol).
\item Pourquoi la fermentation s'arrête-t-elle spontanément avant la consommation totale des sucres si le taux d'éthanol dépasse \SI{12-15}{\% vol} ?
\end{enumerate}
\end{exercice}

\begin{exercice}[Effet de la température sur la production d'éthanol]
On a réalisé trois fermentations identiques (même jus, même inoculum) à trois températures constantes : \SI{15}{\celsius}, \SI{30}{\celsius} et \SI{45}{\celsius}. La figure ci-dessous représente la concentration massique en éthanol (\SI{}{\gram\per\liter}) en fonction du temps (heures).
\begin{popfigure}
\begin{tikzpicture}
\begin{axis}[
width=\linewidth, height=7cm,
xlabel={Temps (h)}, ylabel={[\ce{C2H5OH}] (\SI{}{\gram\per\liter})},
xmin=0, xmax=50, ymin=0, ymax=100,
legend pos=north west, grid=major,
]
\addplot[popblue, thick, mark=*] coordinates {(0,0) (10,10) (20,35) (30,60) (40,80) (50,85)};
\addlegendentry{30 °C}
\addplot[poporange, thick, mark=square*] coordinates {(0,0) (10,5) (20,15) (30,30) (40,45) (50,55)};
\addlegendentry{15 °C}
\addplot[popgreen, thick, mark=triangle*] coordinates {(0,0) (10,20) (20,30) (30,32) (40,32) (50,32)};
\addlegendentry{45 °C}
\end{axis}
\end{tikzpicture}
\legende{Évolution de la concentration massique en éthanol lors de la fermentation du jus de canne à trois températures.}
\end{popfigure}
\begin{enumerate}
\item Déterminer la température optimale de production d'éthanol parmi celles testées. Justifier par deux arguments tirés du graphique.
\item Expliquer pourquoi la courbe à \SI{45}{\celsius} montre une production rapide initiale puis un arrêt précoce à faible concentration finale.
\item Expliquer pourquoi la courbe à \SI{15}{\celsius} montre une cinétique plus lente mais une production finale encore en augmentation à \SI{50}{\hour}.
\end{enumerate}
\end{exercice}

\begin{exercice}[Comparaison de matières premières pour le bioéthanol au Cameroun]
Le tableau ci-dessous donne des données moyennes pour trois filières candidates au Cameroun.

\begin{center}
\begin{tabularx}{\linewidth}{|l|X|X|X|}\hline
\rowcolor{popblueL}
\textbf{Paramètre} & \textbf{Canne à sucre} & \textbf{Manioc} & \textbf{Maïs} \\ \hline
Rendement surfacique brut (\SI{}{\tonne\per\hectare\per\an}) & 80 (tiges) & 25 (racines) & 3 (grains) \\
Teneur en matière fermentescible (\% masse fraîche) & 14\,\% (saccharose) & 30\,\% (amidon) & 70\,\% (amidon) \\
Rendement éthanol théorique (\SI{}{\liter\per\tonne} matière première) & 70 & 280 & 400 \\
Coût de production estimé (\SI{}{\FCFA\per\liter} éthanol) & 350 & 420 & 480 \\
Besoin en eau (\SI{}{\milli\meter\per\an}) & 1500 & 1000 & 600 \\
Concurrence alimentaire & Faible (sucre export) & Forte (aliment de base) & Forte (aliment de base/poulet) \\ \hline
\end{tabularx}
\end{center}
\begin{enumerate}
\item Calculer le rendement surfacique en éthanol (\SI{}{\liter\per\hectare\per\an}) pour chaque filière.
\item En tenant compte du rendement surfacique, du coût de production, de la concurrence alimentaire et des besoins en eau, argumenter pour recommander la filière la plus adaptée au contexte camerounais (zones soudano-sahéliennes vs zones forestières humides).
\end{enumerate}
\end{exercice}

\begin{exercice}[Dimensionnement d'un digesteur à biogaz]
Un éleveur porcin dispose de 200 porcs produisant au total \SI{500}{\kg} de fumier frais par jour. La matière sèche (MS) représente \SI{20}{\%} du fumier frais. La matière organique (MO) représente \SI{80}{\%} de la MS. Le pouvoir méthanigène de ce fumier est de \SI{0,25}{\cubic\meter\per\kg} MO introduite. Le biogaz contient \SI{60}{\%} de $\ce{CH4}$.
\begin{enumerate}
\item Calculer la quantité de matière organique (\SI{}{\kg} MO) apportée quotidiennement au digesteur.
\item Calculer le volume journalier de biogaz produit (en \SI{}{\cubic\meter}).
\item Sachant qu'un foyer rural consomme en moyenne \SI{1,5}{\cubic\meter} de biogaz par jour pour la cuisson, combien de foyers cet éleveur peut-il approvisionner ?
\item Citer deux conditions indispensables au bon fonctionnement du digesteur (paramètres physico-chimiques).
\end{enumerate}
\end{exercice}

\courssub{Je me prépare au Bac}

\begin{exercice}[Sujet type Bac : Fermentation industrielle du jus de canne (12 pts)]
On souhaite optimiser la production de bioéthanol à partir de jus de canne à sucre dans une unité pilote au Cameroun. On dispose des résultats suivants obtenus dans des fermenteurs de \SI{100}{\liter}, à pH 4,5, avec un inoculum de \emph{Saccharomyces cerevisiae} constant.

\textbf{Document 1 :} Cinétique de fermentation à \SI{30}{\celsius}.
\begin{center}
\begin{tabular}{|c|c|c|c|c|c|c|}\hline
Temps (h) & 0 & 12 & 24 & 36 & 48 & 60 \\ \hline
Concentration massique glucose + fructose (\SI{}{\gram\per\liter}) & 180 & 120 & 60 & 15 & 2 & 0 \\ \hline
Concentration massique éthanol (\SI{}{\gram\per\liter}) & 0 & 35 & 70 & 85 & 88 & 88 \\ \hline
Concentration massique biomasse levures (\SI{}{\gram\per\liter} MS) & 2 & 8 & 15 & 18 & 18 & 17 \\ \hline
Volume $\ce{CO2}$ dégagé cumulé (\SI{}{\liter}) & 0 & 1500 & 3100 & 3900 & 4300 & 4400 \\ \hline
\end{tabular}
\end{center}

\textbf{Document 2 :} Rendement final en éthanol (\% du théorique) en fonction de la température.
\begin{center}
\begin{tabular}{|c|c|c|c|c|}\hline
Température (°C) & 20 & 30 & 35 & 40 \\ \hline
Rendement (\% théorique) & 78 & 92 & 85 & 40 \\ \hline
\end{tabular}
\end{center}

\textbf{Questions :}
\begin{enumerate}
\item \textbf{Analyse du Document 1 :}
\begin{enumerate}
\item Tracer sur un même graphique (papier millimétré) les courbes d'évolution de la concentration en sucres, en éthanol et en biomasse levures en fonction du temps. (3 pts)
\item Déterminer la phase de croissance exponentielle des levures. Calculer le temps de génération (doublement) moyen pendant cette phase. (2 pts)
\item Calculer le rendement de conversion sucre $\rightarrow$ éthanol (en \%) à 48 h. Conclure sur l'efficacité de la souche. (2 pts)
\end{enumerate}
\item \textbf{Interprétation du Document 2 :}
\begin{enumerate}
\item Identifier la température optimale de fermentation. (0,5 pt)
\item Expliquer la baisse du rendement à \SI{40}{\celsius} par rapport à \SI{30}{\celsius} en invoquant l'effet de la température sur les enzymes glycolytiques et la membrane cellulaire des levures. (2 pts)
\end{enumerate}
\item \textbf{Synthèse et application :}
\begin{enumerate}
\item Écrire l'équation bilan de la fermentation alcoolique. Préciser le rôle de l'ATP produit. (1,5 pt)
\item L'unité pilote traite \SI{50}{\cubic\meter} de jus par jour (\SI{180}{\gram\per\liter} sucre). En considérant le rendement optimal (92\,\%), calculer la production journalière d'éthanol en litres (masse volumique éthanol = \SI{0,79}{\kg\per\liter}). (1 pt)
\end{enumerate}
\end{enumerate}
\end{exercice}

\begin{exercice}[Sujet type Bac : Biogaz et gestion durable des déchets à Dschang (10 pts)]
La ville de Dschang (Ouest Cameroun) produit environ \SI{50}{\tonne} de déchets ménagers par jour, dont \SI{60}{\%} de matière organique putrescible (déchets de cuisine, marchés). La municipalité envisage d'installer une unité de méthanisation centralisée pour produire du biogaz et du digestat (engrais).

\textbf{Document :} Caractéristiques de la matière organique des déchets de Dschang.
\begin{itemize}
\item Humidité : \SI{75}{\%} (masse eau / masse fraîche).
\item Matière sèche (MS) : \SI{25}{\%} de la masse fraîche.
\item Matière organique (MO) : \SI{85}{\%} de la MS.
\item Pouvoir méthanigène (BMP) : \SI{0,35}{\cubic\meter\per\kg} MO.
\item Composition du biogaz : \SI{55}{\%} $\ce{CH4}$, \SI{45}{\%} $\ce{CO2}$.
\item PCI du méthane : \SI{35,9}{\mega\joule\per\cubic\meter}.
\end{itemize}

\textbf{Questions :}
\begin{enumerate}
\item Calculer la masse de matière organique (en \SI{}{\tonne}) disponible quotidiennement pour la méthanisation. (1,5 pt)
\item Calculer le volume journalier de biogaz produit (en \SI{}{\cubic\meter}). (1 pt)
\item Calculer l'énergie primaire journalière contenue dans le biogaz produit (en \SI{}{\mega\joule} puis en \SI{}{\kWh}). (1,5 pt)
\item Si un groupe électrogène au biogaz a un rendement électrique de \SI{35}{\%}, quelle puissance électrique continue (en \SI{}{\kW}) cette unité peut-elle fournir ? (1 pt)
\item Le digestat sortant contient \SI{3}{\%} d'azote (N) sur matière sèche. Si \SI{80}{\%} de la MS entrée ressort en digestat, calculer la production journalière d'azote fertilisant (en \SI{}{\kg} N). (1,5 pt)
\item Discuter deux avantages environnementaux et un risque sanitaire potentiel lié à l'épandage du digestat non hygiénisé. (2,5 pts)
\item Proposer une mesure technique pour valoriser le $\ce{CO2}$ du biogaz plutôt que de le rejeter à l'atmosphère. (1 pt)
\end{enumerate}
\end{exercice}

\begin{exercice}[Sujet type Bac : Synthèse argumentée — Biocarburants et sécurité énergétique (8 pts)]
\textbf{Sujet :} « Les biocarburants de première génération (bioéthanol, biodiesel) constituent-ils une solution durable pour réduire le déficit énergétique et la dépendance aux importations pétrolières du Cameroun ? »

À partir de vos connaissances et des données du cours, rédiger un texte argumenté structuré (introduction, développement en deux parties, conclusion) de 15 à 20 lignes maximum. Votre réponse doit obligatoirement comporter :
\begin{itemize}
\item Une introduction définissant le contexte et annonçant le plan.
\item Une première partie présentant \textbf{deux avantages majeurs} (économique, environnemental, social/agricole) avec des exemples camerounais précis.
\item Une deuxième partie présentant \textbf{deux limites ou risques majeurs} (concurrence alimentaire, bilan énergétique, usage des terres, coûts) avec des arguments chiffrés ou qualitatifs.
\item Une conclusion nuancée proposant une piste pour les générations futures (biocarburants 2G/3G, mix énergétique).
\end{itemize}
La qualité de la rédaction scientifique, la précision du vocabulaire et la logique de l'argumentation seront évaluées.
\end{exercice}

\newpage

\coursec{Corrigés des exercices}

\begin{corrige}[Exercice 1 — QCM : Biomasse et biocarburants]
L'affirmation exacte est la \textbf{3} : \og Le biodiesel (EMHV) est obtenu par transestérification d'huiles végétales ou de graisses animales avec un alcool (généralement le méthanol) en présence d'un catalyseur basique. \fg

\textbf{Analyse des autres propositions :}
\begin{enumerate}
\item \textbf{Fausse.} Les biocarburants de \emph{première} génération sont précisément produits à partir de parties \emph{comestibles} des plantes cultivées (sucres, amidon, huiles) : canne à sucre, manioc, maïs, palmier à huile. Les déchets agricoles et lignocellulose relèvent de la \emph{deuxième} génération.
\item \textbf{Fausse.} La fermentation alcoolique est une voie \emph{anaérobie} : elle se déroule en \textbf{l'absence de dioxygène}. En présence d'$\ce{O2}$, les levures pratiquent la respiration et produisent peu d'éthanol.
\item[\textbf{4.}] \textbf{Fausse.} Le biogaz est majoritairement composé de \textbf{méthane} $\ce{CH4}$ (50--70\,\%) et de $\ce{CO2}$ (30--50\,\%), avec des traces de $\ce{H2S}$. L'hydrogène n'est pas un constituant majeur.
\end{enumerate}
\end{corrige}

\begin{corrige}[Exercice 2 — Vrai / Faux justifié : Contexte énergétique camerounais]
\begin{enumerate}
\item \textbf{Vrai.} Le Cameroun extrait du pétrole brut (région de Logone, Rio del Rey) mais ses capacités de raffinage (SONARA, Limbe) sont insuffisantes et souvent indisponibles, ce qui impose d'importer l'essentiel des produits raffinés (super, gasoil) — d'où l'intérêt stratégique des biocarburants.
\item \textbf{Vrai.} Ces trois filières disposent au Cameroun de surfaces cultivées importantes et de matières fermentescibles ou estérifiables : saccharose (canne), amidon (manioc) et huile (palme).
\item \textbf{Faux.} La fermentation ne dégrade le glucose qu'en éthanol et $\ce{CO2}$ : elle ne libère que \textbf{2 ATP} par molécule de glucose, contre environ \textbf{30 à 32 ATP} pour la respiration aérobie complète. Son rendement énergétique est donc très inférieur.
\item \textbf{Vrai.} Le bilan carbone est quasi neutre : le $\ce{CO2}$ relâché à la combustion du biocarburant correspond au $\ce{CO2}$ fixé par photosynthèse pendant la croissance de la biomasse, contrairement aux énergies fossiles qui libèrent un carbone séquestré depuis des millions d'années. (Nuance : il faut déduire les émissions liées aux engrais, au transport et à la distillation.)
\end{enumerate}
\end{corrige}

\begin{corrige}[Exercice 3 — Définitions et équations bilan]
\begin{enumerate}
\item \textbf{Biomasse} : ensemble de la matière organique d'origine végétale ou animale (plantes cultivées, résidus agricoles et forestiers, déchets organiques, effluents d'élevage) susceptible d'être valorisée comme source d'énergie (combustion directe, fermentation, méthanisation, estérification).
\item \textbf{Fermentation alcoolique} (par \emph{Saccharomyces cerevisiae}) :
\[ \ce{C6H12O6 -> 2 C2H5OH + 2 CO2} + \text{énergie (2 ATP)} \]
Conditions de milieu : \textbf{anaérobiose} (absence de $\ce{O2}$), température optimale voisine de \SI{28-32}{\degreeCelsius}, pH acide (4--5), présence de levures vivantes.
\item \textbf{Méthanisation} (équation simplifiée) :
\[ \ce{C6H12O6 -> 3 CH4 + 3 CO2} \]
Le biogaz est constitué principalement de \textbf{méthane} $\ce{CH4}$ (50--70\,\%, le combustible utile) et de \textbf{dioxyde de carbone} $\ce{CO2}$ (30--50\,\%).
\end{enumerate}
\end{corrige}

\begin{corrige}[Exercice 4 — Calcul direct : Rendement surfacique]
\textbf{Étape 1 : masse de saccharose par hectare.}
\[ m_{\text{sucre}} = \SI{80}{t} \times \frac{14}{100} = \SI{11,2}{t} = \SI{11200}{kg} \]

\textbf{Étape 2 : rendement théorique de conversion.} D'après l'équation, \SI{180}{g} de sucre donnent \SI{92}{g} d'éthanol, soit un rapport massique :
\[ r_{\text{th}} = \frac{92}{180} = \num{0,511} \]
Éthanol théorique : $11200 \times \num{0,511} = \SI{5724}{kg}$.

\textbf{Étape 3 : rendement réel (90\,\% du théorique).}
\[ m_{\text{éthanol}} = 5724 \times \num{0,90} = \SI{5152}{kg} \approx \SI{5,2}{t\per\hectare} \]

\textbf{Conclusion :} un hectare de canne peut produire environ \SI{5150}{kg} d'éthanol, soit (à \SI{0,79}{kg\per\liter}) environ \SI{6500}{L} de bioéthanol par hectare et par an.

\textbf{Point de vigilance :} ne pas oublier d'appliquer les \emph{deux} facteurs (taux de saccharose, puis rendement de 90\,\%) et de convertir les tonnes en kilogrammes.
\end{corrige}

\begin{corrige}[Exercice 5 — Expérience de fermentation : identification et suivi]
\begin{enumerate}
\item Le gaz qui trouble l'eau de chaux est le \textbf{dioxyde de carbone} $\ce{CO2}$. Il forme un précipité blanc de carbonate de calcium :
\[ \ce{CO2 + Ca(OH)2 -> CaCO3 v + H2O} \]
\item Le ballon perd de la masse car le $\ce{CO2}$ produit par la fermentation \emph{s'échappe} du milieu réactionnel (il traverse le tube barboteur). Or, d'après l'équation $\ce{C6H12O6 -> 2 C2H5OH + 2 CO2}$, une partie de la matière du glucose quitte le ballon sous forme gazeuse : la masse du ballon diminue donc au fur et à mesure (la perte de masse est même proportionnelle à l'avancement de la fermentation).
\item Rendement théorique maximal : \SI{180}{g\per\liter} de sucre $\rightarrow$ \SI{92}{g\per\liter} d'éthanol.
\[ \text{Rendement} = \frac{85}{92} \times 100 = \num{92,4}\,\% \]
La conversion atteint environ \textbf{92\,\%} du rendement théorique : la fermentation est efficace (les pertes correspondent à la biomasse levurienne, au glycérol et aux sous-produits).
\item L'éthanol est un \textbf{alcool toxique} pour les levures elles-mêmes : au-delà de 12--15\,\% vol, il perturbe la perméabilité de la membrane plasmique, dénature les protéines et inhibe les enzymes glycolytiques. Les levures meurent ou deviennent métaboliquement inactives : la fermentation s'auto-inhibe avant épuisement des sucres.
\end{enumerate}
\end{corrige}

\begin{corrige}[Exercice 6 — Effet de la température sur la production d'éthanol]
\begin{enumerate}
\item La température optimale est \textbf{\SI{30}{\degreeCelsius}}. Deux arguments tirés du graphique :
\begin{itemize}
\item la concentration finale d'éthanol y est la plus élevée : \SI{85}{g\per\liter} contre \SI{55}{g\per\liter} à \SI{15}{\degreeCelsius} et \SI{32}{g\per\liter} à \SI{45}{\degreeCelsius} ;
\item la cinétique y est la plus rapide : la pente initiale est la plus forte (environ \SI{35}{g\per\liter} atteints dès \SI{20}{h}).
\end{itemize}
\item À \SI{45}{\degreeCelsius}, la production initiale est rapide car l'agitation moléculaire accélère les réactions enzymatiques ; mais cette température \textbf{dénature progressivement les enzymes} glycolytiques (déstructuration de leur conformation spatiale) et \textbf{altère la membrane plasmique} des levures (fluidité excessive, perte d'étanchéité). Les levures meurent : la fermentation s'arrête prématurément à \SI{32}{g\per\liter}.
\item À \SI{15}{\degreeCelsius}, les enzymes restent \textbf{intactes mais peu actives} (cinétique enzymatique ralentie par le froid) : la pente est faible. En revanche, il n'y a ni dénaturation ni mortalité : les levures restent vivantes et continuent à fermenter lentement, d'où une courbe encore croissante à \SI{50}{h} (\SI{55}{g\per\liter}, non plafonnée).
\end{enumerate}
\textbf{Point de vigilance :} distinguer clairement l'effet du froid (ralentissement \emph{réversible}) de l'effet de la chaleur forte (dénaturation \emph{irréversible}).
\end{corrige}

\begin{corrige}[Exercice 7 — Comparaison de matières premières pour le bioéthanol au Cameroun]
\begin{enumerate}
\item Rendement surfacique = rendement brut $\times$ rendement en éthanol par tonne :
\begin{align*}
\text{Canne : } & 80 \times 70 = \SI{5600}{\liter\per\hectare\per\an} \\
\text{Manioc : } & 25 \times 280 = \SI{7000}{\liter\per\hectare\per\an} \\
\text{Maïs : } & 3 \times 400 = \SI{1200}{\liter\per\hectare\per\an}
\end{align*}
\item \textbf{Argumentation :}
\begin{itemize}
\item \emph{Rendement surfacique :} le manioc est le plus productif (\SI{7000}{\liter\per\hectare\per\an}), suivi de la canne (\SI{5600}{}) ; le maïs est très faible (\SI{1200}{}) et de plus le plus coûteux (\SI{480}{FCFA\per\liter}) : il est à écarter.
\item \emph{Concurrence alimentaire :} le manioc est un \textbf{aliment de base} des populations camerounaises (bâtons, farine, foufou) : le détourner vers l'éthanol menacerait la sécurité alimentaire. La canne, destinée au sucre (filière SOSUCAM, largement exportatrice ou industrielle), entre moins en compétition directe avec l'alimentation de subsistance.
\item \emph{Coût :} la canne est la moins chère (\SI{350}{FCFA\per\liter}).
\item \emph{Besoins en eau :} la canne exige \SI{1500}{mm\per\an} : elle convient aux \textbf{zones forestières humides du Sud} (pluviométrie élevée). Le manioc (\SI{1000}{mm\per\an}) et le maïs (\SI{600}{mm\per\an}) tolèrent mieux les zones soudano-sahéliennes, mais la concurrence alimentaire y est rédhibitoire.
\end{itemize}
\textbf{Recommandation :} développer prioritairement la \textbf{filière canne à sucre} dans les zones forestières humides (bon compromis rendement/coût/faible concurrence alimentaire), en réservant le manioc à l'alimentation ; à terme, orienter la recherche vers les biocarburants de 2\textsuperscript{e} génération (résidus de canne, tiges de manioc, pailles) qui échappent au dilemme \og assiette ou réservoir \fg.
\end{enumerate}
\end{corrige}

\begin{corrige}[Exercice 8 — Dimensionnement d'un digesteur à biogaz]
\begin{enumerate}
\item Matière sèche quotidienne : $500 \times \num{0,20} = \SI{100}{kg}$ MS. Matière organique :
\[ m_{\text{MO}} = 100 \times \num{0,80} = \SI{80}{kg\ MO\per\day} \]
\item Volume de biogaz :
\[ V = 80 \times \num{0,25} = \SI{20}{m^3\per\day} \]
(dont $20 \times \num{0,60} = \SI{12}{m^3}$ de méthane $\ce{CH4}$.)
\item Nombre de foyers approvisionnables :
\[ N = \frac{20}{\num{1,5}} = \num{13,3} \quad \Rightarrow \quad \textbf{13 foyers} \text{ (arrondi par défaut)} \]
\item Deux conditions indispensables :
\begin{itemize}
\item une \textbf{température stable} dans la gamme mésophile (\SI{30-40}{\degreeCelsius}), car les bactéries méthanogènes sont thermosensibles ;
\item une \textbf{anaérobiose stricte} (digesteur parfaitement étanche à l'air), car les méthanogènes sont des anaérobies stricts. (On accepte aussi : pH proche de la neutralité 6,8--7,5 ; humidité suffisante du substrat ; rapport C/N équilibré.)
\end{itemize}
\end{enumerate}
\end{corrige}

\begin{corrige}[Exercice 9 — Sujet type Bac : Fermentation industrielle du jus de canne (12 pts)]
\textbf{Barème indicatif :} 1a : 3 pts ; 1b : 2 pts ; 1c : 2 pts ; 2a : 0,5 pt ; 2b : 2 pts ; 3a : 1,5 pt ; 3b : 1 pt.

\textbf{1.a. Tracé des courbes (3 pts).} Attendus : repère orthonormé, temps en abscisse (0 à \SI{60}{h}), concentrations en ordonnée (0 à \SI{200}{g\per\liter}), trois courbes légendées :
\begin{itemize}
\item \textbf{sucres} : décroissante de 180 à 0 \SI{}{g\per\liter} (consommation du substrat) ;
\item \textbf{éthanol} : croissante en S de 0 à \SI{88}{g\per\liter} (palier à partir de 48 h) ;
\item \textbf{biomasse} : croissante de 2 à \SI{18}{g\per\liter}, puis stationnaire/légère décroissance (17 à 60 h).
\end{itemize}
Points de vigilance : axes gradués et légendés avec unités, courbes identifiées, tracé régulier passant par les points du tableau.

\textbf{1.b. Phase exponentielle et temps de génération (2 pts).} La croissance exponentielle s'observe entre 0 et \SI{24}{h} (biomasse multipliée par 7,5 : 2 $\rightarrow$ \SI{15}{g\per\liter}). Sur l'intervalle 0--12 h, la biomasse passe de 2 à 8, soit \textbf{deux doublements} ($2 \rightarrow 4 \rightarrow 8$) en \SI{12}{h} :
\[ t_g = \frac{12\ \text{h}}{2} = \SI{6}{h} \]
Le temps de génération moyen est d'environ \textbf{6 heures}. (Vérification sur 12--24 h : 8 $\rightarrow$ 15, proche d'un doublement, cohérent avec $t_g \approx$ 6--7 h.)

\textbf{1.c. Rendement de conversion à 48 h (2 pts).} Rendement théorique : \SI{180}{g\per\liter} de sucre $\rightarrow$ \SI{92}{g\per\liter} d'éthanol.
\[ \text{Rendement} = \frac{88}{92} \times 100 = \num{95,7}\,\% \approx 96\,\% \]
\textbf{Conclusion :} la souche est \textbf{très efficace} : presque tout le sucre fermentescible est converti en éthanol ; les pertes ($\sim$4\,\%) correspondent à la production de biomasse, de glycérol et d'autres métabolites.

\textbf{2.a. Température optimale (0,5 pt).} \textbf{\SI{30}{\degreeCelsius}} (rendement maximal : 92\,\% du théorique).

\textbf{2.b. Baisse du rendement à \SI{40}{\degreeCelsius} (2 pts).} À \SI{40}{\degreeCelsius}, la chaleur provoque une \textbf{dénaturation partielle des enzymes glycolytiques} (perte de la structure tertiaire, donc du site actif), ce qui ralentit puis bloque la glycolyse ; simultanément, la \textbf{membrane plasmique} devient hyperfluide et perméable (fuites d'ions, perte du potentiel membranaire), et l'éthanol déjà produit en diffuse plus facilement, amplifiant sa toxicité. Il en résulte une mortalité levurienne élevée et une fermentation incomplète : le rendement chute à 40\,\%.

\textbf{3.a. Équation bilan et rôle de l'ATP (1,5 pt).}
\[ \ce{C6H12O6 -> 2 C2H5OH + 2 CO2} + 2\ \text{ATP} \]
L'ATP produit (2 molécules par glucose, lors de la glycolyse) fournit l'\textbf{énergie nécessaire au métabolisme de la levure} : maintien de l'intégrité cellulaire, synthèses (multiplication), transport actif — c'est ce qui permet à la levure de survivre et de se multiplier en anaérobiose.

\textbf{3.b. Production journalière d'éthanol (1 pt).}
\begin{align*}
m_{\text{sucre}} &= 50\,000\ \text{L} \times 180\ \text{g/L} = 9\,000\,000\ \text{g} = \SI{9000}{kg} \\
m_{\text{éthanol théorique}} &= 9000 \times \frac{92}{180} = \SI{4600}{kg} \\
m_{\text{éthanol réel}} &= 4600 \times \num{0,92} = \SI{4232}{kg} \\
V_{\text{éthanol}} &= \frac{4232}{\num{0,79}} = \SI{5357}{L} \approx \SI{5,4}{m^3\per\day}
\end{align*}
\textbf{Point de vigilance :} convertir les \SI{}{m^3} en litres ($\SI{1}{m^3} = \SI{1000}{L}$), appliquer le rendement de 92\,\% \emph{après} le rendement stœchiométrique, et diviser par la masse volumique (et non multiplier).
\end{corrige}

\begin{corrige}[Exercice 10 — Sujet type Bac : Biogaz et gestion durable des déchets à Dschang (10 pts)]
\textbf{Barème indicatif :} 1 : 1,5 pt ; 2 : 1 pt ; 3 : 1,5 pt ; 4 : 1 pt ; 5 : 1,5 pt ; 6 : 2,5 pts ; 7 : 1 pt.

\begin{enumerate}
\item \textbf{Masse de matière organique disponible :}
\begin{align*}
m_{\text{putrescible}} &= 50 \times \num{0,60} = \SI{30}{t\per\day} \\
m_{\text{MS}} &= 30 \times \num{0,25} = \SI{7,5}{t\per\day} \\
m_{\text{MO}} &= 7,5 \times \num{0,85} = \SI{6,375}{t\per\day} \approx \SI{6,4}{t\ MO\per\day}
\end{align*}
\item \textbf{Volume de biogaz :}
\[ V_{\text{biogaz}} = 6375\ \text{kg} \times \SI{0,35}{m^3\per\kg} = \SI{2231}{m^3\per\day} \approx \SI{2230}{m^3\per\day} \]
\item \textbf{Énergie primaire :} seul le méthane est combustible :
\[ V_{\ce{CH4}} = 2231 \times \num{0,55} = \SI{1227}{m^3\per\day} \]
\[ E = 1227 \times \num{35,9} = \SI{44053}{MJ\per\day} \approx \SI{44}{GJ\per\day} \]
Conversion : $\SI{1}{kWh} = \SI{3,6}{MJ}$, donc $E = \dfrac{44053}{\num{3,6}} = \SI{12237}{kWh\per\day} \approx \SI{12,2}{MWh\per\day}$.
\item \textbf{Puissance électrique continue :}
\[ E_{\text{élec}} = 12237 \times \num{0,35} = \SI{4283}{kWh\per\day} \]
\[ P = \frac{4283\ \text{kWh}}{24\ \text{h}} = \SI{178}{kW} \]
L'unité peut fournir environ \textbf{\SI{180}{kW}} en continu (de quoi alimenter un quartier ou des équipements municipaux).
\item \textbf{Azote fertilisant :}
\begin{align*}
m_{\text{digestat (MS)}} &= 7,5 \times \num{0,80} = \SI{6}{t\per\day} = \SI{6000}{kg\per\day} \\
m_{\text{N}} &= 6000 \times \num{0,03} = \SI{180}{kg\ N\per\day}
\end{align*}
\item \textbf{Avantages environnementaux} (deux attendus) :
\begin{itemize}
\item \textbf{réduction des émissions de méthane} : sans méthanisation, la décharge anaérobie des déchets organiques libère du $\ce{CH4}$, gaz à effet de serre environ 25 fois plus puissant que le $\ce{CO2}$ ; le capter et le brûler réduit fortement l'empreinte climatique ;
\item \textbf{substitution d'engrais} : le digestat riche en azote remplace les engrais minéraux importés (économie de devises, fertilisation des sols agricoles de l'Ouest) et \textbf{assainit la ville} (moins de dépôts sauvages, moins de miasmes et de mouches).
\end{itemize}
\textbf{Risque sanitaire :} un digestat \emph{non hygiénisé} peut contenir des \textbf{agents pathogènes} (coliformes, \emph{Salmonella}, œufs d'helminthes) et des métaux lourds : son épandage brut sur des cultures maraîchères peut contaminer les denrées et les nappes phréatiques. Il faut donc un traitement d'hygiénisation (compostage thermophile, pasteurisation) avant usage.
\item \textbf{Valorisation du $\ce{CO2}$ :} procéder à l'\textbf{épuration du biogaz} (lavage à l'eau sous pression ou absorption sur amines) pour séparer le $\ce{CO2}$ du méthane : on obtient du \textbf{biométhane} injectable/utilisable comme carburant, et le $\ce{CO2}$ capté peut être valorisé (serres agricoles, carbonatation des boissons, culture de microalgues) au lieu d'être rejeté.
\end{enumerate}
\textbf{Points de vigilance :} bien appliquer les pourcentages \emph{en cascade} (60\,\% puis 25\,\% puis 85\,\%) ; ne calculer l'énergie que sur la fraction $\ce{CH4}$ ; convertir MJ en kWh en divisant par 3,6 ; distinguer puissance (kW) et énergie (kWh).
\end{corrige}

\begin{corrige}[Exercice 11 — Sujet type Bac : Synthèse argumentée (8 pts)]
\textbf{Barème indicatif :} introduction (1 pt) ; partie 1 — deux avantages argumentés et contextualisés (2,5 pts) ; partie 2 — deux limites argumentées (2,5 pts) ; conclusion nuancée avec perspective (1 pt) ; qualité rédactionnelle et vocabulaire scientifique (1 pt).

\textbf{Éléments de corrigé (texte attendu, 15--20 lignes) :}

\emph{Introduction.} Le Cameroun, bien que producteur de pétrole brut, importe l'essentiel de ses produits raffinés et connaît des déficits énergétiques récurrents (délestages, coût des carburants). Les biocarburants de première génération — bioéthanol issu de la fermentation des sucres (canne, manioc) et biodiesel issu de la transestérification des huiles (palme) — sont présentés comme une réponse. Sont-ils pour autant une solution \emph{durable} ? Nous montrerons d'abord leurs atouts réels, puis leurs limites structurelles, avant de nuancer.

\emph{Partie 1 — Avantages.} (1) \textbf{Avantage économique et énergétique} : produire localement de l'éthanol à partir de la canne (filière sucrière déjà industrialisée : SOSUCAM à Nkoteng et Bandjock) réduirait la facture d'importation de carburants, préserverait les devises et créerait des emplois ruraux (planteurs, distilleries). (2) \textbf{Avantage environnemental} : le bilan carbone des biocarburants est quasi neutre — le $\ce{CO2}$ dégagé à la combustion a été fixé par photosynthèse durant la croissance — contrairement aux carburants fossiles ; leur usage atténuerait donc les émissions nettes de gaz à effet de serre et la pollution urbaine (moins de soufre que le gasoil).

\emph{Partie 2 — Limites.} (1) \textbf{Concurrence alimentaire} : détourner le manioc ou le maïs (aliments de base) vers l'éthanol ferait monter les prix alimentaires et menacerait la sécurité alimentaire ; de même, étendre le palmier à huile peut accélérer la déforestation. (2) \textbf{Bilan énergétique et coûts} : la production d'éthanol exige de l'énergie (distillation), de l'eau (\SI{1500}{mm\per\an} pour la canne) et des intrants ; le coût de production (\SI{350-480}{FCFA\per\liter}) reste élevé et le rendement surfacique limité (quelques milliers de litres/ha), de sorte que les biocarburants 1G ne peuvent couvrir qu'une fraction des besoins nationaux.

\emph{Conclusion.} Les biocarburants de première génération sont un \emph{complément} utile mais non une solution durable à eux seuls. La voie d'avenir réside dans les générations 2G (résidus agricoles : bagasse, pailles, tiges de manioc, sans concurrence alimentaire) et 3G (microalgues), intégrées à un mix énergétique diversifié (hydroélectricité, solaire, biogaz) et à une politique d'efficacité énergétique.

\textbf{Points de vigilance pour l'examinateur :} exiger une structure visible (intro/développement/conclusion), au moins un exemple camerounais précis dans chaque partie, un vocabulaire scientifique exact (fermentation, transestérification, bilan carbone, biomasse), et une conclusion qui ne se contente pas d'un avis tranché mais propose une perspective (2G/3G, mix énergétique).
\end{corrige}

\newpage

\coursec{Fiche bilan}

\begin{popfigure}
\begin{center}
\begin{center}\tcbox[colback=popblue,colframe=popblue,coltext=white,arc=9pt,boxsep=4pt,left=6pt,right=6pt]{\bfseries\small Recherche de nouvelles sources d'énergie}\end{center}
\vspace{-2pt}
\begin{tcbraster}[raster columns=3,raster equal height=rows,raster column skip=6pt,raster row skip=6pt,enhanced,blankest]
\begin{tcolorbox}[enhanced,colback=popgreen,colframe=popgreen,coltext=white,boxrule=0.9pt,arc=3pt,left=4pt,right=4pt,top=3pt,bottom=3pt,boxsep=1pt,halign=left]\bfseries\scriptsize Biomasse et biocarburants \scriptsize matière végétale, déchets organiques\end{tcolorbox}
\begin{tcolorbox}[enhanced,colback=poporange,colframe=poporange,coltext=white,boxrule=0.9pt,arc=3pt,left=4pt,right=4pt,top=3pt,bottom=3pt,boxsep=1pt,halign=left]\bfseries\scriptsize Fermentation alcoolique \scriptsize \ce{C6H12O6 -> 2C2H5OH + 2CO2}\end{tcolorbox}
\begin{tcolorbox}[enhanced,colback=poppurple,colframe=poppurple,coltext=white,boxrule=0.9pt,arc=3pt,left=4pt,right=4pt,top=3pt,bottom=3pt,boxsep=1pt,halign=left]\bfseries\scriptsize Bioéthanol \scriptsize canne à sucre, manioc\end{tcolorbox}
\begin{tcolorbox}[enhanced,colback=popgold,colframe=popgold,coltext=white,boxrule=0.9pt,arc=3pt,left=4pt,right=4pt,top=3pt,bottom=3pt,boxsep=1pt,halign=left]\bfseries\scriptsize Biodiesel et biogaz \scriptsize huile de palme, méthanisation\end{tcolorbox}
\begin{tcolorbox}[enhanced,colback=poppink,colframe=poppink,coltext=white,boxrule=0.9pt,arc=3pt,left=4pt,right=4pt,top=3pt,bottom=3pt,boxsep=1pt,halign=left]\bfseries\scriptsize Avantages et limites \scriptsize renouvelable vs fossile\end{tcolorbox}
\begin{tcolorbox}[enhanced,colback=popblue!70,colframe=popblue!70,coltext=white,boxrule=0.9pt,arc=3pt,left=4pt,right=4pt,top=3pt,bottom=3pt,boxsep=1pt,halign=left]\bfseries\scriptsize Citoyenneté \scriptsize éduquer, informer, agir\end{tcolorbox}
\end{tcbraster}

\end{center}
\legende{Carte mentale de la leçon : les six idées forces autour des biocarburants et des énergies renouvelables.}
\end{popfigure}

\begin{bilan}
\textbf{Définitions clés}
\begin{itemize}
  \item \cle{Biomasse} : ensemble de la matière organique d'origine végétale ou animale (végétaux cultivés, résidus agricoles, déchets organiques) utilisable comme source d'énergie.
  \item \cle{Biocarburant} : combustible liquide ou gazeux produit à partir de la biomasse par transformation biologique (fermentation) ou physico-chimique (transestérification).
  \item \cle{Énergie renouvelable} : énergie dont la source se régénère naturellement à l'échelle humaine (soleil, biomasse, eau, vent), par opposition aux énergies fossiles (pétrole, charbon, gaz naturel) dont les stocks s'épuisent.
\end{itemize}

\textbf{Les trois grandes familles de biocarburants}
\begin{center}
\begin{tabularx}{\linewidth}{|l|X|X|X|}
\hline
\rowcolor{popblueL} \textbf{Biocarburant} & \textbf{Matière première (Cameroun)} & \textbf{Procédé} & \textbf{Usage} \\
\hline
Bioéthanol & canne à sucre, manioc (sucres, amidon) & fermentation alcoolique puis distillation & carburant (mélange essence) \\
\hline
Biodiesel & huile de palme & transestérification des huiles végétales & carburant diesel \\
\hline
Biogaz & déchets organiques, fumiers & méthanisation (fermentation anaérobie) & cuisson, chauffage, électricité \\
\hline
\end{tabularx}
\end{center}

\textbf{Formules et bilans à connaître par c\oe ur}
\begin{itemize}
  \item Fermentation alcoolique (anaérobie, levures) : \ce{C6H12O6 -> 2C2H5OH + 2CO2} + énergie.
  \item Respiration cellulaire (aérobie) : \ce{C6H12O6 + 6O2 -> 6CO2 + 6H2O} + énergie.
  \item La fermentation alcoolique est une oxydation \emph{incomplète} du glucose : elle dégage moins d'énergie que la respiration.
\end{itemize}

\textbf{Méthodes express}
\begin{itemize}
  \item \textbf{Expliquer une production de bioéthanol} : broyage de la matière première $\rightarrow$ hydrolyse de l'amidon en sucres simples $\rightarrow$ fermentation par des levures en anaérobiose $\rightarrow$ distillation pour concentrer l'éthanol.
  \item \textbf{Comparer deux énergies} : toujours croiser trois critères — renouvelabilité de la source, bilan en gaz à effet de serre, coût et disponibilité locale.
  \item \textbf{Argumenter} : citer un avantage (renouvelable, moins polluant, valorise les déchets, crée des emplois ruraux) puis une limite (compétition avec les cultures vivrières, surface agricole, rendement énergétique).
\end{itemize}

\textbf{Avantages et limites des biocarburants}
\begin{itemize}
  \item Avantages : source renouvelable ; réduction des émissions polluantes ; valorisation des déchets ; réduction de la dépendance aux importations de pétrole.
  \item Limites : concurrence avec l'alimentation (manioc, canne) ; besoin de grandes surfaces ; coût de transformation ; le bilan carbone n'est pas nul (machines, engrais, transport).
\end{itemize}
\end{bilan}

\begin{aretenir}[Checklist avant le Bac]
\begin{itemize}
  \item $\square$ Je sais définir biomasse, biocarburant et énergie renouvelable en une phrase chacun.
  \item $\square$ Je connais le bilan de la fermentation alcoolique : \ce{C6H12O6 -> 2C2H5OH + 2CO2}.
  \item $\square$ Je sais que la fermentation se déroule en anaérobiose grâce à des levures.
  \item $\square$ Je distingue fermentation alcoolique (peu d'énergie) et respiration (oxydation complète).
  \item $\square$ Je cite les trois biocarburants : bioéthanol, biodiesel, biogaz.
  \item $\square$ J'associe chaque biocarburant à sa matière première camerounaise : canne à sucre et manioc, huile de palme, déchets organiques.
  \item $\square$ Je décris les étapes de production du bioéthanol : broyage, hydrolyse, fermentation, distillation.
  \item $\square$ Je sais que le biodiesel s'obtient par transestérification des huiles végétales.
  \item $\square$ Je sais que le biogaz (riche en méthane) provient de la méthanisation des déchets.
  \item $\square$ Je compare biocarburants et énergies fossiles selon renouvelabilité, pollution et coût.
  \item $\square$ Je cite au moins deux avantages et deux limites des biocarburants.
  \item $\square$ Je sais argumenter sur l'intérêt des énergies renouvelables face aux déficits énergétiques du Cameroun.
\end{itemize}
\end{aretenir}

\findecours

\end{document}
